\documentclass[11pt]{article}
\usepackage{graphicx,xcolor}
\usepackage[T1]{fontenc}
\usepackage[utf8]{inputenc}
\usepackage{lmodern,amsmath,amssymb,amsthm,mathtools,booktabs,longtable}
\usepackage[margin=1in]{geometry}
\usepackage{microtype}
\usepackage[hidelinks]{hyperref}
\newtheorem{theorem}{Theorem}[section]
\newtheorem{proposition}[theorem]{Proposition}
\newtheorem{lemma}[theorem]{Lemma}
\theoremstyle{definition}\newtheorem{definition}[theorem]{Definition}
\theoremstyle{remark}\newtheorem{remark}[theorem]{Remark}
\newcommand{\C}{\mathbb C}\newcommand{\Z}{\mathbb Z}\newcommand{\R}{\mathbb R}
\DeclareMathOperator{\tr}{tr}\DeclareMathOperator{\Rea}{Re}
\newcommand{\status}{\textbf{Research draft. GRH is not proved here.}}
\date{November 2, 2026}


\newtheorem{corollary}[theorem]{Corollary}
\newcommand{\gshift}[2]{#1\langle #2\rangle}
\date{November 5, 2026}
\title{Infinite dessin orders, projective resolutions and spectral determinants\\\large Constructions, invariant theory and the arithmetic comparison problem}
\author{Research working draft}
\begin{document}
\maketitle
\begin{abstract}
We work throughout with infinite-dimensional locally gentle algebras and completed surface orders arising from explicit cellular clean dessins on compact Riemann surfaces. We construct their projective resolutions, Cartan determinants and dual dessins; develop determinant semi-invariants, Schur modules and cyclic-quiver character kernels; and calculate actual lattice almost-split sequences and an affine-type diamond quotient. A constrained-path Fock algebra retains the gentle relations and has a compact extension with incidence boundary map. Its degree spectral triple and an inserted arithmetic Fredholm family are explicit. Jacobi sums provide theta-dessin cellular Frobenius determinants with a positive local metric; inert primes obstruct the same real two-dimensional carrier. Recent spectral approaches are compared with their actual hypotheses. Selected power-series, module and matrix statements are verified in Lean4. A derived branch functor repairs projective determinant multiplicity; tame descent recovers a hereditary order but forces multiplicities for full-deck coefficient kernels. We also give a requested-error adaptive theta/readout construction for three fixed completed Dirichlet functions, with executed rational training and explicit sampling budgets, and derive a cutoff-uniform original-family type-I remainder envelope, retaining its unresolved exponent loss. The arithmetic positive determinant comparison remains open; no proof of GRH is claimed.
\end{abstract}
{\small\tableofcontents}
\section{Cellular dessins and their infinite orders}
We start with the surface-order constructions of Schreiber \cite{SurfaceOrders,SurfaceDuality}, retaining the original sources unchanged. The generic-resolution source \cite{Persistence} is background, not a substitute for exactness of a particular complex. All primary algebra examples below come from explicit clean Belyi maps. Finite-dimensional objects occur as modules or specified quotient strata of those infinite algebras, not as replacements for them. The initial arithmetic target is a primitive nonprincipal Dirichlet $L$-function. Finite-order Hecke characters, Dedekind zeta functions and general Artin representations require separate comparison theorems; Artin holomorphy is not assumed.

A clean dessin on a compact oriented surface is specified by a nonempty finite dart set $H$, a fixed-point-free involution $\alpha$, and a permutation $\sigma$ such that $\langle\alpha,\sigma\rangle$ is transitive. Contract the white vertices of degree two. The resulting cellular graph has vertices the $\sigma$-orbits, edges the $\alpha$-orbits, and faces the $\phi$-orbits, where
\[
 \tau=\sigma^{-1},\qquad \phi=\alpha\sigma^{-1},\qquad
 2-2g=\#\operatorname{cyc}(\sigma)-|H|/2+\#\operatorname{cyc}(\phi).
\]
Write $[h]=\{h,\alpha h\}$. Our quiver has vertices these edges and arrows $a_h:[h]\to[\tau h]$. Paths are composed in travel order: $ab$ traverses $a$ first. Set
\[
 A=kQ/(a_h a_{\phi h}:h\in H),\quad J=A_{>0},\quad
 \Lambda=\varprojlim A/J^m.
\]
Each vertex has two incoming and two outgoing arrows. Each arrow has precisely one permitted successor $a_{\tau h}$ and one forbidden successor $a_{\phi h}$; these are distinct because $\alpha$ has no fixed points. The same holds for predecessors. This is the \emph{locally gentle} condition. Every permitted cycle repeats indefinitely, so $A$ is infinite dimensional. We reserve ``gentle'' without qualification for the convention additionally imposing finite dimension.

\begin{proposition}[The actual split order]\label{thm:order}
For a permitted cycle $c$ of length $\ell$, let $B_c$ be the hereditary cycle order in $M_\ell(k[[u_c]])$: diagonal and forward matrix entries are in $k[[u_c]]$, and entries crossing the chosen cycle cut are in $u_c k[[u_c]]$. Then $\Lambda$ is the suborder of $B=\prod_c B_c$ whose two diagonal constant terms labelled $h,\alpha h$ agree. With $R=k[[u]]$ acting by $u\mapsto(u_c I)_c$, it is finite free of rank $\sum_c\ell(c)^2$; the conductor contains $uB$.
\end{proposition}
\begin{proof}
In a cycle order a path of length less than $\ell$ is the corresponding matrix unit, multiplied by $u_c$ exactly when it crosses the cut; another full turn multiplies it by $u_c$. Thus completed paths in each permitted branch give all entries of $B_c$. An arrow belongs to only that branch, whereas $e_{[h]}$ has diagonal entries at both occurrences $h,\alpha h$. Hence the only condition on branchwise formal coefficients is the displayed equality of diagonal constants. This proves both injectivity and surjectivity of the path-normal-form identification. Multiplication respects it: switching branches across a paired edge is precisely a forbidden pair. The quotient $B/\Lambda$ is killed by $u$, so the submodule has the same generic rank as $B$. It is torsion free and finite over the DVR $R$, hence free of that rank. The path and $u$-adic topologies agree: $uA\subset J$, and if $L$ is the longest permitted cycle, any path of length at least $L+1$ is $u$ times a shorter nonempty path. Thus $J^{L+1}\subset uA$. This also justifies the completion used in the identification.
\end{proof}
The path-degree variable $q$ below is not $u$: one full turn has path degree $\ell(c)$ but $u$-degree one. Formal, Laurent and operator-norm completions will likewise remain distinct.

Here are five explicit global inputs. Permutations list darts starting at zero; $E=|H|/2$.
\begin{center}\small
\begin{tabular}{llll}
\toprule
Surface and Belyi map $\beta$ & $\sigma$ & $\alpha$ & $(g,\operatorname{rank}_R\Lambda)$\\\midrule
$\mathbb P^1$, $4z(1-z)$ & $(0)(1)$ & $(01)$ & $(0,2)$\\
$\mathbb P^1$, $(z+1)^2/(4z)$ & $(01)$ & $(01)$ & $(0,4)$\\
$y^2=x^3-1$, $x^3$ & $(012)(345)$ & $(03)(14)(25)$ & $(1,18)$\\
$\mathbb P^1$, $4z^3/(1+z^3)^2$ & $(012)(345)$ & $(03)(15)(24)$ & $(0,18)$\\
$y^2=x^3-x$, $x^2$ & $(0123)$ & $(02)(13)$ & $(1,16)$\\\bottomrule
\end{tabular}
\end{center}
The first is the nodal interval: $A=k[x,y]/(xy)$ and $\Lambda=k[[x,y]]/(xy)$. The second has $A=k\langle x,y\rangle/(x^2,y^2)$, still infinite dimensional. The two theta dessins have the same two permitted three-cycles but different face cycles, respectively $(051324)$ and $(04)(13)(25)$. The elliptic square has one permitted and one forbidden four-cycle. For example $x$ has pole order two on the cubic $y^2=x^3-1$; $x^3$ has ramification $(3,3)$ above zero, $(2,2,2)$ above one since $x^3-1=y^2$, and $(6)$ above infinity. The rational theta map has passport $(3,3),(2,2,2),(2,2,2)$. These ramification computations, and the analogous direct zeros/poles in the other rows, verify cellular Belyi inputs rather than merely abstract quivers.

\section{Projective resolutions, dual dessins and determinants}
Use right modules $P_i=e_iA$, $S_i=P_i/P_iJ$. Internal shift $\gshift{M}{n}$ places its old generator in degree $n$. Let $\mathcal R_n(i)$ be the \emph{indexing words} of length $n$ from $i$ with every adjacent pair forbidden, including $e_i$ for $n=0$. Such a word is usually zero in $A$; it still labels a summand.
\begin{theorem}[Exact resolutions]\label{thm:resolution}
There is a minimal graded projective resolution, with $P_0(i)=e_iA$ and for $n\ge1$,
\[
 P_n(i)=\bigoplus_{p=a_1\cdots a_n\in\mathcal R_n(i)}
 \gshift{e_{t(a_n)}A}{n},\qquad
 d_n(x)_p=a_nx\ \text{in the summand }a_1\cdots a_{n-1}.
\]
The augmentation $P_0(i)\to S_i$ is the quotient. For the completed order the identical finite-rank projective terms and formal-series maps resolve the residue simple. Ext and Tor against residue simples count forbidden words with the prescribed endpoints. In the graded case Ext has matching internal shift $n$.
\end{theorem}
\begin{proof}
The maps are right linear and degree zero. Adjacent differentials multiply by a forbidden pair. The augmentation kernel is the direct sum of $aA$ over initial arrows. For any arrow $a$, multiplication $e_{t(a)}A\to e_{s(a)}A$ kills exactly the permitted paths beginning with its forbidden successor $f(a)$. All other basis paths survive to distinct basis paths. Thus its kernel is $f(a)A$. For a fixed prefix, different last arrows have disjoint surviving supports; there is no cancellation between summands. This proves exactness inductively. For formal series the same assertion holds coefficient by coefficient and identifies the entire closed kernel, not just its dense polynomial subspace. Terms are idempotent projectives and all entries are arrows. Consequently applying Hom or tensor with a residue simple kills every differential, proving minimality and the Ext/Tor count. In $A$ minimality is graded; it does not require that $J$ be its ordinary Jacobson radical.
\end{proof}
The two arms never terminate in this closed class. Their labelled patterns follow $\phi$-cycles. An ungraded positive syzygy tail can have a shorter period by permutation of summands: the nodal interval has labelled period two but unordered tail period one; the elliptic theta has labelled period six and unordered tail period three. A simple need not be isomorphic to a positive syzygy; for the node it is $R$-torsion, whereas the positive syzygies are lattices.

Let $C_{ij}(q)=\sum_{d\ge0}\dim_k(e_iA_de_j)q^d$ and $E_{ij}(q)=\sum_{n\ge0}(-1)^n\dim_k\operatorname{Ext}_A^n(S_i,S_j)q^n$. Degreewise Euler characteristic of the resolution gives $EC=I$: every coefficient is finite, because a summand in homological degree $n$ starts in internal degree $n$. Constant terms are $I$, so the inverse is two-sided. Already signed $E(q)$ must not be evaluated at $-q$ again.

\begin{proposition}[Dual dessin and cycle determinant]\label{thm:cartan}
The complementary-relation algebra is another closed dessin algebra, with $\alpha^\#=\alpha$, $\sigma^\#=\phi^{-1}$, permitted successor $\phi$, and forbidden successor $\tau$. It exchanges graph vertices and faces and preserves genus. Its Cartan matrix satisfies $E(q)=C_{A^\#}(-q)$, with the opposite convention if required for Yoneda composition. For independent commuting arrow weights $w_a$,
\[
 \det C_A(\mathbf w)=
 \frac{\prod_{f\in\operatorname{cyc}(\phi)}(1-(-1)^{|f|}w(f))}
 {\prod_{c\in\operatorname{cyc}(\sigma)}(1-w(c))}.
\]
\end{proposition}
\begin{proof}
$\alpha(\sigma^\#)^{-1}=\alpha\phi=\tau$, so the claimed relations and face data follow directly; the connectedness group is unchanged. To prove the determinant, put $S_{v,a}=w_a$ when $s(a)=v$, and $T_{a,v}=1$ when $t(a)=v$, zero otherwise. On arrow states let $U_{a,b}=w_b$ for permitted adjacency and $V_{a,b}=w_b$ for forbidden adjacency. Then $TS=U+V$ and the path sum is $C=I+S(I-U)^{-1}T$. Sylvester's identity gives $\det C=\det(I+V)/\det(I-U)$. Each transition is a weighted permutation, whose cycle determinant is respectively $1-(-1)^\ell w(c)$ and $1-w(c)$. Everything holds in formal power series, hence as rational functions. This supplies a direct proof of the Bessenrodt--Holm formula in the closed dessin class \cite{BessenrodtHolm}.
\end{proof}
The five univariate determinants are, in the order above,
\[
 \frac{1+q}{1-q},\quad\frac{1+q}{1-q},\quad
 \frac{1-q^6}{(1-q^3)^2},\quad
 \frac{(1-q^2)^3}{(1-q^3)^2},\quad 1.
\]
The last cancellation does not imply finite global dimension: its forbidden resolution is infinite. These rational functions have cycle poles/zeros, not arithmetic prime Euler factors.

For clarity, a complex is a sequence with $d^2=0$, its homotopy category identifies maps differing by $dh+hd$, and $D(A)$ inverts maps inducing cohomology isomorphisms. A perfect complex is represented by a bounded complex of finite projectives. Our infinite simple resolutions compute derived Hom and tensor but are generally not perfect. The two-term presentation $[P_1\to P_0]$ is perfect even if its kernel is nonzero. Quadratic linearity is a Koszul statement; it is not by itself an equivalence between all completed derived categories. Likewise a finite-dimensional gentle graded-curve theorem \cite{GeometricGentle} cannot simply be transferred to these closed infinite orders. Here the surface walks explicitly compute the resolution endpoints: $\tau$ turns round a graph vertex and $\phi$ follows a face. For longer string objects all maps can be computed from the occurrence matrices; a universal geodesic-intersection classification for this completed category remains a separate theorem.

\subsection{Auxiliary staircase resolutions over the order}
The two filtration parameters below are not dessin arrows. Let $R=k[x,y]$
for a field $k$, and let $I=(g_0,\ldots,g_{m-1})$, $m\ge2$, where
$g_i=x^{a_i}y^{b_i}$,
\[
 a_0>\cdots>a_{m-1}=0,\qquad 0=b_0<\cdots<b_{m-1}.
\]
The positive outer exponents make $R/I$ finite length. Write $R(-a,-b)$
for a free generator born in bidegree $(a,b)$.
\begin{proposition}[Adjacent syzygies]\label{thm:staircase}
A minimal bigraded free resolution of $M=R/I$ is
\[
0\longrightarrow\bigoplus_{i=0}^{m-2}R(-a_i,-b_{i+1})
\xrightarrow{d_2}\bigoplus_{i=0}^{m-1}R(-a_i,-b_i)
\xrightarrow{d_1}R\longrightarrow M\longrightarrow0,
\]
where $d_1(e_i)=g_i$ and
$d_2(f_i)=-y^{b_{i+1}-b_i}e_i+x^{a_i-a_{i+1}}e_{i+1}$.
Deleting row $i$ from $d_2$ gives determinant $(-1)^ig_i$.
The Hilbert series numerator is
$1-\sum_i s^{a_i}t^{b_i}+\sum_{i=0}^{m-2}s^{a_i}t^{b_{i+1}}$.
\end{proposition}
\begin{proof}
Each column maps to zero. In a relation $\sum h_i g_i=0$, all terms except
the last contain $x^{a_{m-2}}$. Coprimality with $y^{b_{m-1}}$ implies
$h_{m-1}=x^{a_{m-2}}q$. Subtract $q d_2(f_{m-2})$ from the coefficient
vector; its last component vanishes. Factor $x^{a_{m-2}}$ from the remaining
generators. They form a smaller staircase with the same strict inequalities,
ending in a pure $y$-power. Induction expresses the remaining relation in
the preceding adjacent columns. The base of one generator has no relations,
since $R$ is a domain. Hence $\ker d_1=\operatorname{im}d_2$.
If $d_2(q)=0$, its first row gives $-y^{b_1}q_0=0$; successive rows then
force all $q_i=0$. This proves injectivity. Every differential entry has
positive degree, so reduction modulo $(x,y)$ kills the differentials and
the resolution is minimal. In the row-deleted matrix the unique nonzero
permutation term uses the negative $y$ entries before row $i$ and the
positive $x$ entries after it. Their exponents telescope to $b_i,a_i$,
respectively, and the sign is $(-1)^i$. Finally take the degreewise Euler
characteristic; the free summands have Hilbert denominator $(1-s)(1-t)$.
\end{proof}
For $I=(x^2,xy,y^2)$ the underlying sequence is also kernel checked
in Lean~4 over the actual ring $\operatorname{MvPolynomial}(\mathrm{Fin}\,2,k)$.
The verification proves the kernel/image equalities, the stated ideal image,
surjective quotient augmentation, projective free terms, and vanishing of
both induced Hom differentials to the constant-coefficient residue field.
It uses regular multiplication on the free modules, distinct from the
residue action on $k$. General staircases, graded minimality and the
surface-order coefficient cone remain separately scoped prose arguments.

At grade $\nu\in\mathbb N^2$, $M_\nu$ has dimension one exactly when
no generator degree is componentwise at most $\nu$. Its comparable-grade
map has rank one exactly when both endpoint grades satisfy this condition.
A CW realization takes one loop and attaches a disk once along it at each
generator degree. The active indices form an interval: the $x$ inequality
selects a final interval, the $y$ inequality an initial interval. With $c$
active disks, the cellular map $k^c\to k$ has every entry one. Thus
$\dim H_1=\mathbf1_{c=0}$ and $\dim H_2=\max(c-1,0)$, the latter counting
active adjacent syzygies. This auxiliary CW complex is not the dessin surface.

For the actual nodal order $\Lambda=k[[\xi,\eta]]/(\xi\eta)$ put
$u=\xi+\eta$ and $E=\Lambda/u^2\Lambda$.
The basis $(1,\xi,u,u\xi)$ of $E$ satisfies $\xi^2=u\xi$, $u^2=0$,
$\eta=u-\xi$. Tensoring the resolution above with
$0\to\Lambda\xrightarrow{u^2}\Lambda\to E\to0$ gives, over
$\mathcal B=\Lambda[x,y]$, free ranks $(1,m+1,2m-1,m-1)$.
Indeed the scalar resolution stays exact after tensoring over the field;
its degree-zero module is free over $\Lambda$ on the standard monomials.
Multiplication by $u^2$ is injective there and on all free terms, so the
long exact homology sequence of its mapping cone gives positive-degree
homology zero and degree-zero module $E\otimes_k M$.
In contrast, as an $R$-module this object has ranks $4(1,m,m-1)$ in its
minimal resolution. Neither finite-length $E$ nor these parameter rings
replace the infinite parent dessin algebra or its lattice category.

\subsection{A monotone quantum comparison with exact finite ranks}
\begin{proposition}[One common channel]\label{thm:channel-filtration}
Let $T$ be a positive trace-preserving linear map between finite complex
matrix algebras. For density matrices $\rho,\sigma$, their trace distance
$d(\rho,\sigma)=\tfrac12\|\rho-\sigma\|_1$ satisfies
$d(T\rho,T\sigma)\le d(\rho,\sigma)$.
For finitely many labelled states, keep fixed vertex birth grades $e_i$.
At $(r,e)$ include vertices with $e_i\le e$ and edges with distance at most
$r$, and take the flag complex. Then $T$ gives an identity-on-labels inclusion
from the input to the output bifiltration, natural in $(r,e)$.
\end{proposition}
\begin{proof}
Write the Hermitian difference as $A=A_+-A_-$ with $A_\pm\ge0$.
Positivity, the trace norm triangle inequality and trace preservation give
$\|T A\|_1\le\|T A_+\|_1+\|T A_-\|_1
=\operatorname{tr}A_++\operatorname{tr}A_-=\|A\|_1$.
Thus every input edge is an output edge; the fixed births preserve vertices,
and all cliques are preserved. Increasing either parameter preserves the
simplices, and every map is the same label inclusion, so the squares commute.
Complete positivity is sufficient but is not needed for this norm argument.
\end{proof}
For a finite inclusion $K\subset L$, work over a field and embed a basis
$Z_K$ of source one-cycles into the target edge space. If $B_L$ is the
matrix of target two-boundaries, the induced $H_1$ rank is
\[
 \rho_1(K,L)=\operatorname{rank}[B_L\ Z_K]-\operatorname{rank}B_L.
\]
Its image is $(Z_K+\operatorname{im}B_L)/\operatorname{im}B_L$; source
boundaries already map into target boundaries. This proves the formula,
including empty complexes. Computing a flag two-skeleton suffices for
$H_0,H_1$, but does not compute the full flag complex's $H_2$.
Integer ceilings of scaled exact distances preserve the inclusion.
For numerical distances, frozen integer labels define an exact combinatorial
diagram only after their inequality is checked; this is not an interval
certificate of the underlying quantum distances. None of these finite ranks
identifies a completed arithmetic determinant or proves GRH.

\subsection{A correct DG model of the infinite node}
Use cohomological degrees and an uncurved differential of degree $+1$.
An odd element $\varepsilon$ with $d\varepsilon=1$ contracts the underlying $k$-complex of every unital
left DG module: $h(m)=\varepsilon m$ satisfies $dh+hd=1$ by the graded
Leibniz rule. Thus this example has zero cohomology, contrary to the
original ground-field-resolution claim. An isomorphism on $H^0$ alone is
also insufficient for a DG quasi-isomorphism; all degrees must be checked.

\begin{proposition}[Nodal Koszul model]\label{thm:nodal-dg}
Let $k$ be a field, $R=k[\xi,\eta]$ and $f=\xi\eta$ and $K=R\oplus R\varepsilon$, with
$|\varepsilon|=-1$, $\varepsilon^2=0$, $dR=0$ and $d\varepsilon=f$.
The augmentation $K\to A=R/(\xi\eta)$, $a+b\varepsilon\mapsto[a]$,
is a DG-algebra quasi-isomorphism onto the actual infinite nodal dessin
algebra. Branch rescaling $\xi\mapsto q\xi$, $\eta\mapsto q^{-1}\eta$,
$\varepsilon\mapsto\varepsilon$ for $q\in k^\times$ lifts to a DG automorphism.
\end{proposition}
\begin{proof}
Multiplication is $(a+b\varepsilon)(c+e\varepsilon)=ac+(ae+bc)\varepsilon$.
The differential kills degree zero and sends $b\varepsilon$ to $fb$.
Its square is zero. For two odd inputs the Leibniz right side is
$fbc\varepsilon-bfc\varepsilon=0$; the other homogeneous cases follow
by $R$-linearity. Multiplication by $f$ is injective in the polynomial
domain. Hence $H^{-1}(K)=0$, $H^0(K)=R/(f)$, and all other groups vanish.
The augmentation preserves products and kills boundaries. Rescaling fixes
$f$, so its coefficientwise extension commutes with $d$ and has inverse
obtained by replacing $q$ with $q^{-1}$. This is a one-relation Koszul
construction \cite{StacksKoszulDG,StacksKoszulRegular}.
\end{proof}
Replacing $R$ by $k[[\xi,\eta]]$ gives the completed order: setting each
variable separately to zero maps onto the branch pullback, and the kernel
consists precisely of mixed monomials, hence equals $(\xi\eta)$.
Rescaling is generally semilinear over the coefficient ring, rather than
linear over the fixed diagonal DVR $k[[u]]$, $u=\xi+\eta$.
It preserves that DVR as a subring exactly when $q^2=1$, and fixes it
pointwise exactly when $q=1$. No Hilbert norm or arithmetic spectrum follows.

The corresponding Lean file checks the actual exterior-line algebra,
differential, sign-Leibniz, homotopy and polynomial ambient-resolution
identities. Full graded/category/completion objects remain separately scoped.
The two theta dessins likewise require endpoint-sensitive Ext: self
$\operatorname{Ext}^1$ is zero in both; for the elliptic theta self
$\operatorname{Ext}^3$ has dimension two, while for the spherical theta
positive even self Ext has dimension two. Outdegree alone does not count
self extensions. Minimal $A_\infty$ models use additional graded data
\cite[\S3.3]{KellerAinf}; $m_3$ is not automatically a physical curvature.


\section{Determinant semi-invariants and Schur constructions}
Assume $k$ algebraically closed of characteristic zero for invariant-ring statements. A finite-dimensional representation of the \emph{infinite} $A$ has spaces $N_i$, arrow matrices $N(a)$ satisfying the monomial relations, and base-change group $G=\prod_i\mathrm{GL}(N_i)$. With column vectors, $N(a)\mapsto g_{t(a)}N(a)g_{s(a)}^{-1}$. Let $P_0=\bigoplus P_i^{a_i}$, $P_1=\bigoplus P_i^{b_i}$, and $f:P_1\to P_0$ be a finite path-matrix map. When $\sum a_i\dim N_i=\sum b_i\dim N_i$, define
\[
 c_f(N)=\det D_f(N),\qquad D_f(N)=\operatorname{Hom}_A(f,N).
\]
It is intrinsically a determinant-line section; chosen bases give this scalar.
\begin{proposition}[Weights and exact homological scope]\label{thm:semi}
$c_f(gN)=\prod_i\det(g_i)^{b_i-a_i}c_f(N)$. Restricted determinants of this form span all semi-invariants on a closed $G$-stable relation or nilpotence locus. Set $M=\operatorname{coker}f$. On the square locus,
\[
 c_f(N)\ne0\quad\Longleftrightarrow\quad\operatorname{Hom}(M,N)=0
 \quad\Longleftrightarrow\quad\operatorname{Hom}(M,N)=\operatorname{Ext}^1(M,N)=0.
\]
If $f$ is injective, $\operatorname{coker}D_f=\operatorname{Ext}^1(M,N)$ and higher Ext vanish. Without injectivity neither that cokernel equality nor higher vanishing follows.
\end{proposition}
\begin{proof}
Naturality gives $D_f(gN)=G_1D_f(N)G_0^{-1}$, where $G_j$ repeats the relevant $g_i$. Take determinants. The contragredient coordinate-ring action has the inverse character. Schofield--van den Bergh's spanning theorem allows oriented cycles \cite[Theorem2.3]{SchofieldVdB}. Exactness of invariants for reductive $G$ in characteristic zero lifts a fixed-weight function from any closed stable locus to the ambient representation space; restrict the spanning determinants. Injectivity of the path map need not survive that restriction. Finally $\ker D_f=\operatorname{Hom}(M,N)$. Writing $K=\operatorname{im}f$, the exact sequence with projective $P_0$ expresses $\operatorname{Ext}^1(M,N)$ as $\operatorname{Hom}(K,N)/\operatorname{im}\operatorname{Hom}(P_0,N)$; the epimorphism $P_1\to K$ embeds this quotient in $\operatorname{coker}D_f$. A square finite matrix with zero kernel has zero cokernel. If $f$ is injective it is the actual length-one resolution, giving the stronger assertion.
\end{proof}
For the node, $f=\times x$, $M=\Lambda/(x)$, $N=\Lambda/(x,y^2)$ give $\dim\operatorname{coker}D_f=2$ but $\dim\operatorname{Ext}^1(M,N)=1$. For the elliptic square name arrows so that $a,b:0\to1$, $c,d:1\to0$, permitted cycle $acbd$ and forbidden cycle $adbc$. With $f=a$ and $N(a)=1$, $N(b)=N(c)=N(d)=0$, its Hom matrix is $[1]$, yet the subsequent forbidden resolution gives $\operatorname{Ext}^{2,3}(M,N)=k$. Thus nonvanishing is not all-degree derived orthogonality.

\begin{proposition}[What completion does to invariants]\label{thm:nilpotent}
A continuous $d$-dimensional representation of $\Lambda$ kills $J^d$. On the reduced nilpotent representation variety of fixed dimension vector every invariant regular function is constant. Nonzero-weight semi-invariants need not be constant.
\end{proposition}
\begin{proof}
Continuity with discrete finite-dimensional target kills some $J^m$; its Loewy filtration has at most $d$ nonzero strict steps. Choose vertexwise complements to this filtration. An arrow sends a layer $p$ into layers $q\ge p+1$. Scaling layer $p$ by $z^p$ conjugates each component by $z^{q-p}$. At $z=0$ the representation tends to the zero-arrow representation. Every invariant therefore has the same value at all points; reducedness and the Nullstellensatz give the ring assertion. In the square example the scalar $N(a)$ varies on a nilpotent locus and has weight $g_1/g_0$, so semi-invariants survive. Formal entries in a projective map are evaluated on this locus by truncating modulo $J^d$, a legitimate finite polynomial operation.
\end{proof}
On the node, $u=x+y$ is regular and $A$ is free of rank two over $k[u]$. The polynomial presentation $\times(u-\lambda)$ for $\lambda\ne0$ is injective and its determinant on $N$ is $\det(u_N-\lambda I)$, the characteristic polynomial of $u_N$ up to the sign $(-1)^d$. But $u-\lambda$ is a unit in $\Lambda$, with its geometric-series inverse, so those modules do not extend continuously to the completion. An $R$-free lattice and its $u$-nilpotent quotients are different categories. On a lattice a separately specified endomorphism $F$ has a meaningful relative characteristic polynomial over $R$; the scalar base action $uI$ alone does not create Frobenius.

\subsection{Young tableaux and a cyclic representation stratum}
For a nonempty partition $\lambda\vdash m$, choose a numbering $T$ of its boxes. On $V^{\otimes m}$ let $Y_T=R_TC_T$, the unnormalized row sum followed by the signed column sum. The classical identity $Y_T^2=h_\lambda Y_T$, where $h_\lambda$ is the hook product, makes $Y_T/h_\lambda$ an idempotent in characteristic zero. Its image is $S_\lambda(V)$, one Schur copy, rather than the whole symmetric-group isotypic component \cite{Etingof,SchurComplexes}. Its character is $s_\lambda$, computed by semistandard tableaux (weak rows, strict columns), and
\[
 \dim S_\lambda(k^d)=\prod_{(i,j)\in\lambda}\frac{d+j-i}{h_{ij}}.
\]
Since $f^{\otimes m}$ commutes with the place permutations, it induces $S_\lambda(f)$ and respects composition and zero maps. Applying this at every vertex and arrow therefore gives a genuine representation of the monomial dessin algebra: $S_\lambda(N(b))S_\lambda(N(a))=S_\lambda(0)=0$ for a forbidden pair. The empty partition does not have this property: it sends zero arrow maps to the identity of $k$.

The Littlewood--Richardson rule gives natural intertwiners
\[
 S_\lambda(V)\otimes S_\mu(V)\cong
 \bigoplus_\nu S_\nu(V)^{\oplus c_{\lambda\mu}^{\nu}}.
\]
Here $c_{\lambda\mu}^{\nu}$ counts semistandard tableaux of skew shape $\nu/\lambda$ and content $\mu$ whose row-reading word is Yamanouchi. Naturality makes the same decompositions intertwine all the arrows. The tensor operation is vertexwise with synchronous arrow maps, not $\otimes_A$ and not a Hopf-algebra tensor structure. There is no automatic tensor unit preserving these zero relations. These functors are not additive or exact; they cannot simply be applied degreewise to a projective resolution. They also need not preserve the fixed $R$-action on lattices.

In the elliptic theta set one of the two permitted three-cycle branches to zero. The other has arbitrary maps $X_1:V_0\to V_2$, $X_2:V_2\to V_1$, $X_3:V_1\to V_0$, with all spaces of dimension $d$. This is a closed representation stratum of our infinite dessin algebra. Put $M=X_3X_2X_1$.
\begin{theorem}[Cyclic invariants and a local positive kernel]\label{thm:schur}
The invariant ring on this stratum is $k[c_1(M),\ldots,c_d(M)]$, where $c_j$ are characteristic coefficients. Its total-arrow-degree Hilbert series is $\prod_{j=1}^d(1-t^{3j})^{-1}$. Common multidegree $(m,m,m)$ has basis $\chi_\lambda(M)=\tr S_\lambda(M)$, $\lambda\vdash m$, $\ell(\lambda)\le d$; other multidegrees have no invariants. Products obey the Littlewood--Richardson rule. Over $\C$, if $\|M_i\|\le R_0$ and $0\le r$ and $rR_0^2<1$, the kernel
\[
 K_r(M_i,M_j)=\det(I-rM_i\otimes\overline{M_j})^{-1}
 =\sum_\lambda r^{|\lambda|}\chi_\lambda(M_i)\overline{\chi_\lambda(M_j)}
\]
is absolutely convergent and positive semidefinite.
\end{theorem}
\begin{proof}
Cycle-trace generation for quivers, including relation quotients, is the Le Bruyn--Procesi theorem \cite{LeBruynProcesi,CrawYamagishi}. Every cycle here is a power of the full three-cycle, up to cyclic rotation. In characteristic zero Newton identities give the generators $c_j$. They are algebraically independent: take $X_1=X_2=I$ and arbitrary $X_3$. Scalar changes at the three vertices force equal arrow multidegrees. Restriction to that same section identifies degree $(m,m,m)$ with degree-$m$ symmetric polynomials in eigenvalues, whose Schur basis and product rule give the assertion. For the kernel use the ordinary Cauchy identity, first as a formal identity in eigenvalues. On Schur modules the operator norm of $S_\lambda(M)$ is at most $R_0^{|\lambda|}$, so its trace is bounded by $\dim S_\lambda(\C^d)R_0^{|\lambda|}$. The dimension Cauchy identity bounds the absolute series by $(1-rR_0^2)^{-d^2}$. Each summand is a rank-one Gram kernel, proving positivity by absolute summation.
\end{proof}
The term for the empty partition is a constant scalar feature, not an allowed module-valued tensor unit. In dimension two,
\[
 (\tr M)^2=\tr\operatorname{Sym}^2M+\det M,\qquad
 \tr\operatorname{Sym}^2M=\tfrac12((\tr M)^2+\tr M^2).
\]
More generally $\det(I-zM)=\sum_j(-z)^j\tr\bigwedge^jM$ and its reciprocal is $\sum_m z^m\tr\operatorname{Sym}^mM$. These are actual determinant/Schur connections, not identifications with an arithmetic $L$-function. On continuous finite-length $\Lambda$-modules the monodromy is nilpotent and all positive-degree characters vanish. The displayed Cauchy positivity is local representation-theoretic positivity, not the global Weil form.

\section{Lattice almost-split sequences and genuine diamond quotients}
For the interval dessin let $X=(x)$, $Y=(y)$ in $\Lambda=k[[x,y]]/(xy)$, $R=k[[u]]$, $u=x+y$. A lattice is a finitely generated $\Lambda$-module free over $R$.
\begin{theorem}[The nodal lattice calculation]\label{thm:nodear}
Every lattice is $\Lambda^r\oplus X^a\oplus Y^b$. The almost-split sequences in this lattice category are
\[
 0\to Y\lhook\joinrel\longrightarrow\Lambda\xrightarrow{\,\times x\,}X\to0,
 \qquad0\to X\lhook\joinrel\longrightarrow\Lambda\xrightarrow{\,\times y\,}Y\to0.
\]
In particular $\tau X=Y$, $\tau Y=X$. Multiplication by $u$ between two copies of either row, followed by reduction modulo $u$, gives an actual nonzero exact three-by-three diagram.
\end{theorem}
\begin{proof}
Normalize $M$ \emph{inside its total quotient module}: $\overline M=BM$, $B=k[[x]]\times k[[y]]$. The conductor gives $u\overline M\subset M\subset\overline M$. Both branch projections of $W=M/u\overline M\subset k^p\oplus k^q$ are onto. Intersect $W$ with the two factors; its remaining quotient identifies their two residual quotients. Separate bases for these pure factors and the common quotient put $W$ into pure first lines, pure second lines and diagonal pairs. Lift branch bases over the DVRs. These are respectively the summands $X,Y,\Lambda$. Their endomorphism rings are local.

Exactness of the first sequence is $\operatorname{ann}x=Y$. It does not split: a map $X\to\Lambda$ lands in $\operatorname{ann}y=X$, and its composite with $\times x$ is a nonunit in $\operatorname{End}X=k[[x]]$. Every map $\Lambda\to X$ factors through this epimorphism, $\operatorname{Hom}(Y,X)=0$, and every nonunit endomorphism of $X$ is divisible by $x$ and factors through $X\hookrightarrow\Lambda\to X$. Classification handles all finite sums; a unit $X$ component would be a retraction. This is right almost splitness. The dual argument for nonsections out of $Y$ proves left almost splitness. Endomorphisms fixing either structure map have constant term one, hence are units, proving minimality. Swap $x,y$ for the other sequence. Finally all three terms are $R$-free. Multiplication by $u$ is injective and reduction preserves exactness, proving the diagram assertion.
\end{proof}
The lattice AR quiver has arrows $Y\to\Lambda\to X$ and $X\to\Lambda\to Y$; the mesh composites are zero. Its projective vertex $\Lambda$ is not the endpoint of another nonsplit AR sequence. AR meshes do not mean that all rows and columns of any drawn grid are short exact.

For a second, explicit family fix $n\ge3$ and the rational Belyi map
\[
 \beta_n(z)=\frac1{1-T_n(z)^2}
 =\frac1{(1-z^2)U_{n-1}(z)^2},
\]
where $T_n,U_{n-1}$ are Chebyshev polynomials. It has passport $(2n),(2^n),(2^{n-1},1,1)$. Its contracted dual graph is a path, identifying its genus-zero dessin. The associated infinite algebra has vertices $1,\ldots,n$, endpoint loops $l,r$, opposite arrows $a_i:i\to i+1$, $b_i:i+1\to i$, and precisely the relations $l^2,r^2,a_ib_i,b_ia_i$. Its single permitted cycle is $l a_1\cdots a_{n-1}r b_{n-1}\cdots b_1$; the order rank is $(2n)^2$.

Retain one arrow at each edge in an orientation $\Omega$, and kill its opposite. The quotient is exactly the string algebra $H_n(\Omega)$ of Huang--Lin--Su \cite{HuangLinSu}; the killed-arrow ideal contains $J^{n+2}$, so it is also an actual quotient of $\Lambda_n$. Inflation identifies its module category with the full abelian quotient-annihilator subcategory $\mathcal C_\Omega$ of $\Lambda_n$-modules. This is generally not extension closed in the whole order category. The imported AR property is intrinsic to $\mathcal C_\Omega$, not automatically an ambient lattice or finite-length AR property. The extended edition gives nine explicit string modules, all twelve maps and the four exact meshes from the supplied diamond resource, now situated in this actual dessin quotient.

The other supplied resources also have a precise home. For the interval order the subcategory annihilated by $J^2$ is $\operatorname{Mod}B_0$, $B_0=\Lambda/J^2$. This is a domestic string quotient, not itself a gentle algebra. Inflation is full and preserves purity, and $\operatorname{Zg}B_0$ is a closed subspace of $\operatorname{Zg}\Lambda$. Ringel's infinite-word/pure-injective program and its proof by Puninski--Prest apply there \cite{Ringel,PuninskiPrest}; the isolated band $xy^{-1}$ gives Krull--Gabriel and relative Cantor--Bendixson dimension two by \cite[Theorem10.19]{LakingPrestPuninski}. This is not an ambient rank calculation for all $\Lambda$-modules. The parameter in $x=J_2,y=\lambda J_2$ survives even though all cycle traces vanish; the pp condition $\exists v\ (z=xv\ \wedge\ (y-\lambda x)v=0)$ detects it. These modules are $u$-torsion and are not nonzero order lattices. Infinite-string and locally finite AR theorems must retain their distinct finiteness hypotheses; a surviving permitted cycle has infinitely many paths at a vertex.

\subsection{Polynomial strings, completed lattices and the Fukaya category}
Crawley--Boevey's theorem \cite[Theorems1.1--1.4]{CrawleyBoeveyStrings} applies directly to the infinite polynomial dessin algebra: its finitely controlled modules decompose into strings and finite-dimensional bands; its pointwise artinian modules additionally allow primitive injective bands. For our finite quiver, these conditions mean finitely generated and artinian respectively. The words may be infinite, but their modules have algebraic direct-sum bases. This is not a classification of every module or of arbitrary formal completions.

In fact no nonzero completed-order lattice, restricted to $A$, meets either finiteness condition. The polynomial $A$ is finite over $k[u]$: remove full turns from sufficiently long paths as in Proposition~\ref{thm:order}. If an $R$-free lattice were finitely generated over $A$, it would be finite torsion free, hence free, over the PID $k[u]$. Multiplication by $1-u$ would then fail to be onto, whereas it is a unit over $R=k[[u]]$. Also $M\supsetneq uM\supsetneq u^2M\supsetneq\cdots$ excludes artinianity. For a concrete indecomposable outside the ordinary string/band list take $k[[x]]$ with $y=0$: its endomorphisms are multiplication by formal series, so its endomorphism ring is local. Its $x$ is injective and nonsurjective, while $1-x$ is invertible; these properties exclude every pure-$x$ ordinary string or band. The extended edition gives the full word argument.

The Fukaya connection requires equally precise categories. Lekili--Polishchuk explicitly allow infinite-dimensional locally gentle algebras, but their equivalence is for \emph{perfect dg modules} of a homologically smooth algebra, with specified grading, line field and stops \cite[Definition3.1.1, Lemma3.1.3, Theorem3.2.2]{LekiliPolishchuk}. Every nonempty closed dessin here has both permitted and forbidden cycles, so its polynomial algebra is neither proper nor smooth by their criterion. This does not exclude appropriate singular or completed variants; it prevents a direct use of that particular theorem. Infinite word, graded curve and geodesic are not interchangeable without a surface realization and boundary/grading data.

There is a genuine positive example inside our setting: the polynomial interval quotient $A/(y)=k[x]$ is still infinite dimensional, is smooth, and has no forbidden cycle. Its perfect category has the partially wrapped annulus model with one stop total. In degree zero $\operatorname{Perf}(k[x])=D^b(\operatorname{mod}_{\rm fg}k[x])$. Inflation of this quotient is fully faithful on ordinary modules, but not on derived categories: $\operatorname{Ext}_A^2(k[x],k[x])=k$, whereas $\operatorname{Ext}_{k[x]}^2(k[x],k[x])=0$. The former follows from the alternating $y,x$ resolution. Thus this quotient bridge does not identify the whole closed dessin's perfect category with the annulus category.

Finally, Leavitt path algebras \emph{can} fail invariant basis number, but need not. The permitted-state Laurent algebras here have finite-dimensional matrix quotients, and therefore invariant basis number. The polynomial and completed dessin algebras have their vertex-simple quotients, giving the same conclusion. The complex Leavitt algebra of $G_\alpha$ below is a dense *-subalgebra of its graph C*-algebra, on that specified graph; it is not the original gentle algebra with ordinary CK1 imposed.


\section{Leavitt relations, the constrained Fock algebra and spectrum}
Work now over $\C$. Ordinary Leavitt and graph C* constructions \cite{AbramsArandaPino,BPRS} cannot simply impose our original forbidden products: if $ab=0$ and CK1 says $a^*a=p_{t(a)}$, then $0=a^*ab=b$. Every arrow is a forbidden successor, hence every arrow and then every vertex is zero. CK2 is not needed. The ordinary graph-Toeplitz CK1 setting has the same collapse.

The valid permitted-arrow-state graph instead consists of the disjoint directed $\sigma$-cycles. Its Leavitt algebra is $\bigoplus_cM_{|c|}(k[z,z^{-1}])$, and its graph C* algebra is $\bigoplus_cM_{|c|}(C(\mathbb T))$. These forget the $\alpha$-gluing; both theta dessins give the same answer. They are not the completed original order. Nor is the ordinary Toeplitz algebra of a cycle of length $\ell>1$ equal to $M_\ell(\mathcal T)$: their $K_0$ ranks are $\ell$ and one respectively \cite{BrownloweLacaRobertsonSims}.

Let $\mathcal H=\ell^2(\mathcal P)$ for all permitted finite paths, including vacua $e_v$. Define $L_a\delta_p=\delta_{ap}$ when the prefix is permitted and composable, and zero otherwise. Let $P_v$ project onto paths starting at $v$, $Q_a=L_aL_a^*$ and $v_v$ onto the vacuum. Then
\[
 L_aL_{\phi(a)}=0,\quad
 L_a^*L_a=P_{t(a)}-Q_{\phi(a)}=v_{t(a)}+Q_{\tau(a)},\quad
 P_v-\sum_{s(a)=v}Q_a=v_v.
\]
These follow on basis vectors. The different initial projection avoids CK1 collapse. Write $\mathcal T_A=C^*(P_v,L_a)$, a constrained-path Fock C* algebra in the monomial/subproduct setting \cite{KakariadisShalit,ShalitSolel}. The polynomial $A$ acts faithfully: applying a finite path sum to the sum of vacua recovers its coefficients. This representation need not extend to all formal $\Lambda$ series: $\sum_{m\ge0}c^m$ along a cycle would send its vacuum to a nonsquare-summable sum.

\begin{theorem}[Compact extension and incidence index]\label{thm:fock}
Let $\mathcal H_v$ be the terminal-path sector. There is an exact sequence
\[
 0\longrightarrow\bigoplus_{v\in Q_0}\mathcal K(\mathcal H_v)
 \longrightarrow\mathcal T_A\longrightarrow
 \bigoplus_{c\in\operatorname{cyc}(\sigma)}M_{|c|}(C(\mathbb T))\longrightarrow0.
\]
With $\delta[u]=[1-W^*W]-[1-WW^*]$, its boundary map is
$\delta(c)=-\sum_{h\in c}e_{[h]}$, minus the unsigned incidence map $B_\Gamma$ of the cellular graph. Thus, for its $E$ edges,
\[
\begin{array}{c|cc}
 &K_0(\mathcal T_A)&K_1(\mathcal T_A)\\\hline
\Gamma\text{ bipartite}&\Z^{E+1}&\Z\\
\Gamma\text{ nonbipartite}&\Z^E\oplus\Z/2&0.
\end{array}
\]
\end{theorem}
\begin{proof}
All prefixes preserve the terminal sectors. For paths $p,q$ ending at $v$, $L_pv_vL_q^*=|\delta_p\rangle\langle\delta_q|$. Therefore the vacuum ideal is exactly the displayed sum of compacts: it contains these matrix units and that sum is itself an ideal. In the quotient, $L_a$ has initial projection $Q_{\tau(a)}$ and range $Q_a$, a graph CK family for the permitted-state cycles. The $Q_a$ have infinite rank and nonzero quotient images; the length gauge action descends. Gauge-invariant uniqueness \cite[Theorem2.1]{BPRS} identifies the quotient as claimed.

For a cycle $c$ based at arrow $a$, lift its corner unitary by $W=L_cQ_a+(1-Q_a)$. Then $W^*W=1$ and $1-WW^*=Q_a-L_cQ_aL_c^*$ projects onto the first $|c|$ paths starting with $a$. Their terminal vertices count exactly its incident darts, including twice for a loop. This proves the index formula. The six-term exact sequence gives $K_1=\ker B_\Gamma$ and $0\to\operatorname{coker}B_\Gamma\to K_0\to\Z^{\#\operatorname{cyc}\sigma}\to0$. The free last term splits abstractly. Spanning-tree integer elimination gives Smith entries $1$ repeated $V-1$ times and one zero in the bipartite case; an odd cycle instead adds the entry two. This proves the groups.
\end{proof}
The nodal interval has $K_0=\Z^2$, $K_1=\Z$; the single-loop sphere has $\Z\oplus\Z/2$, zero. Both theta dessins have $\Z^4$, $\Z$. Coarse $K$-theory does not recover genus or the arithmetic factors.

\subsection{A graph presentation and an actual pullback}
The Fock algebra itself has a graph presentation, on a different graph from the original gentle quiver. Let $G_\alpha$ have vertices $q_h$ for $h\in H$ and sinks $b_e$ for $e\in Q_0$, and edges
\[
 s_h:q_h\to q_{\tau h},\qquad t_h:q_h\to b_{[\tau h]}.
\]
\begin{proposition}\label{thm:graphmodel}
$C^*(G_\alpha)\cong\mathcal T_A$, by $q_h\mapsto Q_{a_h}$, $b_e\mapsto v_e$, $s_h\mapsto L_{a_h}Q_{a_{\tau h}}$ and $t_h\mapsto L_{a_h}v_{[\tau h]}$. The integer graph relation at $h$ is $q_h-q_{\tau h}-b_{[\tau h]}$. Eliminating these relations recovers $K_1=\ker B_\Gamma$ and $K_0\cong\operatorname{coker}B_\Gamma\oplus\Z^{\#\operatorname{cyc}\sigma}$.
\end{proposition}
\begin{proof}
The two edges at $q_h$ have orthogonal initial pieces and orthogonal ranges summing to $Q_{a_h}$, by the Fock initial-projection identity. All vertex projections are nonzero. The length gauge acts with degree one on both edge families, so gauge-invariant uniqueness applies. They generate since $L_{a_h}=s_h+t_h$ and $P_e=v_e+\sum_{[h]=e}Q_{a_h}$. In the integer matrix a kernel coefficient is constant around each permitted cycle; its sink equations are exactly $B_\Gamma c=0$. In the cokernel, eliminate all but one state per cycle; its remaining relation is the sum of its incident sink generators. The remaining cycle generators are free, proving the stated groups.
\end{proof}
For the nodal interval, $G_\alpha$ has two loop vertices, each with an exit to the common vacuum sink. Chirvasitu--Hajac--Tobolski's actual pullback theorem \cite[Theorem3.4]{GraphPullbacks} gives
\[
 \{T\in\mathcal T_A:\pi(T)\in\C^2\}
 \cong C^*(E_1),\qquad \C^2\hookrightarrow C(\mathbb T)^2\text{ by constants},
\]
where $E_1$ has the same three vertices and countably many arrows from each loop vertex to the sink, with no loops. The arrow numbered $n\ge0$ maps to loop$^n$ followed by exit. This algebra is AF, with $K_0=\Z^3$, $K_1=0$. The extended edition verifies every hypothesis. The theorem does not turn arbitrary cycle-vertex identifications into pullbacks: its first graph is acyclic and its path functor must cover all pointed paths. Its sink-extension Theorem4.3 has further injection and path-coverage requirements.

\subsection{A separate diagonal-fiber C* completion retaining the pairing}
The formal order's residue gluing cannot be copied by taking diagonals of arbitrary full matrices: $\operatorname{diag}(E_{12}E_{21})\ne\operatorname{diag}(E_{12})\operatorname{diag}(E_{21})$. For $\ell>1$ there is no nonzero *-homomorphism from $M_\ell(C(\mathbb T))$ to a commutative C*-algebra, as restriction to its constant matrix units shows. Restrict the distinguished fiber first. Define
\[
 D_\ell=\{f\in C(\mathbb T,M_\ell(\C)):f(1)\text{ is diagonal}\},\qquad
 \mathcal D_H=\Bigl(\prod_cD_{|c|}\Bigr)\mathop{\times}_{\C^H}\C^{Q_0},
\]
where the first map evaluates diagonal entries at one and the second repeats $\lambda_e$ at the two darts paired by $\alpha$.
\begin{theorem}[A faithful polynomial pullback model]\label{thm:diagonalmodel}
There is a split exact sequence
\[
 0\to\bigoplus_c C_0(\mathbb T\setminus\{1\},M_{|c|})
 \to\mathcal D_H\to\C^{Q_0}\to0,
\]
so $K_0(\mathcal D_H)=\Z^{|Q_0|}$ and $K_1(\mathcal D_H)=\Z^{\#\operatorname{cyc}\sigma}$. The polynomial dessin algebra embeds faithfully and generates this C*-algebra with its adjoints. This embedding does not extend to a unital algebra map from all of $\Lambda$.
\end{theorem}
\begin{proof}
The section consists of the constant paired diagonal projections. The kernel is the displayed vanishing fiber, a sum of matrix suspensions, giving the groups by split exactness and suspension. For an explicit embedding set $g(z)=|1-z|^2/4$, choose one cut arrow per permitted cycle, and send $a_h$ to $g(z)z^{\epsilon_h}E_{h,\tau h}$ in its branch, where $\epsilon_h=1$ at the cut and zero otherwise. Send $e_e$ to the constant paired diagonal projection. Every forbidden product is zero. A full cycle at $h$ maps to $w_c(z)E_{hh}$, $w_c=g^{|c|}z$. This scalar function is injective on the circle: away from its unique zero at one its argument recovers $z$. Continuous functional calculus therefore supplies every $f(z)E_{hh}$ with $f(1)=0$. Multiplying by path images and approximating division away from one supplies every vanishing off-diagonal entry. Together with the constant paired diagonals these generate $\mathcal D_H$. Path normal forms in a fixed entry are a nonzero prefactor times a polynomial in $w_c$; its infinite range proves independence. Evaluation at one first separates the zero-length coefficients. Finally for $z_0\ne1$, $1-w_c(z_0)^{-1}c_h$ is a unit in the formal order, but its proposed image is singular at $z_0$. A unital map preserves units, proving nonextension.
\end{proof}
This is the C* algebra generated by a faithful polynomial representation and its adjoints, distinct from the completed order and from $\mathcal T_A$. The norm closure of the polynomial image alone need not contain those adjoints. Their $K$ groups already differ. The pairing belongs to the labelled presentation; an unmarked isomorphism class need not recover the ribbon ordering. No arithmetic descent, Frobenius action or completed-$L$ positive determinant is inferred from this additional valid gluing model.


\begin{proposition}[The degree spectral triple]\label{thm:degree}
Define $N\delta_p=|p|\delta_p$ on $\{\xi:\sum_p|p|^2|\xi_p|^2<\infty\}$. With the polynomial creation/annihilation *-algebra this is an ungraded spectral triple, $p$-summable exactly for $p>1$. With $N_+^{-s}$ zero on the kernel,
\[
 \operatorname{Tr}e^{-tN}=E+\frac{2E}{e^t-1},\quad
 \operatorname{Tr}N_+^{-s}=2E\zeta(s),\quad
 \operatorname{Tr}\bigl(\chi(N)N_+^{-s}\bigr)=2E L(s,\chi)
\]
for $t>0$ and $\Rea s>1$ in the last two formulas, where $\chi$ is an externally chosen Dirichlet character.
\end{proposition}
\begin{proof}
There are $E$ vacua and $2E$ paths of every positive length: each starting arrow has a unique continuation. The real diagonal operator is self-adjoint on the stated maximal domain; finite supports are a core and its resolvent is compact. Prefixes preserve the domain and $[N,L_a]=L_a$, $[N,L_a^*]=-L_a^*$, so polynomial commutators are bounded. The trace formulas and summability reduce to the corresponding scalar series. The character multiplier is bounded; membership in $\mathcal T_A$ is not asserted.
\end{proof}
Length parity anticommutes with the creators but commutes with $N$, so it is not an even-spectral-triple grading. The ordinary odd Fredholm class of this positive $N$ is trivial: $N(1+N^2)^{-1/2}-I$ is compact. The spectrum depends only on $E$; zeros of the scalar zeta trace are not eigenvalues of $N$.

\section{Arithmetic attachment, descent and the remaining comparison}
The compact ideal gives a precise home for an arithmetic determinant, with no claim that the arithmetic is derived from the dessin. Fix one terminal sector. Its two paths of each length $n\ge1$ give orthonormal vectors $\xi_n=(\delta_{p_{n,1}}+\delta_{p_{n,2}})/\sqrt2$ and an isometry $U:\ell^2(\mathbb N_{\ge1})\to\mathcal H_v$.

For a nonprincipal character $\chi$ put $A_n=\sum_{j\le n}\chi(j)$, $B_0=\sup|A_n|\le q_\chi$, and, on $2^k\le n<2^{k+1}$, set
\[
 w_n=\sqrt n(k+1),\quad u_n=A_n/w_n,\quad
 v_n(s)=w_n(n^{-s}-(n+1)^{-s}),\quad
 \varphi_s(z)=\sum_n v_n(s)z_n-z_1.
\]
Here $u_1=1$. For $\Rea s>0$, $u,v(s)\in\ell^2$ and the pairing is complex linear, with no conjugation of $s$. Define
\[
 \widetilde K_s=U(u\otimes\varphi_s)U^*,\qquad
 \det_F(I+\widetilde K_s)=L(s,\chi).
\]
Indeed dyadic blocks give $\|u\|^2\le2B_0^2$, and on a compact set $\Rea s\ge\sigma_0>0$, $|s|\le M$, the squared $v$ tail is bounded by $M^2\sum_{k\ge J}(k+1)^2 2^{-2\sigma_0k}$. Thus the rank-one family is trace-norm holomorphic. Its determinant is $1+\varphi_s(u)$, equal by Abel summation to $L$. Replacing $v$ by $P_Nv$ and adding $A_N(N+1)^{-s}e_1^*$ to the functional gives exactly every raw partial sum $S_N$, with compact-uniform trace-norm convergence; $\|UXU^*\|_1=\|X\|_1$. For primitive $\chi$, completion replaces the functional by $g(s)\sum v_n(s)z_n-z_1$, with the standard nonzero gamma factor $g$, rather than multiplying the whole old functional. The extended edition supplies the quantitative estimates.

This inserts a scalar encoding into one actual dessin compact ideal; it is not known to be $A$-linear, a derived chain map or Frobenius compatible. A common fixed positive metric cannot make this same fixed-range family symmetric for all critical heights: symmetry forces its functionals to be proportional, whereas the coefficient ratio at heights $0,t\ne0$ is asymptotic to $(\sigma+it)n^{-it}/\sigma$, with two opposite subsequential limits. Different positive realizations remain possible.

\begin{proposition}[Complex dessin data do not determine arithmetic descent]\label{thm:descent}
The rational Belyi pairs $E_d:y^2=x^3-d$, $\beta_d=x^3/d$, for $d=1,2$, are isomorphic over $\C$ but have distinct Frobenius polynomials at the good prime seven: $1-4T+7T^2$ and $1-T+7T^2$.
\end{proposition}
\begin{proof}
Choose $a^3=d$, $b^2=d$. The projective isomorphism $(x,y)\mapsto(ax,by)$ from $E_1$ to $E_d$ preserves the Belyi map. The degree-six clean passports are the elliptic theta passport already computed. Directly counting $y$ for $x=0,\ldots,6$ gives respectively $(0,1,1,0,1,0,0)$ and $(0,0,0,2,0,2,2)$. Including infinity yields four and seven points, hence $a_7=8-\#E_d(\mathbb F_7)=4,1$. Both discriminants $-432d^2$ are nonzero modulo seven.
\end{proof}
Any invariant of the complex dessin or its undecorated order/Fock algebra alone is the same for these two inputs, so cannot recover both factors without descent data. This does not obstruct a theory including a specified rational model and compatible Galois action. Elliptic $H^1$ has weight one and unshifted center $s=1$, unlike Dirichlet center $1/2$; the example is not a counterexample to GRH.

\subsection{A constructive arithmetic benchmark on the theta dessin}
The preceding twists can be distinguished constructively once their finite-field descent data are retained. This closes the elementary local-factor computation, not the global positive-determinant problem. Let $\kappa=\mathbb F_q$, $\operatorname{char}\kappa\ne2,3$ and $d\in\kappa^*$. Put $E_d:y^2=x^3-d$ and $a_q(d)=q+1-\#E_d(\kappa)$. For $q\equiv1\pmod3$, fix a nontrivial cubic character $\rho$ and the quadratic character $\eta$, both extended by zero at zero. Define
\[
 J=\sum_{u\in\kappa}\rho(u)\eta(1-u),\qquad
 j=-\eta(-1)J,\qquad \lambda_d=\eta(d)\rho(d),\quad \pi_d=\lambda_dj.
\]
\begin{theorem}[Uniform local factor]\label{thm:jacobitheta}
For every such prime power and nonzero $d$,
\[
 a_q(d)=\pi_d+\overline{\pi_d},\quad |\pi_d|^2=q,\quad
 P_d(T):=1-a_q(d)T+qT^2=(1-\pi_dT)(1-\overline{\pi_d}T).
\]
If $q\equiv2\pmod3$, instead $a_q(d)=0$. The split-case values lie in $\Z[\omega]$, $\omega^2+\omega+1=0$.
\end{theorem}
\begin{proof}
Cube fibers have size $1+\rho(t)+\overline{\rho(t)}$, including at zero; the first term is the literal constant one, not the zero-extended trivial character. Counting square fibers gives $a_q(d)=-\sum_x\eta(x^3-d)$. Grouping by $t=x^3$ removes the constant term since $\sum_t\eta(t-d)=0$. Substitution $t=du$ gives $\sum_t\rho(t)\eta(t-d)=\rho(d)\eta(-d)J$; the other term is its conjugate. This proves the signed trace formula.

For a nontrivial additive character $\psi$, write $G(\chi)=\sum_x\chi(x)\psi(x)$. Character orthogonality gives $|G(\chi)|^2=q$ for nontrivial $\chi$: substitute $x=uy$ in its squared norm; the inner sum over $y\ne0$ is $q-1$ at $u=1$ and $-1$ otherwise. Grouping $G(\rho)G(\eta)$ by $x+y$ gives $JG(\rho\eta)$, the zero-sum fiber vanishing since $\rho\eta$ is nontrivial. The three Gauss norms yield $|J|^2=q$. Neither purity nor RH was assumed. The remaining signs and character values have modulus one. For $q\equiv2\pmod3$ cubing is bijective, so the quadratic-character sum is zero. Finally all split-case summands are signs times cube roots of unity. These are the classical Gauss--Jacobi methods \cite[pp.499--502]{WeilJacobi}, with our normalization proved explicitly.
\end{proof}
At $q=7$, choose $\rho(3)=\omega$. Then $j=3+2\omega$, $\lambda_2=\omega^2$, and $\pi_2=-1-3\omega$, giving the two factors already counted. At $q=13$, $\rho(2)=\omega$ gives $J=1+4\omega$ but $j=-1-4\omega$; this detects the essential sign hidden by the first example.

\subsection{A genuine cellular carrier, and what still fails for the order}
Assume $q\equiv1\pmod3$ in this cellular construction. The elliptic theta dessin has cellular chains $C_2=\mathbb{Q}$, $C_1=\mathbb{Q}^3$, $C_0=\mathbb{Q}^2$. Orient the three edges from $v_+$ to $v_-$; then $\partial_1(e_i)=v_--v_+$ and $\partial_2=0$. Thus $H_1$ is the sum-zero plane, with basis $e_0-e_1,e_1-e_2$. Rotation $\sigma$ acts there by
\[
 C=\begin{pmatrix}0&-1\\1&-1\end{pmatrix},\qquad
 G=\begin{pmatrix}2&-1\\-1&2\end{pmatrix}.
\]
$G$ is the actual Euclidean edge Gram matrix in this basis. The deck involution $\alpha$ acts as $-I$. Consequently the six-element deck group $\langle\sigma,\alpha\rangle$ acts on this homology and permutes the original quiver and its entire projective resolution.

Write $\pi_d=A_d+B_d\omega$. Its multiplication action on this \emph{actual cellular} plane is
\[
 M_d=A_dI+B_dC=\begin{pmatrix}A_d&-B_d\\B_d&A_d-B_d\end{pmatrix},\qquad
 \det(I-TM_d)=P_d(T),\qquad M_d^tGM_d=qG.
\]
This constructs the local matrix from finite-field character sums and the dessin's deck action, rather than fitting its trace. Positivity follows from $2x^2-2xy+2y^2=(x-y)^2+x^2+y^2$. The normalized matrix is $G$-unitary, not asserted self-adjoint. The extended edition supplies a rational cellular chain lift, its exact determinant cancellation, and Kummer descent compatibility on homology. Its odd-over-even determinant ratio is the classical finite-field curve zeta function
\[
 Z(E_d/\kappa,T)=\frac{P_d(T)}{(1-T)(1-qT)}.
\]
Identification with the geometric $\ell$-adic Frobenius polynomial uses the standard elliptic trace and determinant theorem; a canonical integral comparison or an order-derived Frobenius functor has not been constructed here.

There is also a precise obstruction to extending this same real two-dimensional carrier to $q\equiv2\pmod3$. Frobenius then conjugates the cubic deck action. The relation $MC=C^2M$ forces $M=\left(\begin{smallmatrix}a&b\\a+b&-a\end{smallmatrix}\right)$, whose determinant is $-a^2-ab-b^2\le0$, contradicting the required positive determinant $q$. This excludes that real carrier with its fixed deck action, not complex or enlarged Hilbert--P\'olya constructions. The extended edition gives a compatible complex local replacement.

For the following insertion test use the characteristic-zero theta order $\Lambda_{\mathbb{Q}}$ over $R_{\mathbb{Q}}=\mathbb{Q}[[u]]$, with residue simple $S_e=\mathbb{Q}$. The limitation is concrete. The deck action has finite order and cannot alone supply the norm-$\sqrt q$ Jacobi coefficient. Inserting $M_d$ into $S_e\otimes\mathbb{Q}^2$ gives an actual residue-module endomorphism with determinant $P_d$, but that simple is not perfect over the order and choosing it breaks full deck symmetry. Inserting the same matrix $\Lambda_{\mathbb{Q}}$-linearly into the regular projective instead gives the $R_{\mathbb{Q}}$-determinant $P_d^{18}$, or $P_d^6$ on $e\Lambda_{\mathbb{Q}}\otimes\mathbb{Q}^2$. Taking a root is not an established determinant-line comparison; a deck-semilinear lift is a different map. These surviving multiplicity/equivariance obligations prevent promoting the local computation to the required arithmetic order functor.

For the global comparison fix a nonzero rational $d$ and use its good reductions. The finite-field curve zeta has center $1/2$ after $T=q^{-s}$. The \emph{global elliptic} $L$-function over $\mathbb{Q}$ has unshifted center one, whereas Dirichlet $L$-functions have center $1/2$. Local weight-one purity does not determine either global zero set. The elementary local bound proves absolute convergence of the elliptic Euler product on $\Rea s>3/2$; the new local identities give no determinant-preserving continuation with the desired global positive spectrum.


\subsection{Arbitrary projective actions: an obstruction and a genuine determinant repair}
The preceding tensor test has a stronger replacement. Keep the characteristic-zero completed elliptic theta order $\Lambda$ over $R=\mathbb Q[[u]]$, and write $K=\mathbb Q((u))$. All determinants in the next theorem are \emph{ordinary underlying-$R$ determinants}; they are not reduced norms.
\begin{theorem}[Projective determinant multiplicity]\label{thm:theta-multiplicity}
Let $P$ be any finitely generated projective right $\Lambda$-module and $f\in\operatorname{End}_\Lambda(P)$. Then $P\simeq\bigoplus_{i=0}^2(e_i\Lambda)^{n_i}$. Let $A_i\in M_{n_i}(\mathbb Q)$ be the induced maps on its three top multiplicity spaces. Then
\[
 \left.\det_R(1-Tf)\right|_{u=0}
 =\left(\prod_{i=0}^2\det(1-TA_i)\right)^6.
\]
For any bounded complex of these projectives and any $\Lambda$-linear chain endomorphism, its alternating determinant ratio specializes to a sixth power in $\mathbb Q(T)$. Generically, the ordinary ratio is a cube in $K(T)$.
\end{theorem}
\begin{proof}
The arrow ideal is the Jacobson radical: its elements admit geometric-series inverses after subtracting from one, and the quotient is $\mathbb Q^3$. A projective's top lifts to a map from the displayed sum; Nakayama gives surjectivity, projectivity splits it, and another application of Nakayama kills the kernel.

Put $B=\Lambda/u\Lambda$. Each $e_iB$ has path layers of dimensions $1,2,2,1$: the identity, two length-one paths, two length-two paths, and one length-three path. The two full cycles based at $i$ add to $ue_i$, so become negatives; all length-four paths vanish. A positive-length entry of a projective homomorphism strictly raises path length. Consequently the associated graded map has six copies of $A=\operatorname{diag}(A_0,A_1,A_2)$, giving the formula by triangular block determinants. Base change of determinants justifies specialization. Multiplicativity along filtrations is ordinary linear algebra, consistent with \cite[Tags0FJ9,0GSX]{StacksDeterminants}; no purity assumption enters.

Each term's determinant has constant term one in $T$, so every denominator has nonzero specialization. Taking the finite alternating product therefore preserves the sixth power, even after cancellation. Finally $\Lambda_K=M_3(K)\times M_3(K)$; a projective's branch modules are sums of three-dimensional simples. A commuting endomorphism acts on each multiplicity space, repeating every factor three times. This proves the generic assertion.
\end{proof}
Thus no such ordinary perfect-complex ratio can equal the constant-in-$u$ quadratic $P_d(T)$, or the full local curve zeta $P_d(T)/((1-T)(1-qT))$. For $P_d$, the valuation at infinity is $-2$, not a multiple of six. For the zeta ratio, $P_d(1)=q+1-a_q>0$ by the local bound $|a_q|\le2\sqrt q<q+1$, so its pole at $T=1$ is simple. This also covers square prime powers where the numerator may have a double root. It is an obstruction for the specified ordinary determinant and linear category, not all spectral or determinant theories.

There is nevertheless a \emph{genuine integral functor} removing this multiplicity. Choose a permitted branch and its cycle cut. Its projection $\Lambda\to B_+\subset M_3(R)$ makes the column lattice $W_+=R^3$ a left $\Lambda$-module. Define
\[
 \Phi_+ :\operatorname{Perf}(\Lambda)\longrightarrow\operatorname{Perf}(R),
 \qquad P^\bullet\longmapsto P^\bullet\otimes_\Lambda W_+.
\]
Bounded projectives compute this derived tensor functor; flatness of $W_+$ on all modules is unnecessary. Since $e_iW_+\simeq R$, its determinant on $e_i\Lambda\otimes_{\mathbb Q}V$ with external action $1\otimes M_d$ is exactly $P_d$, whereas on the regular module it is $P_d^3$. This is a defined change of determinant functor, not an arbitrary polynomial root. It does not construct the arithmetic action independently: the already computed $M_d$ is still the input.

The price of this repair is precise. Branch swap $\alpha$ exchanges $W_+$ and $W_-$; pairing them yields $P_d^2$ on that basic-projective input. The chosen branch determinant cannot intertwine $\alpha$ on all endomorphisms: the corner unit $e_i+c_i^+$ acting on $e_i\Lambda$, with $c_i^+$ a plus full cycle, acts on its determinant by $1+u$, whereas its $\alpha$-image acts by one. Also no rank-one-per-vertex column lattice supported on one branch admits an $R$-linear lift of the rotation $\sigma$ fixing $u$. Its three nonzero arrow scalars have nonnegative $u$-valuations summing to one; an invertible cyclic lift would make all three equal. This contradiction concerns that lattice, not every determinant-level rotation comparison. Ramification $u=v^3$ \emph{and a new lattice} with all three arrow scalars $v$ admits such a rotation, changing the stipulated base/lattice data. A descent-compatible arithmetic determinant retaining the required symmetry and a global positive realization remains open.


\subsection{Tame descent: a surviving order and a sharper obstruction}
We now test the proposed ramified repair rather than assume that its rotation descends. Extend coefficients explicitly to $k=\mathbb Q(\omega)$, $\omega^2+\omega+1=0$. Put $R=k[[u]]$, $S=k[[v]]$, $u=v^3$, and let $g(v)=\omega v$, fixing $k$. On the balanced plus column $W=S^3$, let $C\epsilon_i=\epsilon_{i+1}$ and
\[
 \Theta=Dg,\qquad D=\operatorname{diag}(1,\omega^2,\omega).
\]
The plus arrows act by $vE_{i,i-1}$, and minus arrows by zero. Direct calculation gives $\Theta^3=1$, compatibility of $\Theta$ with this $\Lambda_R$-action, and
\[
 \Theta C=\omega^2C\Theta,\qquad
 L=W^\Theta=R\epsilon_0\oplus Rv\epsilon_1\oplus Rv^2\epsilon_2.
\]
Thus invariants restore precisely the original cut lattice. The extension $R\to S$ is finite \emph{free} of rank three and faithfully flat; the group quotient is not a torsor at $v=0$. The comparison $S\otimes_R L\to W$ has matrix $\operatorname{diag}(1,v,v^2)$, determinant $u$ and cokernel $S/(v)\oplus S/(v^2)$. Hence this semilinear group action alone is not effective descent of the column to $R$. This is the explicit local root-stack situation, consistent with the trivial-stabilizer criterion for vector-bundle descent \cite[Theorem10.3]{AlperDescent}; no stack theorem is needed for the displayed calculation.

Two positive constructions survive. The determinant line $\bigwedge_S^3W$ has trivial $D$ and $C$ actions and does descend, but the natural map
\[
 \det_R(L)\longrightarrow(\det_S W)^\Theta
\]
has image $uR$ inside $R$, not an isomorphism. Determinant and coarse invariants cannot be interchanged while dropping this nonunit comparison. Also the \emph{carrier} endomorphism algebra $\operatorname{End}_S(W)^\Theta$, conjugated to the invariant basis, is exactly the hereditary order with entries $R$ on and below the diagonal and $uR$ above. Its rotation descends as $\operatorname{Ad}(A)$, where
\[
 A=\begin{pmatrix}0&0&u\\1&0&0\\0&1&0\end{pmatrix},\qquad A^3=uI.
\]
The algebra automorphism has order three although $A$ is not an integral unit. This is not $\operatorname{End}_\Lambda(W)$ or a descended arithmetic cohomology group.

The natural arithmetic attempt fails: $aI+bC$ does not commute with $\Theta$ when $b\ne0$, as in both prime-seven Jacobi examples. The original cellular sum-zero plane is not $\Theta$-stable either. Replacing $C$ by $vC$ does give the descended operator $A$ on $L$, but
\[
 \det(1-T(aI+bA))=(1-aT)^3-ub^3T^3,
\]
whose special fiber is $(1-aT)^3$, not the required local quadratic. These carrier maps are not automatically $\Lambda$-linear arithmetic endomorphisms.

There is a stronger category-level conclusion.
\begin{theorem}[Full-deck coefficient kernels]\label{thm:full-deck-kernel}
Let $W_0$ be a finite $R$-free left $\Lambda_R$-module on which the central $u$ acts by the given scalar. Suppose it has invertible $R$-linear intertwiners lifting the order automorphisms $\sigma$ and $\alpha$. Then, for some $m\ge0$, every $e_iW_0$ has rank $6m$. For any finite projective $P$ and $\Lambda_R$-linear $f:P\to P$, let $A_i$ be its top blocks. The genuine coefficient functor $\Phi_{W_0}= -\otimes_\Lambda^{\mathbf L} W_0$ satisfies
\[
 \left.\det_R(1-T\Phi_{W_0}(f))\right|_{u=0}
 =\left(\prod_i\det_k(1-TA_i)\right)^{6m}.
\]
Bounded perfect ratios retain this exponent. Thus this whole class of full-deck, fixed-center finite free coefficient kernels cannot supply either the single local quadratic or the full local curve zeta.
\end{theorem}
\begin{proof}
In the generic fiber intersect $W_0$ with its two central branch summands to obtain finite free integral lattices $W_\pm$; an integral direct-sum decomposition is not assumed. Rotation preserves them and branch swap exchanges them. On a branch let the common vertex rank be $r$. The three arrow matrices have cyclic product $uI_r$. Their determinant valuations are nonnegative, sum to $r$, and rotation intertwiners make them equal. Therefore $3\mid r$. Branch swap makes the two ranks equal, so the full vertex rank is $6m$.

For the determinant formula, filter $W_0/uW_0$ by the arrow radical, whose fourth power vanishes. A positive-length entry of $f$ raises the filtration; a top block $A_i$ repeats across its $i$-th graded pieces a total of $\dim_k e_i(W_0/uW_0)=6m$ times. Triangular determinants give the assertion. Finite alternating products preserve it. Quadratic degree and the uncancelled simple pole of the full local zeta exclude the targets as before.
\end{proof}
This theorem is sharp: $W_0=\Lambda_R$ has vertex rank six. It requires neither left projectivity nor left flatness of a general kernel. It excludes this coefficient category, not arbitrary semilinear correspondences, twisted determinant lines or infinite spectral constructions. The extended edition identifies the full character-orbit repair with the hereditary normalization and explains why Schur commutator cancellation is insufficient. All determinant multiplicities here are relative to the stated $k$-valued base; restricting a rational quadratic factor to the former $\mathbb Q$-base squares it. No arithmetic descent or global positive determinant has been recovered by this ramification.


The recent NCG literature supplies concrete comparison targets, with strong hypotheses. Connes--Consani construct finite abelian covers of the adelic Riemann sector with prime-orbit length $\log p$ and arithmetic Frobenius monodromy \cite[Theorem3.12]{ConnesConsaniKnots}. The prolate/local-factor and moment constructions \cite{Prolate,LocalMoments} give particular finite-prime analytic models, not a uniform positive comparison for all completed $L$-functions. The zeta-spectral-triple determinant theorem \cite[Theorem5.10]{ZetaTriples} assumes a specified restricted Weil form, simple even bottom eigenfunction and boundary normalization. Its infinite comparison is still an additional step. The distributional real-zero theorem \cite[Theorem6.1]{QuadraticForms} concerns a particular interval quadratic form, essential self-adjointness on its trigonometric core and an even simple ground state; a Fock vacuum alone verifies none of those arithmetic identifications. Deninger's periodic-packet construction and real-cohomology obstruction \cite{DeningerPackets,DeningerReal} further specify, rather than remove, these issues. Full theorem-level reading coverage and exact versions are recorded alongside this draft.

A complete approach must construct arithmetic descent and prime actions on actual order/derived or analytic objects, prove determinant-line and local-factor compatibility, control regularization and conductor/gamma normalization, and establish a positive realization for that \emph{same} function. Finite-field cohomological trace formulas and Deligne purity supply the classical model \cite{DeligneWeilI,DeligneWeilII}; alternating duality alone is not positivity or purity. In a number-field construction these are obligations, not imported conclusions.

For a sharply testable final analytic target, put $F_\chi(z)=\Lambda(1/2+iz,\chi)$. Suppose entire functions $h_m(z)\det(zI-H_m)$, with $H_m$ finite self-adjoint and $h_m$ entire nowhere zero, converge locally uniformly on $\C$ to this nonzero completed function. Hurwitz on every disk disjoint from $\R$ then proves all its zeros real, hence the Dirichlet critical-line assertion. Establishing this convergence with independently constructed $H_m$ is exactly the missing load-bearing claim; a rank-one encoding or a reformulation does not establish it. Hecke, Dedekind and Artin extensions must state their analytic scope anew.

\subsection{Whole components, derived endpoints and determinant relations}
The full node lattice category has precisely the three indecomposables $X,\Lambda,Y$ and arrows $X\to\Lambda\to Y$, $Y\to\Lambda\to X$. Its stable category has only $X,Y$, no irreducible arrows, and suspension exchanging them; its AR triangles have zero middle term. For the polynomial node at a smooth support $x=\lambda+t$, $y=0$, $\lambda\ne0$, the modules $T_r=k[t]/t^r$ form an ambient homogeneous tube:
\[
 0\longrightarrow T_r\xrightarrow{(\pi,t)}T_{r-1}\oplus T_{r+1}
 \xrightarrow{(t,-\pi)}T_r\longrightarrow0,\quad T_0=0,\quad\tau T_r=T_r.
\]
The shift changes the derived placement, not the quasi-length. These nonzero-central-support objects do not extend to continuous finite-length modules of the arrow-adic completion. At its singular point even the candidate pure-branch tube fails ambient almost-splitness: $T_2^y\to k$ cannot factor through $T_2^x\to k$, since every intervening map lands in the $x$-socle, killed by the last map.

The nine HLS strings above sit in an entire intrinsic quotient component $T_{p,q}$, $(p,q)\in\Z^2$, with arrows increasing one coordinate and $\tau T_{p,q}=T_{p-1,q-1}$ \cite{HuangLinSu}. Its central object $E=M(ab)$ has infinite projective dimension: restriction to the first corner gives the dual-number simple with its infinite loop resolution, while restricted projectives are free or zero. Its module AR sequence therefore yields a distinguished derived triangle but not a derived AR triangle ending at $E$. Ringel's four-quarter wind-wheel tile \cite{Ringel}, with infinite-word diagonals and mixed center, is a different category's quilt; it is not silently identified with that HLS component or the Belyi surface.

There is no general equation ``AR mesh equals polynomial semi-invariant syzygy.'' The intrinsic $k(1\to2)$ quotient of the infinite Chebyshev algebra has AR row $0\to S_2\to E\to S_1\to0$ and a free presentation determinant $z$ on dimension $(1,1)$, so its semi-invariant ring $k[z]$ has no generator relation. A genuine framed elliptic-square stratum has $Z=[x_1,x_2,y_1,y_2]\in\operatorname{Mat}_{2\times4}$ with six minors and relation
\[
 p_{12}p_{34}-p_{13}p_{24}+p_{14}p_{23}=0.
\]
These minors have the same determinant character for the target $\mathrm{GL}_2$ with framed source; they are not all semi-invariants for the full unframed vertex group. In contrast, the nodal lattice AR row gives a genuine relative determinant identity over $R$: multiplication by $a+bx$ on the basis $(y,1)$ has matrix
$\bigl(\begin{smallmatrix}a&-b\\0&a+ub\end{smallmatrix}\bigr)$ and characteristic factor $(1-Ta)(1-T(a+ub))$. The endpoints have infinite order-projective dimension but are $R$-perfect.

Finally, the three-vertex theta algebra has no unital surjection onto a six-vertex radical-square-zero circular algebra: its semisimple image has at most three primitive idempotent images. Duplicating face vertex spaces instead gives an exact faithful unfolding which is not full; it does not inherit quotient or AR-transfer claims. Zwara Hom-control means that the Hom-dimension defect is bilinear on $K_0$, not merely a covering direct-sum formula \cite{ZwaraCorrected,RiedtmannZwaraCorrected}. The book supplies explicit fixed-dimension counterexamples, corrected circular-complex incidence resolutions, complete drawn component windows and all category hypotheses. These additions do not supply the arithmetic determinant comparison.


\section{A genuine theta covering and its categorical limits}
\label{cover:section}
Use column matrices, hence covariant quiver representations (left modules);
opposite categories translate the right-module convention in
\cite{CoveringLecture2026}. Let $A$ be the infinite elliptic theta dessin
algebra: vertices $i\in\Z/3$, arrows $a_i,b_i:i\to i+1$, and all mixed
length-two paths zero. For the arrow ideal $J$ and $L\geq2$ put $B_L=A/J^L$.
This is a finite string quotient of our actual infinite gentle algebra;
for $L>2$ its length-$L$ relations prevent treating it as quadratic gentle.

\begin{proposition}[A locally bounded integer cover]\label{cover:hom}
The category $\widetilde B_L$ with vertices $r\in\Z$, two arrows
$a_r,b_r:r\to r+1$, and the same mixed and length-$L$ zero relations,
has orbit $B_L$ under $r\mapsto r+3$. For finite-support representations,
push-down $F$ is exact and
\[
 \operatorname{Hom}_{B_L}(FM,FN)
 \simeq\bigoplus_{g\in\Z}
       \operatorname{Hom}_{\widetilde B_L}(M,T_gN),\qquad
 (T_gN)_r=N_{r-3g}.
\]
Only finitely many summands are nonzero. In general $F$ is neither full,
dense, nor Hom-controlled in Zwara's sense.
\end{proposition}
\begin{proof}
A nonzero path is constant-branch with length less than $L$.
Thus the incoming and outgoing total Hom dimensions at a cover vertex are
$2L-1$, and $\dim B_L=3(2L-1)$. Paths and relations lift uniquely; the
deck action is free and has the stated orbit. Push-down sums spaces of equal
residue. A morphism decomposes uniquely into blocks with target lift index
minus source index equal to $-3g$. The arrow commutation equations preserve
that difference, proving the formula; finite supports bound the possible
differences. Direct sums of vector spaces preserve exact sequences.
The following explicit witnesses prove the three failures.
\end{proof}

For $L=5$, let $P$ be the cover interval on $0,1,2,3,4$ with all its
$a$-maps one and $b$-maps zero, and let $S$ be the sum of the five simples
at the same vertices. They have the same $K_0$ class, but
\[
\begin{array}{c|ccc}
 &\dim\operatorname{End}_{\widetilde B_5}
 &\dim\operatorname{End}_{B_5}(F(-))&\text{defect}\\\hline
 S&5&9&4\\ P&1&2&1
\end{array}
\]
Indeed the pushed-down simple multiplicities are $(2,2,1)$, whereas
$\operatorname{End}_{B_5}(FP)=kI\oplus kX^3$ for its single nilpotent
five-step chain $X$. The extra endomorphism $X^3$ proves non-fullness;
the unequal defects at the same class exclude a bilinear $K_0$ defect.
To use Zwara's finite-algebra definition literally, restrict the source to
the finite convex interval $[0,4]$, then extend by zero before push-down.
This exact functor has the same endomorphism spaces and $K_0$ comparison.

\begin{proposition}[An indecomposable module outside this cover image]
\label{cover:band}
Over a characteristic-zero field, $B_3$ has a nine-dimensional
indecomposable representation which is not the push-down of any
finite-support $\widetilde B_3$ representation.
\end{proposition}
\begin{proof}
Label its basis by $0,\ldots,8$ with vertex residues
$(0,1,0,1,2,1,2,0,2)$, and use matrix units $E_{ij}$ acting on columns.
Set
\begin{align*}
 X&=E_{10}+E_{32}+E_{43}+E_{65}+E_{76}+E_{08},\\
 Y&=E_{12}+E_{45}+E_{78}.
\end{align*}
Their restrictions between adjacent vertex spaces give the $a_i,b_i$ maps.
They satisfy $XY=YX=X^3=Y^3=0$, hence every length-three path is zero.
The endomorphism ring preserving vertex spaces is $kI\oplus\mathfrak r$,
where
\[
 \mathfrak r=\langle E_{02}+E_{13},E_{15},E_{35}+E_{46},
                       E_{48},E_{68}+E_{70},E_{72}\rangle,
 \qquad \mathfrak r^3=0.
\]
Solving the arrow commutation equations gives exactly these seven basis
elements. The ring is local: an element $cI+n$ is invertible exactly when
$c\ne0$, by the finite geometric series. It has no nontrivial idempotents,
so the representation is indecomposable.

A finite-support lift would define an Euler operator $D$ acting by $r$ on
the summand at lift vertex $r$. Then $[D,X]=X$ and $[D,Y]=Y$.
Every displayed nonzero arrow is the unique nonzero entry in its row and
column in its own matrix. Its commutator entry is therefore
$D_{ii}-D_{jj}$. Add the six $X$ equations and subtract the three $Y$
equations. Each diagonal term cancels, leaving $0=6-3=3$, a contradiction.
This argument allows arbitrary $D$, not just a diagonal matrix in the
displayed basis.
\end{proof}

\begin{corollary}[Higher-genus and derived non-density]\label{cover:higher}
For the clean dessin with rotations
$\sigma=(0\,\cdots\,m-1)(m\,\cdots\,2m-1)$ and
$\alpha=\prod_{i=0}^{m-1}(i\,i+m)$, odd $m\geq3$, and genus $(m-1)/2$,
the analogous $A_m/J^3$ has a $3m$-dimensional representation
outside its integer-cover image in characteristic zero. Its bounded-derived
push-down is not dense either.
\end{corollary}
\begin{proof}
Use basis $v_0,\ldots,v_{3m-1}$, assign vertices $i,i+1,i$ to
$v_{3i},v_{3i+1},v_{3i+2}$ modulo $m$, and set
\[
 Xv_{3i}=v_{3i+1},\quad Xv_{3i+2}=v_{3i+3},\quad
 Yv_{3i+2}=v_{3i+1},
\]
with indices modulo $3m$ and all other values zero. Again mixed products
and cubes vanish. An Euler operator would yield $2m-m=m=0$ by the same
telescoping calculation, impossible in characteristic zero. If a bounded
finite-support complex pushed to this module in degree zero, exactness would
make its zeroth cohomology a module lift, a contradiction.
\end{proof}

The source's derived result is a \emph{precovering} in the stated bounded
categories \cite[Theorem~4.28]{CoveringLecture2026}; density remains an
additional condition. Neither this result nor the GLS equivariant
classification \cite{TameGLS2026} identifies our completed order with a GLS
algebra. Graded surface/Fukaya models \cite{SchrollGeometric2026} likewise
require their own dg category, stops and grading. These categorical tools
do not construct the independent positive arithmetic determinant.


\section{Verification and research frontier}
The two nodal/invariant Lean files construct the node as an actual subring of $k[[X]]\times k[[X]]$ with equal constants, prove its annihilators, regular common parameter, augmentation, alternating projective resolution and a bijective pair of power-series coordinates. A second file proves CK1 collapse under the stated permutation/coverage hypotheses, determinant base change, actual $2\times2$ symmetric-square functoriality and relation preservation, a Schur character identity and the fixed-parameter lattice failure. These thirty printed theorem types have only standard Lean foundations. They do not formalize general dessins, abstract Ext/AR/Schur--Weyl, C* exactness, Fredholm determinants, arithmetic descent or GRH. Earlier formal layers remain preserved with their literal scopes.

Exact-rational diagnostics independently check the displayed dessins, path identities, Schur and Young examples and the nine-object grid. These finite checks supplement the proofs, rather than prove infinite spectral or arithmetic assertions. The next arithmetic target is an exact descent/Frobenius-compatible construction on the labelled order/Fock data, with a determinant-preserving positive comparison. The local Schur kernel and the simple degree spectrum alone do not supply it. All source clones and earlier research are retained; the earlier general focused PDFs are archived at the preceding commit.

Ten further declarations in \texttt{DessinPullbackComponents.lean} check actual matrix-unit identities, nonmultiplicativity of raw diagonal extraction, impossibility of a unital two-state matrix residue into a nonzero commutative ring, multiplicativity on diagonal fibers, and genuine equal-image pullback subring/projection components. The analytic suite then contained187printed declarations at unchanged pins. The circle-function C* model, functional calculus, generation, graph uniqueness and K-theory remain prose/external-theorem arguments, not hidden assumptions of these component checks.

Sixteen additional declarations in \texttt{ThetaJacobi.lean} check genuine finite-field MulChar Jacobi conjugation, norm and rescaling with explicit nontriviality hypotheses; cubic--quadratic noncancellation; the cyclotomic multiplication determinant and positive Gram identity; and actual theta cellular boundary, homology-inclusion, restriction and rational chain-map matrices. The full analytic suite now contains203declarations. The point-count proof, deck/Kummer comparison, cohomological trace interpretation and global arithmetic comparison are not thereby fully formalized.

The final three \texttt{ThetaJacobi.lean} declarations verify the actual conjugation matrix identities and the real two-dimensional inert-prime determinant obstruction. They do not assume any unknown zero location or exclude complex/enlarged spectral realizations.


\texttt{ThetaProjectiveMultiplicity.lean} adds five declarations. Arbitrary upper-block matrices with six identical diagonal blocks have the sixth-power determinant and characteristic-factor formulas. A literal power-series constant-coefficient map specializes matrix determinants, for each scalar parameter; the coefficientwise polynomial/order interface remains prose. A nonzero quadratic cannot be a rational-function sixth power, proved by polynomial degrees. Three genuine power series with product $X$ cannot have the cyclically equivalent residue-vanishing conditions forced by invertible rank-one arrow intertwiners. The derivation of the path blocks from the full theta order, projective decomposition, derived branch functor and deck intertwiners is not fully formalized. The complete analytic suite now has208printed declarations; no hidden arithmetic positivity assumption occurs.


\texttt{ThetaTameDescent.lean} adds twelve printed component declarations (and one auxiliary polynomial lemma). They use genuine formal-power-series rescaling and its coefficient operator, proving the coefficient action, invariant-coefficient vanishing and the tame operator's cube. Actual three-state matrices satisfy the clock/rotation commutator, comparison determinant and closed-fiber equations, weighted-cycle cube and full characteristic polynomial, noncommutation of a nonconstant rotation polynomial, trivial whole-column determinant actions, and the same-operator comparison for $vC$. The classification of all fixed series by residues, full order action and hereditary invariant algebra, stack interpretation, valuation/rank theorem, general coefficient-functor filtration and GRH are not thereby fully formalized. No unknown positivity premise is assumed.

\section{Why an absolute coprime pair-square estimate cannot close the short rows}
\label{sec:pair-barrier}
Keep the original sextic family
\[
 M_u(D)=D^{-1/2}\sum_n\mu(n)\nu(n)W(Nn/D)\chi_n(u),
 \qquad \chi_n(u)=(u/n)_6,
\]
with squarefree primary column ideals outside the fixed excluded set,
their full annular coefficients (fix $W$ supported in $[1,2]$), and all
nonzero Eisenstein element rows
$Nu\le H$. The cubic symbol is $\chi_n^2$, not $\chi_n$.
The following classical frame inequality \cite{WelchGeometry} gives a new
obstruction to one proposed estimate in this arithmetic program.

\begin{proposition}[Coprime pair-square obstruction]\label{prop:pair-barrier}
Let $A$ be an $R$-by-$N$ complex matrix with $|A_{ui}|=1$, and put
$K_{ij}=\sum_u A_{ui}\overline{A_{uj}}$. Then
\[
 \sum_{i\ne j}|K_{ij}|^2\ge NR(N-R).
\]
For the original character kernels, let $B_1(H,D)$ be the sum of
$|K_H(n,m)|^2$ over all distinct coprime squarefree column ideals in
$[D,2D]$. At $H=D^{3/4}$,
\[
 B_1(H,D)\gg \frac{D^{11/4}}{(\log D)^2}.
\]
In particular the proposed sufficient input
$B_1(H,D)\ll H^2D^\varepsilon$ fails for each fixed $\varepsilon<5/4$.
This does not refute the signed M\"obius mean square.
\end{proposition}
\begin{proof}
Set $F_{uv}=\sum_i A_{ui}\overline{A_{vi}}$. Reordering four finite sums
proves $\sum_{i,j}|K_{ij}|^2=\sum_{u,v}|F_{uv}|^2$.
The row diagonal contributes $RN^2$; the column diagonal contributes
$NR^2$. Removing the latter and retaining the former gives the inequality.

Restrict the column pairs in the nonnegative sum $B_1$ to distinct prime
ideals $p,q$ of norm in $[D,2D]$, outside the same fixed set. If $H<D$,
no such prime divides a nonzero row element: divisibility would imply
$Nu\ge Np$. Thus all original zero masks are exactly one on this ball,
$|\chi_p(u)|=1$, and the residual primitive finite character conductor is
$pq$, of norm between $D^2$ and $4D^2$. All these pairs have common divisor1.
The lattice count gives $R=(2\pi/\sqrt3)H+O(\sqrt H+1)$.
The classical unconditional fixed-field prime-ideal theorem gives
$N\sim D/\log D$; this is an external analytic input, not a GRH hypothesis.
At the stated endpoint $R=o(N)$, and the bound is therefore
$(1-o(1))N^2R\gg D^{11/4}/(\log D)^2$.
\end{proof}

To see the exact failed step, write $c_n=\mu(n)\nu(n)W(Nn/D)$ and
$C_D=\sum_n|c_n|^2=O(D)$. Cauchy over coprime pairs bounds the absolute
signed contribution by $D^{-1}C_DB_1^{1/2}$. The hoped-for pair-square
upper bound would have closed this particular route, but is false.
A large absolute pair-square mass is not a lower bound on the quadratic
form with the fixed $c_n$. Nor is the prime part a positive summand of the
full M\"obius energy: prime--composite cross terms can cancel. The original
low-gcd estimate and the positive arithmetic determinant bridge remain open.

\section{A signed M\"obius decomposition and its termwise obstruction}
\label{sec:mobius-split}
Keep the original squarefree good-ideal family, its fixed $\nu$, its
full $W(Nn/D)$ and every zero-extended sextic symbol. In ideal Dirichlet
convolution, write $\zeta_0(n)=1$ for every good ideal and let $\delta$
be the convolution unit. For $1\le U<D$ put
$a_U=\mu\mathbf1_{Nn\le U}$ and $b_U=\mu-a_U$. The classical Vaughan
algebra gives
\begin{equation}\label{eq:mobius-split}
 \mu=2a_U-a_U*a_U*\zeta_0+b_U*b_U*\zeta_0.
\end{equation}
Indeed $\mu*\zeta_0=\delta$, and expansion of the last product proves the
identity without an analytic assumption. This is a classical decomposition,
not a new cancellation theorem; compare \cite[Lemma4.1]{GreenTaoVaughan}.

On squarefree $n$ all factors in $n=d c=a b c$ are pairwise coprime.
Define the actual, nonarbitrary coefficient functions
\[
 \beta_U(d)=\sum_{ab=d\atop Na,Nb\le U}\mu(a)\mu(b),\qquad
 g_U(s)=\sum_{b\mid s\atop Nb>U}\mu(b).
\]
For squarefree $d$, $\beta_U(d)$ is $\mu(d)$ times the number of the
displayed ordered factorizations. For $Ns>U$,
$g_U(s)=-\sum_{b\mid s,\,Nb\le U}\mu(b)$. The two regrouped pieces are
\begin{align*}
 I_u&=D^{-1/2}\sum_{d,c\ {\rm squarefree}\atop(d,c)=1}
  \beta_U(d)\nu(dc)\chi_d(u)\chi_c(u)W(NdNc/D),\\
 II_u&=D^{-1/2}\sum_{a,s\ {\rm squarefree}\atop Na,Ns>U,\,(a,s)=1}
  \mu(a)g_U(s)\nu(as)\chi_a(u)\chi_s(u)W(NaNs/D).
\end{align*}
Then $M_u=-I_u+II_u$. All ideals here remain outside the same fixed
excluded set. The short $2a_U$ term vanishes because $U<D$ and $W$ is
supported in $[1,2]$. No phase or coprimality mask was removed.

\begin{proposition}[An all-cutoff separate-piece bound is false]
\label{prop:mobius-unit-cutoff}
Take the allowed principal $\nu$, $U=1$, and a fixed nonnegative smooth
$W$ supported in $[1,2]$, bounded below by a positive constant on a closed
subinterval of $(1,2)$ of positive length.
For each of the two pieces,
\[
 \sum_{0<Nu\le H}|I_u|^2\gg_W\frac{D}{\log^2D},\qquad
 \sum_{0<Nu\le H}|II_u|^2\gg_W\frac{D}{\log^2D}
 \quad(H\ge1).
\]
Consequently a proposed bound $H D^\varepsilon$ for each piece at
$H=D^{3/4}$ fails for $\varepsilon<1/4$. In the exact original notation
$H(HD)^\varepsilon$, it fails for $\varepsilon<1/7$.
This refutes an all-cutoff termwise strengthening, not the original signed
mean square or a specified moving-cutoff argument.
\end{proposition}
\begin{proof}
At $U=1$, $a_U=\delta$, so $a_U*a_U*\zeta_0=\zeta_0$ and
$b_U*b_U*\zeta_0=\mu+\zeta_0-2\delta$.
On the annulus the two coefficients are $1$ and $1+\mu(n)$, respectively.
On squarefree ideals the latter is zero or two. At the actual unit row
$u=1$, every original character equals one. A fixed-field prime-ideal
theorem gives $\gg D/\log D$ good prime columns in the subinterval,
so $I_1\gg_W\sqrt D/\log D$. Fix a good prime ideal $q_0$.
Products $q_0p$ with varying distinct good primes in the rescaled
subinterval have $\mu(q_0p)=1$. Their nonnegative contributions give
$II_1\gg_W\sqrt D/\log D$. Keeping that one row in each nonnegative
energy proves the result. The prime-ideal theorem is a classical
unconditional external input, with the field and $q_0$ fixed.
\end{proof}

The factorization does not shorten the primitive row conductor:
for squarefree coprime $a,s$ it is $as$, not $a$ or $s$ separately.
For two resulting columns $n,n'$ with common factor $h$, the oscillatory
conductor is $nn'/h^2$ and a zero-preserving period is $nn'/h$.
The signed low-gcd estimate and the independent positive arithmetic
determinant comparison therefore remain open.

\section{A full-mask Hecke expansion for the type-I piece}
\label{hti:core}
We now continue the signed decomposition through its actual squarefree
inner sum. Fix $K=\mathbb Q(\sqrt{-3})$, the finite-order Hecke character
$\nu$, and the same finite excluded prime set $S$, containing primes above
$6$ and the conductor of $\nu$. For $u\in\mathcal O_K\setminus\{0\}$,
let $\psi_u$ be the primitive Hecke character induced by
$\nu(n)(u/n)_6$ away from $S\cup\operatorname{supp}(u)$.
It is defined before, and independently of, the type-I divisor $d$.
The abelian Kummer extension $K(u^{1/6})/K$ constructs the sextic
factor; multiplying by $\nu$ gives the stated character.
This uses classical abelian reciprocity, not general Artin holomorphy.
Write $\psi_{u,2}$ for the primitive character inducing its square and
$\mathfrak f_u,\mathfrak f_{u,2}$ for the two primitive conductors.
For a squarefree good ideal $d$ put
\[
 T=T(u,d)=S\cup\operatorname{supp}(u)\cup\operatorname{supp}(d),\qquad
 R_T=\prod_{p\in T}Np,\quad Q_u=|D_K|N\mathfrak f_u.
\]
All primes in $T$ are deleted even if absent from either primitive conductor.

\begin{theorem}[Restoration and a uniform scalar contour bound]
\label{hti:theorem}
For $\Re s>1$ the original type-I inner Dirichlet series is
\begin{align}
 F_{u,d}(s)
 &=\sum_{c\ {\rm squarefree}\atop(c,T)=1}
       \frac{\psi_u(c)}{(Nc)^s}
   =\frac{L(s,\psi_u)}{L(2s,\psi_{u,2})}\,C_T(s),\label{hti:ratio}\\
 C_T(s)
 &=\frac{\displaystyle\prod_{p\in T,\ p\nmid\mathfrak f_u}
             (1-\psi_u(p)(Np)^{-s})}
          {\displaystyle\prod_{p\in T,\ p\nmid\mathfrak f_{u,2}}
             (1-\psi_{u,2}(p)(Np)^{-2s})}.
 \label{hti:restore}
\end{align}
It continues meromorphically to $\Re s>1/2$, with only a possible
simple pole at $s=1$ when $\psi_u=1$. In that case its residue is
\[
 \rho_T=\frac{\operatorname*{res}_{s=1}\zeta_K(s)}{\zeta_K(2)}
             \prod_{p\in T}(1+(Np)^{-1})^{-1}.
\]
Fix $0<\delta<1/2$, $\varepsilon>0$ and
$W\in C_c^\infty((0,\infty))$ supported in $[1,2]$. Set
\[\mathcal M_4(W)=\sum_{j=0}^4\sup_x|(x\partial_x)^jW(x)|.\]
Uniformly for $Y\ge1$, $u$ and $d$,
\begin{align}
 \sum_{c\ {\rm squarefree}\atop(c,T)=1}\psi_u(c)W(Nc/Y)
 ={}&\mathbf1_{\psi_u=1}\rho_TY\widehat W(1)\notag\\
 &+O_{\nu,S,\delta,\varepsilon}\!\left(
 \mathcal M_4(W)Y^{1/2+\delta}
 Q_u^{1/4-\delta/2+\varepsilon}R_T^\varepsilon\right).
 \label{hti:scalar}
\end{align}
Here $\widehat W(s)=\int_0^\infty W(x)x^{s-1}\,dx$.
The continuation and functional equation of degree-one Hecke functions
are classical external inputs; the contour argument is given below.
\end{theorem}

At a good prime with $v_p(u)=3$, $\psi_u$ is ramified while its square
is unramified; the second product in~\eqref{hti:restore} must restore a
factor invisible to squaring the zero-extended primitive value.
At $v_p(u)=6$ both characters can forget the prime, but the original row
still vanishes on columns divisible by it. In particular, for principal
$\nu$, $u=\pi^3$ and $u=\pi^6$ at $N\pi=7$, the deleted local factors are
\[
\begin{array}{c|ccc}
 &\text{primitive ratio at }s=1&\text{restoration}&\text{full masked factor}\\\hline
v_p(u)=3&48/49&49/48&1\\
v_p(u)=6&8/7&7/8&1
\end{array}
\]
These are exact local calculations, not numerical evidence about zeros.

For $1\le U$, $U^2\le D$, insert $Y=D/Nd$ into~\eqref{hti:scalar}
and retain the actual $\beta_U(d)$ and outer $\nu(d)\chi_d(u)$.
Classical tame Kummer conductors give $Q_u\ll_\nu Nu$. Thus the normalized
$I_u$ equals its displayed principal-residue sum plus an error $E_u$ with
\begin{equation}\label{hti:outer}
 |E_u|\ll_{\nu,S,\delta,\varepsilon}\mathcal M_4(W)
 D^\delta U^{1-2\delta+\varepsilon}
 (Nu)^{1/4-\delta/2+\varepsilon}.
\end{equation}
The resulting elementary energy envelope is only
\begin{equation}\label{hti:energy}
 \sum_{0<Nu\le H}|E_u|^2
 \ll\mathcal M_4(W)^2D^{2\delta}U^{2-4\delta+\varepsilon}
 H^{3/2-\delta+\varepsilon}.
\end{equation}
It does not prove the short-row target of order $H$ up to small losses.
This is an insufficient upper bound, not a lower bound on the actual
error or a counterexample to coefficient-sensitive averaging. The
original signed mean square and positive determinant bridge stay open.

\section{Actual renormalization retains the effective action}
\label{sec:actual-rg}
A compression score does not specify an RG transformation. Following the
real-space formulation of Wilson and Kogut \cite{WilsonKogutRG}, we specify
the eliminated variables, effective action, normalization and scale map.
Our exact finite decimation is distinct from an approximate tensor-network
truncation such as Levin--Nave TRG \cite{LevinNaveRG}.

\begin{definition}[Finite classical RG step]
Let the fine configuration be $(s,t)\in S\times E $, with finite nonempty
sets $ S,E $, and dimensionless log Boltzmann weight $ V(s,t)=-\beta H(s,t)$.
Define
\[
 (\mathcal RV)(s)=\log\sum_{t\in E}e^{V(s,t)},\qquad
 Z(V)=\sum_{s,t}e^{V(s,t)}.
\]
The full action includes its additive constant. A truncation or a learned
approximation is a further operation with its own error, not this identity.
\end{definition}
\begin{proposition}[Partition, observables and composition]\label{rg:exact}
One has $ Z(\mathcal RV)=Z(V)$, and the normalized expectation of every
observable depending only on $ s $ agrees in the two descriptions. Eliminating
$ E_1$ and then $ E_2$ equals eliminating $ E_1\times E_2$ at once, with the
same tensor labels and inverse temperature.
\end{proposition}
\begin{proof}
Exponentiating the definition gives $ e^{\mathcal RV(s)}=\sum_te^{V(s,t)}$.
Sum over $ s $, optionally multiplying by an observable $ O(s)$, and divide
both sides by the same strictly positive $ Z $. Reordering finite sums proves
the composition statement. No projection onto pair interactions is used.
\end{proof}

For an even homogeneous periodic Ising chain, write the transfer matrix as
$ T_{ab}=\exp(C+Kab+h(a+b)/2)$, $ a,b\in\{-1,1\}$.
Its partition function is $\operatorname{tr}T^N $. Summing alternate spins
replaces $ T $ by $ T^2$ and $ N $ by $ N/2$; the physical spacing doubles.
Writing $ P=T^2$, its positive entries determine
\[
 K'=\tfrac 14\log\frac{P_{++}P_{--}}{P_{+-}^2},\quad
 h'=\tfrac 12\log\frac{P_{++}}{P_{--}},\quad
 C'=\tfrac 14\log(P_{++}P_{--}P_{+-}^2).
\]
These three parameters reproduce every entry of $ P $, so no normalization
is lost. At $ h=0$, $\tanh K'=\tanh^2 K $ and the eigenvalue-ratio correlation length of the homogeneous extension
in blocked units halves. Its product with the spacing stays fixed. The
fixed points in $ t=\tanh K $ are $ t=0,1$; the latter is zero temperature,
not a positive-temperature transition. Odd dessin examples below use the
full-action definition, not this even-chain shortcut.

\subsection{Genuine dessin incidence models}
For odd $ m $, the clean dessin $ y^2=x^m-1$, $\beta=x^m $ has two black
vertices, $ m $ white vertices and one face; its genus is $(m-1)/2$.
Its incidence graph is $ K_{2,m}$ with $2 m $ clean edges. Its infinite
locally gentle order quiver instead has $ m $ vertices and completed
DVR-rank $2 m^2$. We place physical spins $ B_+,B_-,s_0,\ldots,s_{m-1}$
on the incidence vertices; this is a specified finite spin model, not a
faithful representation of that infinite algebra or order. The abstract
incidence graph alone does not recover the ribbon embedding or genus.

With $\eta_j=(-1)^j $ and fixed nonuniform fields $ h_j $, choose
\[
 V=J_+B_+\sum_js_j+J_-B_-\sum_j\eta_js_j
       +bB_+-\tfrac 35 bB_-+\sum_jh_js_j.
\]
Summing the two black spins gives the \emph{entire} effective action
\[
 V'(s)=\sum_jh_js_j+
 \log\!\left[2\cosh\!\left(b+J_+\sum_js_j\right)\right]+
 \log\!\left[2\cosh\!\left(-\tfrac 35 b+J_-\sum_j\eta_js_j\right)\right].
\]
Every one of the $2 m $ physical edges is present, although its couplings
belong to a tied three-parameter family. The Walsh expansion
$ V'(s)=\sum_{A\subseteq\{0,\ldots,m-1\}}c_A\prod_{j\in A}s_j $,
with $ c_A=2^{-m}\sum_sV'(s)\prod_{j\in A}s_j $, is unique by Boolean
orthogonality. It retains the constant $ c_\varnothing $ and genuinely
generated higher-body interactions. Our genus-two and genus-four examples
use $ m=5,9$ and continue by exact nested marginalization. Dropping all terms
of degree greater than two is displayed as a measured approximation.

\subsection{Quantum thermal effective Hamiltonians}
For a finite Hermitian Hamiltonian $ H $ on nonzero finite-dimensional
$\mathcal H_S\otimes\mathcal H_E $
and $\beta>0$, define
\[
 W'=\operatorname{Tr}_E e^{-\beta H},\qquad
 H_{\mathrm{eff}}(\beta)=-\beta^{-1}\log W'.
\]
The matrix $ W'$ is strictly positive: for $ v\ne 0$, sum the strictly
positive quadratic forms of $ v\otimes e_j $ over an orthonormal basis of
$\mathcal H_E $. Its Hermitian spectral logarithm is therefore defined.
The identity $ e^{-\beta H_{\mathrm{eff}}}=W'$ proves partition and retained
\emph{equal-time} observable preservation by the partial-trace property.
Repeated traces compose. The full Pauli expansion, including its identity
coefficient, is retained. It is generally nonlocal and depends on $\beta $;
this statement does not preserve all fine real-time dynamics or supply a
local zero-temperature/thermodynamic RG fixed point.

\texttt{FiniteRG.lean} checks nine actual finite scalar marginal,
exponential-action, partition, observable, composition, error-sandwich and
transfer-trace statements. The quantum spectral-log and complete
thermodynamic/field-theory interfaces are not formalized by these lemmas.
These physics constructions give no arithmetic determinant comparison or
new GRH conclusion.

\section{Linear relations and functorial branch filtrations}\label{rel:core}
The original linear-relations chapter mixes graphs, partial relations and
stable kernels. We keep these different objects separate. For a field $k$,
a linear relation is a subspace $C\subseteq V\oplus V$; its inverse swaps
coordinates and $DC=\{(x,z):\exists y,(x,y)\in C,(y,z)\in D\}$.
The full graph $\Gamma_T=\{(x,Tx):x\in V\}$ always has dimension
$\dim V$, including pairs with zero output. In particular
$\Gamma_T^n(0)=0$, whereas $(\Gamma_T^{-1})^n(0)=\ker T^n$.
Thus the source's $\kappa(\Gamma_T)=\bigcup_n\Gamma_T^n(0)$ is zero,
not the ordinary stable kernel.

For odd $m\geq3$, use the actual clean dessin of
$y^2=x^m-1$, $\beta=x^m$, on the compact smooth curve, in
characteristic zero. Its genus is $(m-1)/2$; its infinite locally gentle
algebra $A_m$ has vertices $i\in\mathbb Z/m$, arrows
$a_i,b_i:i\to i-1$, and mixed consecutive branch products zero.
Below paths act in their direction of travel on left representations over
any coefficient field $k$. Let $\mathcal E_m$ be the full category of
finite-dimensional representations with every $b_i=0$ and $a_i$ invertible
for $i\ne0$.

\begin{theorem}[A concrete relation category on an infinite dessin algebra]
\label{rel:branch}
The category $\mathcal E_m$ is equivalent to finite-dimensional pairs
$(V,T)$ with intertwining maps. A normal form $N_m(T)$ puts $V$ at each
vertex, $a_0=T$, and all other $a_i=1$. For $n=\dim V$,
\[
 V=K(T)\oplus I(T),\qquad K(T)=\ker T^n,\quad I(T)=\operatorname{im}T^n.
\]
These summands and their kernel/image filtrations are functorial. The first
restriction is nilpotent and the second invertible. The central full-turn
parameter $u$ acts as $T$. A finite $N_m(T)$ extends continuously to the
$u$-adically completed order precisely when $T$ is nilpotent. Such modules
are torsion, not nonzero order lattices.
\end{theorem}
\begin{proof}
Transport all vertex spaces to vertex zero along the invertible arrows;
the remaining map is the full-turn monodromy $T$. Arrow-compatible maps
become exactly $hT=T'h$. Finite kernel/image chains stabilize by step $n$.
If $v=T^nw$ and $T^nv=0$, then $w\in\ker T^{2n}=\ker T^n$, hence $v=0$;
rank-nullity proves the direct sum. Stability gives the two restrictions,
and $hT^r=T'^rh$ proves functoriality. Continuity on a finite-dimensional
discrete space means some $u^r$ kills it. Conversely nilpotence makes every
formal series act by a finite sum. This proves the completion assertion.
\end{proof}

Over an algebraically closed $k$, the indecomposable branch pairs are
$J_r(\lambda)$. The value $\lambda=1$ is permitted for a polynomial band;
it is not a forbidden deformation parameter. The nilpotent blocks are
strings, but are not all rigid: $N_m(J_2(0))$ is a nonsplit self-extension
of $N_m(J_1(0))$. The Hom dimension between two such Jordan pairs is
$\min(r,s)$ for equal eigenvalues and zero otherwise. These assertions
concern this specified full subcategory, not all infinite modules.

If $L$ transports once along the $a$-arrows on $\bigoplus_iV$, then
\begin{equation}\label{rel:det}
 \det(1-zL)=\det(1-z^mT).
\end{equation}
Eliminating the identity arrow blocks leaves the final monodromy block;
the eliminated identity blocks have determinant one. For $J_r(\lambda)$
the answer is $(1-\lambda z^m)^r$. Nilpotent strings therefore contribute
one, and arbitrary bands have no forced unit-circle zero location. This
determinant is an actual cyclic-quiver invariant; it is not the missing
arithmetic positive determinant comparison.

\section{A kernel-checked extension over the infinite dessin algebra}
\label{bre:core}
The preceding branch example can be realized on every element of the
infinite algebra, rather than verified only on sampled arrow matrices.
Use the complex genus-two clean dessin $y^2=x^5-1$, $\beta=x^5$, and any
coefficient field $k$. Put $H\subseteq M_5(k[t])$ equal to the matrices
whose reduction modulo $t$ is upper triangular. Let
\[
 B=H\mathbin{\times}_{k^5}H,
\]
where each map to $k^5$ takes the constant diagonal. Equivalently, $B$ is
generated inside $M_5(k[t])^2$ by
\[
 e_i=(E_{ii},E_{ii}),\quad
 a_i=(w_iE_{i-1,i},0),\quad b_i=(0,w_iE_{i-1,i}),
 \qquad w_0=t,\quad w_i=1\ (i\ne0).
\]
Indices are modulo five. Matrix multiplication acts on left column
modules. The path-normal-form proof below identifies $B$ with $A_5$;
this is a polynomial algebra, not an arrow-truncated finite quotient.

\begin{theorem}[An ambient module extension]\label{bre:extension}
Let $D=k[\epsilon]/(\epsilon^2)$, $M_1=k^5$ and $M_2=D^5$.
For $(C(t),D'(t))\in B$, act on $M_1$ by $C(0)$ and on $M_2$ by
$C(\epsilon)$. These are unital $B$-modules, and
\[
 0\longrightarrow M_1\xrightarrow{\ i\ }M_2
 \xrightarrow{\ p\ }M_1\longrightarrow0,
 \qquad i(x)=\epsilon x,\quad p(x+\epsilon y)=x,
\]
is exact and nonsplit. Moreover $B$ is infinite dimensional over $k$.
The assertions hold in every characteristic, including two.
\end{theorem}
\begin{proof}
Evaluation of polynomials gives both actions. Multiplying
$C(\epsilon)=C_0+\epsilon C_1$ proves $B$-linearity of $i,p$;
their kernel and image follow coordinatewise. If $s$ were a $B$-linear
section, the constant coefficient of $(s e_0)_0$ would be one. But $a_0$ kills
$M_1$, while $(a_0s e_0)_4=\epsilon(s e_0)_0$ has nonzero
$\epsilon$ coefficient, a contradiction. Finally the permitted turn
$c=a_1a_2a_3a_4a_0$ has first component $tE_{00}$.
The unital algebra map $k[t]\to B$, $f\mapsto f(c)$, is injective:
its $(0,0)$ first-component entry is $f(t)$. Thus $B$ cannot be
finite dimensional.
\end{proof}

In the preceding notation these modules are $N_5(J_1(0))$ and
$N_5(J_2(0))$. They are finite torsion modules over an infinite algebra;
they are not nonzero lattices over its completed surface order. The
nonsplit sequence defines a nonzero Yoneda extension class in the ambient
module category. No assertion that it is almost split in that entire
category is needed or inferred. In particular it does not supply an
arithmetic spectral determinant or a zero-location theorem.


\begin{theorem}[Unbounded normal-image faithfulness]\label{bre:normal-faithful}
In $M_5(k[t])^2$, index the stationary images by $i$ and the positive
images by $(b,i,j,q)$, where $b\in\{a,b\}$, $0\leq i,j<5$, $q\geq0$
and $i<j$ or $q>0$. The positive image is $t^qE_{ij}$ in factor $b$,
with zero in the other factor. Its path length is $5q+j-i>0$.
All these images, including the shared stationary images, are linearly
independent over every field. Every finite normal expansion has uniquely
recoverable coefficients, without a bound on degree or path length.
\end{theorem}
\begin{proof}
Read a stationary coefficient from the first factor's constant diagonal;
read a positive coefficient from the indicated factor, entry and degree.
The index restriction ensures that no positive image contributes to a
constant diagonal. These linear selectors are biorthogonal to the images,
so applying each selector to a finite expansion recovers its coefficient.
\end{proof}
This theorem is now kernel checked for actual polynomial matrix pairs.
The membership and spanning steps are supplied by the next theorem.
Comparison with a separately constructed path quotient remains a distinct
algebraic formalization task.


\begin{theorem}[Normal coordinates in the generated algebra]\label{bre:normal-generated}
Let $A\subseteq M_5(k[t])^2$ be the algebra generated by the displayed
vertices and arrows, and let $\mathcal K$ be the normal index set of
Theorem~\ref{bre:normal-faithful}. Every normal image belongs to $A$.
The actual $A$-valued finite expansion
\[
 E:k^{(\mathcal K)}\longrightarrow A,\qquad
 c\longmapsto\sum_{x\in\mathcal K}c_xv_x
\]
is a linear isomorphism, whose inverse is $(E^{-1}a)_x=\lambda_x(a)$.
Thus the normal images form a basis of the infinite generated algebra,
not merely an independent family in its ambient matrix space.
\end{theorem}
\begin{proof}
Write $P_b(i,j)$ for the shortest positive branch walk from $j$ to $i$,
using a full turn when $i=j$. Its image is $t^dE_{ij}$ in branch $b$,
where $d=0$ if $i<j$ and $d=1$ otherwise. The actual arrow product
$P_b(i,i)^nP_b(i,j)$ has image $t^{n+d}E_{ij}$ for every $n\geq0$.
Admissibility gives $q\geq d$, so $n=q-d$ proves membership.

The normal span contains the generators and $1=\sum e_i$ and is
multiplicatively closed: stationary factors select endpoints; positive
factors either multiply to zero or to $t^{q+r}E_{il}$ in the same branch.
When $q+r=0$, the two strict upper inequalities imply $i<l$, so this
output slot is still permitted. Adjoin induction gives spanning of $A$.
The previous coefficient selectors prove uniqueness and the inverse formula.
\end{proof}
All membership, spanning and finite-expansion assertions in this theorem
are now Lean checked over arbitrary fields. The formal linear equivalence
is not asserted to be an algebra map from an independently defined path
quotient, nor a statement about completed orders, derived categories or GRH.

\section{Learned fields, width limits and quantum geometry}\label{lf:core}
Learned functions enter this program as candidate generators, physical
surrogates and control laws. A network may suggest an arithmetic metric,
operator or gluing datum; a successful fit does not establish its positivity,
analytic domain or identification with a completed $L$-function. Those are
independent mathematical obligations. This section gives explicit interfaces
and executed examples rather than treating ``quantum'', ``tropical'' or
``infinite width'' as a substitute for that comparison.

\begin{proposition}[A specified quantum-feature width limit]\label{lf:kernel}
Let $x\in\mathbb R^d$, $w\sim N(0,I_d)$ and
$b\sim\mathrm{Unif}[0,2\pi]$ be independent, and set
$\phi_{w,b}(x)=\sqrt2\langle Z\rangle_{R_Y(w\cdot x+b)|0\rangle}
=\sqrt2\cos(w\cdot x+b)$. For independent features,
\[
 K_m(x,y)=m^{-1}\sum_{r=1}^m\phi_r(x)\phi_r(y)
 \longrightarrow K(x,y)=e^{-\|x-y\|^2/2}
\]
almost surely on every specified finite set of pairs. Every $K_m$ is a
positive semidefinite Gram kernel. For a readout
$f_a(x)=m^{-1/2}\sum_r a_r\phi_r(x)$ with only $a$ trained, its tangent
kernel is exactly $K_m$. These statements do not assert the same limit
for a jointly trained finite quantum transformer.
\end{proposition}
\begin{proof}
Uniform-phase integration gives
$\mathbb E_b[2\cos(u+b)\cos(v+b)]=\cos(u-v)$.
The Gaussian characteristic function then gives $K$. Products are bounded
by two, so the strong law applies to each pair and hence any finite union.
For any finite coefficient vector $c$, the Gram quadratic form is
$m^{-1}\sum_r|\sum_i c_i\phi_r(x_i)|^2\geq0$.
Differentiating with respect to $a_r$ gives the final assertion.
\end{proof}
The frozen-kernel squared-loss flow has residual
$r(t)=\exp(-K_m t)r(0)$ when the loss normalization is absorbed into time.
Null directions do not decay. A complex feature space instead uses the
Hermitian kernel $\mathbb E[\phi(x)\overline{\phi(y)}]$.
A finite coefficient field has neither this Gaussian law nor an ordered
positive Hilbert metric; changing algebraic coefficients is a separate operation.

\begin{proposition}[Tropical limit and a cancellation boundary]\label{lf:tropical}
For real $s_1,\ldots,s_n$ and $\tau>0$,
\[
 0\leq\tau\log\sum_j e^{s_j/\tau}-\max_j s_j\leq\tau\log n.
\]
An attention maximum therefore has a controlled log-sum-exp approximation.
It does not by itself define a tropicalization of a variety or an HMS functor.
In a valued coefficient field, valuation of a sum can increase by cancellation.
\end{proposition}
\begin{proof}
After subtracting $M=\max s_j$, the exponential sum lies between one and $n$.
For the second statement, both $1+T$ and $-1+T$ have valuation zero in
$\mathbb Q((T))$, whereas their sum $2T$ has valuation one. In characteristic
two that sum is zero. Thus replacing every addition by its tropical minimum
without tracking leading coefficients loses information.
\end{proof}
Kontsevich--Soibelman constructions require valued analytic and
integral-affine gluing data; appropriate gentle/Fukaya comparisons require
specified categories, gradings, stops and smoothness hypotheses
\cite{LFAffine,LFQNTK,LFFukaya}. Learned coefficients can propose such data,
but the compatibility equations and arithmetic comparison still need proofs.

\paragraph{Executed physical examples.}
On the actual clean genus-two dessin, put seven quantum spins on its
$K_{2,5}$ incidence graph and compute a normalized Gibbs state. Exact white-spin
partial trace gives a two-spin quantum RG marginal. Twelve RTX4090 runs train
four-qubit-feature diffusion and flow-matching models for its smoothed
$Z$-measurement law: two objectives, entangled/product circuit controls, and
three seeds. They have 5,810 nominal parameters and 2,000 steps each.
A separate four-qubit variational purification flow learns the full known
density interpolation, with 48 controls. Its minimum fidelity on a 33-point
grid is $0.999903$; this is not a uniform theorem or universal noise inverse.
The matched controls show no consistent entanglement advantage.

Actual cuStateVec/cuTensorNet execution also covers 7-, 11- and 23-qubit
finite circuits on genus-two, genus-four and genus-ten dessin incidence
graphs. The 23-qubit case has 8,388,608 amplitudes and 447 gates. These larger
circuits are untrained; the small cases have independent NumPy comparisons,
while the largest has normalization and inverse-circuit checks.
The exact environment and all failed/preliminary records are preserved.

\paragraph{Geometric and hardware interfaces.}
For the contracted $m$-edge dessin, the graph fundamental group is
$F_{m-1}$, and its universal graph cover is the $m$-regular tree.
The theta case $m=3$ has a three-regular cover, distinct from the four-regular
Cayley graph of $F_2$. Neither is the coefficient field $\mathbb F_2$ or the
universal cover of the filled compact surface. Explicit path modules over the
infinite genus-two algebra give Hom, Ext and Tor computations; Frobenius
pants tensors give exact TQFT gluing, with a heralded physical implementation.
CHSH gives a finite game benchmark, distinct from the unrestricted
correlation separation in MIP*$=$RE \cite{LFMIP}. Three executable tropical
hardware references and a source-pinned chip-design survey complement the
circuits. They are arithmetic/control designs, not synthesized chips,
validated CFT/string dualities or a proof of GRH.

\section{Complementary physical datasets and checked dynamics}
The learned program now has two additional, separately constructed corpora:
16,384 effective-material response records and22,951 molecular conformer/force
records in their default configurations. Time samples are not counted as
independent systems. Source-linked and dynamical configurations are additional.
The active HolonForge/Aster and Mora-3 repositories were read only; no data,
vocabulary, checkpoint, source enrollment or training process was changed.

The material corpus covers SSH, Kitaev, QWZ, Rice--Mele, Harper and actual
clean-dessin classical RG models. The molecular corpus includes designed
amides and peptides, with explicit-H MMFF94 calculations, geometric persistence,
and method-labelled conformations. Forty-eight existing Aster MP20 structures,
32 existing Mora PubChem identities and eight public PDB structures supply
separately labelled anchors. Their new computed responses are not new
experimental observations, DFT labels, binding affinities or CISS measurements.

All conformer, gauge, intervention and task descendants share a partition.
Specified physical/chemical similarity edges are closed transitively before
component-level stratification. Large components are kept intact even when
80/10/10 proportions or rare-class coverage become impossible. This gives a
finite, threshold-relative separation statement; it cannot guarantee every
kind of statistical independence or prevent overfitting. Source anchors are
train only. Molecular anchors are also checked against every default heldout
using the same chemical and applicable geometric thresholds; conflicting parent
groups are excluded. Global exposure to the models' previous corpora is a separate
evaluation gate.

For a number-conserving one-particle Hamiltonian $H$, write
$P_{ij}=\langle c_j^\dagger c_i\rangle$ and $U(t)=e^{-itH}$.
The data use $P(t)=U(t)P(0)U(t)^*$, occupations $n_i=P_{ii}$ and
outgoing currents $J_{ij}=-2\operatorname{Im}(H_{ij}P_{ji})$.
Then $\dot n_i=-\sum_jJ_{ij}$, directly from $i\dot P=[H,P]$.
Paired BdG evolution includes its independently computed condensate source;
it is not falsely treated as electron-number-conserving normal transport.
An exact interacting Hubbard dimer supplements Gaussian exchange correlations.
Molecular nuclear dynamics uses MMFF94, with a separately labelled,
uncalibrated effective $\pi$-orbital model for charge modulation.

Parquet configurations retain condition, target and evaluation roles, source
versions/hashes, units, failures and audit evidence. The local
\href{http://127.0.0.1:8765/datasets.html}{dataset explorer} shows actual
structures, currents, time evolution, graph wavelets and checked circuit
compilations. These computations support candidate discovery and model
validation; the independent arithmetic determinant/zero-location bridge remains
open. No new whole-chapter Lean formalization is asserted.

\section{Molecular structure, complex transport and learned chemistry}
The complementary-data laboratory now separates deposited protein structure,
actual finite quantum transport and learned conformation generation. Covalent
bonds use explicit molecular graphs or qualified PDB templates, with depth
ordering and visible atom/bond scales. The ribbon view reads deposited
\texttt{HELIX}/\texttt{SHEET} records; short synthetic peptide backbones remain
unassigned rather than acquiring invented secondary structure.

Five additional finite propagation examples compare matched SSH/QWZ
topological and trivial controls and an SSH domain wall. Complex phase,
probability, current, coherence and selected eigenstates are distinct views.
Analytic hydrogen volumes illustrate nodes, quantum beats and a phase vortex;
they are not molecular electronic-structure calculations. Eight further CPU
runs train classical and two-qubit hybrid score/flow models on periodic MMFF
butane/pentane torsion spaces. All checkpoints, losses, controls and positive
noise endpoints are explicit. These examples neither establish quantum
advantage nor close the independent arithmetic bridge for GRH.

\section{Learning objectives with explicit analytic guarantees}
A useful learned contribution must distinguish approximation of known
$L$-functions from proof of their zero locations. We propose holomorphic
coefficient/readout learning, certified contour losses, arithmetic trace and
functional-equation residuals, Schatten-controlled operator refinement,
certificate search for the original signed moment, and exact-algebra
representation constraints. Each has a different acceptance test.

For an analytic target on $|s-c|\le r$, let $|F|\le M_r$ there and choose
$R<\rho<r$. A degree-$N$ holomorphic readout trained on $m>N$ equally spaced
points of the radius-$\rho$ circle has the explicit uniform error budget
\[
 \|p_k-F\|_{|s-c|\le R}
 \le E_{\rm opt}(k)+E_{\rm oracle}+E_{\rm alias}+E_{\rm Taylor}.
\]
The following proofs give all four terms, a primitive-Dirichlet arithmetic
bound for $M_r$, and an infinite-feature kernel limit with its actual rates.
Thus known analytic families can be approximated to arbitrary requested
accuracy with independently certified evaluation error. This does not prove
GRH. Uniform approximation by an already-real-zero determinant class is a
different, potentially GRH-equivalent requirement; it must not be assumed
through an alleged universal-approximation theorem.

\section{Learning nonlocal strategies and checking their certificates}
\label{qg:section}
A two-player nonlocal game samples questions $x,y$ and scores outputs $a,b$.
The players may share a state, but communicate neither questions nor outcomes
during play. Our finite tensor strategies have
$p(a,b|x,y)=\operatorname{Tr}[\rho(P_a^x\otimes Q_b^y)]$.
Offline joint optimization is allowed and is distinct from communication
during the game. We train local measurement frames and a Schmidt angle using
the exact Born objective, and replay actual checkpoints, probabilities,
correlations, word moments and simulated shots in the interactive laboratory.

For $0\le\alpha\le2$ use the tilted score
\[
 S_\alpha=\alpha\langle A_0\rangle+
 \langle A_0B_0+A_0B_1+A_1B_0-A_1B_1\rangle.
\]
Its classical maximum is $2+\alpha$; its quantum maximum is
$\sqrt{8+2\alpha^2}$ \cite{GameTilted}. At $\alpha=0$ the uniform CHSH winning
probability is $1/2+S_0/8$. This conversion does not apply unchanged to the
tilted objective. The three seeds for each of $\alpha=0,1,3/2$ give nine
trained runs. One $\alpha=3/2$ run remains at the classical value $3.5$;
the other runs reach the stated quantum maxima numerically. This retained
failure rules out a claim of universal convergence of the optimizer.

The magic-square task uses two Bell pairs, local four-dimensional joint
measurements, row parity $+1$ and column parities $(+1,+1,-1)$. Its classical
value is $8/9$ and its quantum value is $1$, as proved below. Three further
trained local-frame runs reach values within $3\cdot10^{-13}$ of one.
All twelve runs are small CPU simulations with saved seeds and parameters,
not hardware Bell experiments or evidence for quantum training advantage.

These examples belong to both the closure of finite tensor correlations
$C_{qa}$ and the commuting-operator set $C_{qc}$. The external theorem
$\mathrm{MIP}^*=\mathrm{RE}$ separates these sets and implies the negative
answer to Connes embedding \cite{GameMIP}. We do not reproduce that theorem.
Finite learned strategies, positive state Gram matrices and finite NPA
constraints do not witness its separation.

\section{Hubs, spokes and the topology retained by a dessin}
\label{hs:section}
The hubs-and-spokes construction in the HoTT book, Sections~6.6--6.8
\cite{HubsSpokes}, gives a useful language for attaching cells. For a loop
$f:S^1\to X$, form the homotopy pushout
$X\cup_f CS^1$: add a fresh hub and paths from every $f(u)$ to that hub,
continuously indexed by $u:S^1$. This attaches a disk, rather than merely a
finite graph of radial edges. The book's Remark~6.7.1 explains why identifying
the hub with an existing boundary point requires extra coherence;
Remark~6.7.2 distinguishes propositional from judgmental computation.

For our dessins the useful distinction is geometric. A vertex star has a
disk neighborhood with a specified cyclic ordering of its incident bands.
Removing the disk center gives a punctured disk, which retains its peripheral
loop. Removing only the vertex of a finite star instead leaves disconnected
open rays. These operations cannot be interchanged. In a cellular dessin,
filling a face imposes its attaching-word relation; leaving its center
punctured retains a boundary component of the ribbon neighborhood.

\begin{center}
\includegraphics[trim=5pt 5pt 5pt 5pt,clip]{/home/iska/Documents/LaTeX-projects/GRH-from-Mac-20261104/outputs/adaptive-theta-articles-20261105/concise-v2/compiled-diagrams/3667c0b10605da38015eb0f87a64df21254f7ab7002b67dae645ab113048e0a6.pdf}
\end{center}
The shaded regions are two-dimensional; the displayed finite spokes sample
them. The annulus shown on the right is a deformation retract of the punctured
disk, not a missing vertex in a one-dimensional graph.

\begin{proposition}[Face filling changes the covering problem]
Let a connected cellular dessin be a spine of an oriented genus-$g$ ribbon
surface with $n>0$ boundary components. Its graph fundamental group is free
of rank $2g+n-1$. Capping all boundary components gives the presentation
\[
 \left\langle a_1,b_1,\ldots,a_g,b_g\ \middle|\
       \prod_{i=1}^g[a_i,b_i]=1\right\rangle.
\]
For the one-face clean dessins $K_{2,3}$ and $K_{2,5}$ used here, these free
ranks are respectively $2$ and $4$, before capping.
\end{proposition}
\begin{proof}
A spanning-tree collapse of the connected spine leaves $E-V+1$ circles.
The ribbon surface retracts onto that spine and has Euler characteristic
$2-2g-n=V-E$, proving the rank formula. Cutting a surface along a standard
system of handle arcs gives the boundary word
$\prod_i[a_i,b_i]\prod_jc_j$. Van Kampen therefore presents its group by
these generators and this one relation. With a boundary present, eliminate
$c_n$, leaving the stated free group. Capping attaches disks along the $c_j$,
so van Kampen sets each $c_j=1$, giving the closed-surface presentation.
For our rotations on $K_{2,m}$, both black vertices have the same cyclic
order on the $m$ incident edges. Suppressing the bivalent white vertices
gives dart rotations $(0\,\cdots\,m-1)(m\,\cdots\,2m-1)$ and involution
$i\leftrightarrow i+m$. The face permutation is one $2m$-cycle for odd $m$.
Thus $V=m+2$, $E=2m$, $F=1$, and $g=(m-1)/2$, giving both examples.
\end{proof}

The once-punctured torus and closed torus both have $H_1\cong\mathbb Z^2$,
but their groups are $F_2$ and $\mathbb Z^2$; only the closed torus has
$H_2\cong\mathbb Z$. The analogous genus-two groups are $F_4$ and the
four-generator surface group, with $H_1\cong\mathbb Z^4$ in both cases.
Indeed the standard closed-surface cellular complex has one vertex,
$2g$ edges and one face. The face word is a product of commutators, so its
abelianized boundary is zero. Removing the face leaves the wedge of circles.
This computes the stated homology and also explains why $H_1$ alone misses
the face relation. The torus group is abelian while $F_2$ is not; in higher
genus the closed surface has nonzero $H_2$ while the graph does not.

Consequently the graph's universal cover is a tree, whereas the universal
cover of the filled genus-two surface includes lifted faces and is a disk.
Collapsing a spanning tree changes the combinatorial cover graph; the Cayley
tree of a free basis is not literally the original biregular dessin cover.
This distinction matters when lifting strings, periodic paths and boundary
conditions in the infinite locally gentle algebra.

The gluing language organizes local patches and their compatibility maps.
It does not reconstruct the ribbon cyclic order, complex structure, Belyi
map, quadratic zero relations, coefficient ring or polynomial/formal/operator
completion from the homotopy type alone. These data are retained explicitly
in the existing dessin and order constructions. Descent of modules, derived
categories, Hom-control, AR meshes or Fukaya categories needs a separately
defined functor and hypotheses; a homotopy pushout of underlying spaces is
not such a theorem. Nor does it imply an arithmetic determinant identity.

The accompanying tool computes the same one-face genus-one/two/four dart
families and draws local disk/annulus patches and their surface polygon
presentations, peripheral loops and face words. These
are exact finite combinatorial checks plus the topological proof above.
HoTT is referenced as a conceptual foundation; no univalence or higher
inductive-type axiom is added to the pinned Lean4 development. Its complete
cell-attachment and categorical formalization remains an explicit task.


% Reviewed analytic targets and original weighted arithmetic moments.
\section{A finite, validated approximation result}
Let $\chi_3=(0,1,-1)$, $\chi_4=(0,1,0,-1)$ and
$\chi_5=(0,1,-1,-1,1)$, extended periodically, with parities $a=1,1,0$.
We approximate the completed functions
\[
 \Lambda_q(s)=(q/\pi)^{(s+a)/2}\Gamma((s+a)/2)L(s,\chi_q).
\]
The primitive theta/Mellin identity is the explicit classical analytic input;
no hypothesis about zeros is used. These three examples do not assert a
uniform result for all conductors, general Hecke functions or Artin holomorphy.

\begin{theorem}[Scope of the supplied finite certificate]
For each $q=3,4,5$, the supplied exact-rational calculation constructs a
complex polynomial $P_q$ of degree $28$ satisfying
\[
 \sup_{|s-1/2|\le7/10}|P_q(s)-\Lambda_q(s)|<5\cdot10^{-15}.
\]
The computation includes enclosures of all $64$ true Fourier-circle nodes,
completed-function sample error, analytic truncation, Fourier aliasing,
the exact closed-form $180$-step gradient-descent iterate and coefficient
midpoint rounding. The Python implementation and informal analytic proof
have independent checks; this is not a kernel-checked Lean theorem.
\end{theorem}
\begin{proof}
Write $F(c+z)=\sum_{n\ge0}a_nz^n$, with $c=1/2$, $R=7/10$,
$\rho=6/5$, $r=12/5$. The independently recomputed rational theta majorants
are strictly below $M_q=4,10,6$, respectively, on $|z|\le r$.
Use true roots $\zeta_j=e^{2\pi i j/64}$ and rational labels $y_j$
with proved error $|y_j-F(c+\rho\zeta_j)|\le\eta$.
Use the real Euclidean gradient on complex coefficient coordinates for
\[
 \mathcal L(b)=\frac1{128}\sum_{j=0}^{63}
 \left|\sum_{n=0}^{28}b_n\zeta_j^n-y_j\right|^2.
\]
Discrete Fourier orthogonality gives the least-squares optimum
$b_n^*=64^{-1}\sum_jy_j\bar\zeta_j^n$, $0\le n\le28$,
and gradient $\nabla\mathcal L(b)=b-b^*$.
For step $1/4$ and zero initialization, the exact $180$-step iterate is
$b_n^{180}=(1-(3/4)^{180})b_n^*$. The code evaluates this closed form;
it does not claim an explicit180-update optimizer run.

Set $u=R/\rho=7/12$, $v=R/r=7/24$ and $w=(\rho/r)^{64}=2^{-64}$.
Let $\widehat b_n$ be an outward rectangle enclosing $b_n^*$ and
$\epsilon_n$ the sum of the coordinate halfwidths of the rectangle
for $b_n^{180}$. Choose each actual polynomial coefficient as that
rectangle's exact rational midpoint $c_n$. The normalized basis is explicit:
\[
 P_q(s)=\sum_{n=0}^{28}c_n\left(\frac{s-1/2}{6/5}\right)^n.
\]
Thus the exported coefficient with index $n$ multiplies the normalized
variable, rather than the raw power $(s-1/2)^n$. Its error on the compact disk is bounded by
\begin{align*}
 E_{\rm opt}&=(3/4)^{180}\sum_{n=0}^{28}|\widehat b_n|_1u^n,&
 E_{\rm arithmetic}&=\sum_{n=0}^{28}\epsilon_nu^n,\\
 E_{\rm oracle}&=\eta\frac{1-u^{29}}{1-u},&
 E_{\rm alias}&=M_q\frac{w}{1-w}\frac{1-v^{29}}{1-v},\\
 E_{\rm Taylor}&=M_q\frac{v^{29}}{1-v}.&&
\end{align*}
Here $|\widehat b_n|_1$ is the sum of the maximum absolute real and imaginary
coordinates, a safe upper bound for every enclosed complex modulus.
Cauchy's estimate $|a_n|\le M_qr^{-n}$ and the Fourier identity
$b_n^*=a_n\rho^n+\sum_{\ell\ge1}a_{n+64\ell}\rho^{n+64\ell}$
for exact labels prove the alias and Taylor bounds. Label perturbations
contribute at most $\eta$ per Fourier coefficient. The remaining two
terms follow directly from the exact iterate and coefficient rectangles.
Each term is rounded upward to a $180$-bit dyadic rational; the reported total
is the exact sum of those reported terms. The certified totals below are
strictly below the asserted threshold. Detailed sample enclosures are proved
in the extended edition and encoded with exact integer endpoints.
\end{proof}

\begin{center}\begin{tabular}{c|cc}Conductor & Maximum sample radius & Uniform polynomial bound\\\hline
3 & $3.026\times10^{-17}$ & $1.786\times10^{-15}$ \\
4 & $7.206\times10^{-17}$ & $4.455\times10^{-15}$ \\
5 & $1.404\times10^{-16}$ & $2.906\times10^{-15}$ \\
\end{tabular}\end{center}
The displayed values are rounded upward; exact rational records establish the strict inequalities.


\section{What this does and does not certify}
The final JSON records exact rational centers, interval endpoints, source
hashes and all five error terms. Float values are display summaries only.
A high-precision comparison is an independent diagnostic, not the reason
an interval contains its mathematical target. Arithmetic uses outward
integer operations, and invalid float/order inputs are rejected, including
through typed caches. Earlier pilots and failed validation probes are retained.

This is a fixed finite tolerance/configuration. Increasing accuracy needs
new truncation and precision choices; the fixed $N=11,T=128$ interface is
not an arbitrary-accuracy oracle. The broader approximation existence theorem
is separate. The certificate says nothing about the location of zeros:
$P_q$ is not required to be real-rooted or a self-adjoint spectral determinant.
The independent arithmetic determinant identity, positivity and original
signed short-row/covariance estimates remain open. No GRH proof or complete
new Lean4 formalization is claimed.

\section{Requested-error readouts from preserved certified labels}\label{ar:requested}
The accepted final-v5 sample certificate is preserved unchanged. For each
$q=3,4,5$, it supplies $m=64$ rational complex labels $y_j$, bounds
$|y_j-\Lambda_q(c+\rho\zeta_j)|\le\eta_j$, and outward rectangles enclosing
the true roots $\zeta_j=e^{2\pi i j/m}$, where $c=1/2$, $\rho=6/5$.
The classical primitive theta/Mellin identity and accepted growth bounds
$M_q=4,10,6$ on $|s-c|\le r=12/5$ remain the analytic inputs.
This extension accepts an exact rational error request $\varepsilon>0$ on
the fixed compact $|s-c|\le R=7/10$. It changes the finite readout, while
retaining the certified sample errors and the $64$-node alias budget.
It is not an arbitrary-accuracy theta-sampling interface.

\begin{theorem}[Executed rational readout with explicit refusal]\label{ar:returned}
For each supported conductor and rational request, the finite-data
selection below either returns a normalized rational polynomial $P$ and an
executed, exactly replayable gradient-descent trace with
\[
 \sup_{|s-c|\le R}|P(s)-\Lambda_q(s)|\le\varepsilon,
\]
or refuses the request with the exact minimum base budget over all
$64$ allowed degrees. Refusal concerns attainability by this particular
fixed-data majorant; it is not an impossibility theorem for approximation.
\end{theorem}
\begin{proof}
Let $B_n$ be an outward rectangle enclosing
$b_n^*=m^{-1}\sum_j y_j\bar\zeta_j^n$, $0\le n<64$, computed from the
preserved root rectangles. Put $t_n=\operatorname{mid}(B_n)$ and let
$e_n$ be the sum of its coordinate halfwidths. These are exact rationals,
and $|t_n-b_n^*|\le e_n$.
For a selected $d<64$ execute, using exact rational operations,
\[
 b_n^0=0,\qquad b_n^{k+1}=\tfrac34 b_n^k+\tfrac14 t_n,
 \qquad 0\le n\le d.
\]
This is genuine gradient descent with step $1/4$ for the explicitly
specified real Euclidean objective
$J(b)=\frac12\sum_{n=0}^d|b_n-t_n|^2$, whose gradient is $b-t$.
All iterates and exact objective values are retained. The output is literally
\[
 P_k(s)=\sum_{n=0}^d b_n^k\left(\frac{s-c}{\rho}\right)^n.
\]
Set $u=R/\rho=7/12$, $v=R/r=7/24$, $w=(\rho/r)^{64}=2^{-64}$,
$S_d(x)=\sum_{n=0}^d x^n$ and $\eta=\max_j\eta_j$. The five budgets are
\begin{align*}
 E_{\rm opt}(k)&=\sum_{n=0}^d|b_n^k-t_n|_1u^n,&
 E_{\rm oracle}&=\eta S_d(u),\\
 E_{\rm alias}&=M_q\frac{w}{1-w}S_d(v),&
 E_{\rm Taylor}&=M_q\frac{v^{d+1}}{1-v},\\
 E_{\rm arithmetic}&=\sum_{n=0}^d e_nu^n.&&
\end{align*}
Here $|x+iy|_1=|x|+|y|$. Cauchy's estimate and true-root Fourier
orthogonality give the alias and Taylor terms; certified label errors give
the oracle term. Triangle inequality applied to
$b_n^k-b_n^*=(b_n^k-t_n)+(t_n-b_n^*)$ gives the other two.
Thus their exact sum bounds the compact-uniform error.

Define $C_d=E_{\rm oracle}+E_{\rm alias}+E_{\rm Taylor}
+E_{\rm arithmetic}$. Select the smallest $d<64$ with $C_d<\varepsilon$;
if none exists, report $\min_{d<64}C_d$ and refuse. For a selected degree,
$E_{\rm opt}(k)=(3/4)^k\sum_n|t_n|_1u^n\longrightarrow0$.
Therefore executing updates until
$E_{\rm opt}(k)\le\varepsilon-C_d$ terminates after finitely many steps.
The serialized coefficients, budgets, targets and trace use exact rational
endpoints; displayed decimal approximations are not proof inputs.
\end{proof}

\section{Requested-error theta samples for three fixed conductors}\label{at:sampling}
For the primitive characters $\chi_3=(0,1,-1)$, $\chi_4=(0,1,0,-1)$,
$\chi_5=(0,1,-1,-1,1)$ with parities $a=1,1,0$, retain the classical
primitive theta/Mellin identity and root number one already stated above.
Put $D=[-1,2]+i[-4/3,4/3]$ and $\delta=3/q$.
The following adaptive extension is implemented with independent interval
and sample checks; it is an informal analytic theorem, not a Lean result.
\begin{theorem}[Finite adaptive sample specification]\label{at:sample}
For every supported conductor, exact rational $s\in D$, and rational
$\eta>0$, the choices below terminate and give a rational complex midpoint
$y$ and exact rational radius $R\le\eta$ with
$|y-\Lambda_q(s)|\le R$.
The same conclusion holds for a specified true circle node when its
root and node enclosures are refined on every precision retry.
\end{theorem}
\begin{proof}
For $x\ge0$ define the exact rational bound $E(x)=2^{-\lfloor x\rfloor}$.
Since $e>2$, $e^{-x}\le E(x)$; this bound is nonincreasing and tends to zero.
Let $Q(x)=x^{-1}+2x^{-2}+2x^{-3}$ and
$C_0=1$, $C_1=(1+\delta^{-1})/2$.
For $N\ge1$, $n_0=N+1$, $J\ge1$, $T=2^J$, use
\begin{align*}
 B_n&=2C_aQ(\delta n_0^2)
       \frac{E(\delta n_0^2/2)}{1-E(\delta n_0)},\\
 B_t&=\sum_{n=1}^N2n^aT^3E(\delta n^2T)Q(\delta n^2T).
\end{align*}
Here $\delta n_0\ge6/5$ ensures a positive denominator; $E(x)$ equals
one for $0\le x<1$, so this safeguard is essential.
On the four panels of each $[2^j,2^{j+1}]$, write midpoint $m$,
halfwidth $h$, $\rho=m/2$, $v=h/\rho\le2/9$,
and $B(m)=\max\{(m+\rho)^2,(m-\rho)^{-2},1\}$.
Use the Cauchy remainder
\[
 B_Q=\sum_{\mathrm{panels}}\sum_{n=1}^N
 \frac{2n^a\,2h\,3B(m)E(\delta n^2m/2)v^{2K+2}}
 {(2K+3)(1-v^2)}.
\]
Choose in turn the first $N\ge1$, first $J\ge1$, and first $K\ge0$
for which their respective budgets are at most $\eta/8$.
Gaussian decay, then exponential time decay, then $v\le2/9$ prove
termination. These are the accepted tail and Cauchy arguments with each
true exponential separately replaced by its upper bound $E$.

There are $P=4JN$ potential pairs. A nonzero-character pair may be
omitted only if its entire two-power contribution has bound
\[
 b=4hn^a(m+h)^2E(\delta n^2(m-h))\le\eta/(8P).
\]
Add every omitted bound to $B_{\rm discard}$, giving
$B_{\rm discard}\le\eta/8$.
Evaluate all other integrated Taylor polynomials by the stated exact
coefficient recurrence using outward rectangles. Refine both interval
precision and all primitive series remainders until the final rectangle's
sum of coordinate halfwidths satisfies $R_{\rm arith}\le\eta/4$.
For this fixed finite expression all input and primitive intervals shrink;
continuity and outward arithmetic imply $R_{\rm arith}\to0$.
Precision doubling therefore terminates, and the sum of the five raw
components is at most $3\eta/4$.
Choose an integer output scale $B$ such that $5\,2^{-B}\le\eta/8$ and
round each component upward by
$\operatorname{up}_B(x)=2^{-B}\lceil2^Bx\rceil$.
The reported radius is their literal exact sum. With
$\mathbf b=(R_{\rm arith},B_n,B_t,B_Q,B_{\rm discard})$, it obeys
\[
 |y-\Lambda_q(s)|\le\sum_{i=1}^5\operatorname{up}_B(b_i)
 \le\frac{3\eta}{4}+5\,2^{-B}\le\frac{7\eta}{8}<\eta.
\]
The output retains exact rational midpoints, rounded components, $B$,
raw interval radius, parameters, input metadata and source hashes.
The raw analytic bounds are exactly reconstructible; they need not be
serialized as potentially enormous fractions. The output precision varies
with the request and introduces no fixed reporting floor.
\end{proof}

This fixed-conductor sample result uses no hypothesis on zeros.
It does not remove a downstream fixed-node Fourier alias floor, supply
an all-conductor theorem or a positive arithmetic determinant, or prove GRH.

\section{An end-to-end adaptive analytic readout}\label{at:family}
The new driver changes the sample truncations, interval precision, circle
size and degree; it does not reuse the frozen labels to assert a smaller
error. Its scope is still the three stated primitive quadratic characters
and the fixed compact $|s-1/2|\le7/10$.

\begin{theorem}[Requested-error construction in the stated finite family]\label{at:requested}
Given a positive rational $\varepsilon$ and $q\in\{3,4,5\}$, the specified
exact-arithmetic algorithm terminates in the mathematical unbounded-resource
model and returns a normalized rational polynomial $P$ with
\[
 \sup_{|s-1/2|\le7/10}|P(s)-\Lambda_q(s)|\le\varepsilon.
\]
Its sample enclosures, rational optimizer updates and all error budgets are
explicit. This is an approximation theorem, with the stated classical theta
identity as an input; it supplies no zero-location or positive determinant
conclusion and is not a Lean theorem or a practical runtime bound.
\end{theorem}
\begin{proof}
Put $u=7/12$, $v=7/24$. Choose $d\ge0$ so
$10v^{d+1}/(1-v)\le\varepsilon/8$, and choose $m>d$ so
\[
 10\frac{2^{-m}}{1-2^{-m}}\frac1{1-v}\le\varepsilon/8.
\]
The driver rounds $m$ upward to a multiple of eight, which only decreases
this bound. Set the requested sample error
$\eta=\varepsilon(1-u)/8$. The adaptive sampler gives all $m$ true circle
nodes and labels with error at most $\eta$. Form their outward Fourier
rectangles through degree $d$, and refine the nodes and finite evaluation
arithmetic if their weighted midpoint radii exceed $\varepsilon/8$.
This refinement terminates: for these fixed finite $d,m$, node uncertainty
and every finite-expression arithmetic width tend to zero with precision.
Merely lifting fixed boxes would not suffice, so the driver regenerates the
nodes and samples on this retry.

Choose rational midpoint targets $t_n$ and execute the previously specified
quadratic-objective updates until
$\sum_{n=0}^d|b_n^k-t_n|_1u^n\le\varepsilon/4$.
Their geometric contraction proves finite termination. The oracle,
aliasing, Taylor and Fourier-midpoint terms are each at most
$\varepsilon/8$, and the optimization term is at most $\varepsilon/4$.
Hence the raw five-term sum is at most $3\varepsilon/4$.
Choose a serialization precision $B$ with $5\,2^{-B}\le\varepsilon/8$,
round each term upward, then add the five reported rationals exactly.
The printed bound is at most $7\varepsilon/8\le\varepsilon$.
The polynomial uses $((s-1/2)/(6/5))^n$, and the preceding Fourier/Cauchy
proof gives its asserted compact-uniform bound.
\end{proof}

The trace is losslessly encoded by the fixed target vector and a rational
scale extracted from each \emph{executed} coefficient state, together with
its objective, weighted residual and full-state hash. Every coefficient
can be recovered exactly; a separate replay expands the states and checks
every update. The producer does not generate its trace from a closed form.
The saved objective remains the rational Fourier-midpoint quadratic, with
its perturbation from ideal-root least squares explicitly budgeted.

% Accepted prose AFM proof; explicit external inputs, not a Lean theorem.
\section{An arithmetic-family moment with every original row retained}
\label{mm:core}
Let $K=\mathbb Q(\sqrt{-3})$, $\mathcal O=\mathbb Z[\omega]$, and write
$N$ for the ideal or element norm. Fix a finite-order Hecke character
$\nu$ and a finite set $S$ of prime ideals containing those above $6$ and
those dividing its conductor. For $0\ne u\in\mathcal O$, let $\kappa_u$
be the character of $K(u^{1/6})/K$ defined using arithmetic Frobenius, and
put $\psi_u=\operatorname{prim}(\nu\kappa_u)$. Here
$\operatorname{prim}$ means the underlying primitive Hecke character:
it is formed before introducing any divisor variable or evaluating
conductor-dependent zero extensions. The rows below are all actual
Eisenstein elements, including units and prime powers supported in $S$.

The external inputs are abelian Kummer reciprocity and character
uniqueness, the primitive Hecke functional equation and its usual
vertical growth, Stirling's formula, scalar Dedekind-zeta convexity, and
Blomer--Goldmakher--Louvel's large sieve for their specified sixth-order
family. We use the pinned
\href{https://arxiv.org/pdf/1112.1650v1}{arXiv:1112.1650v1},
Theorem~1.3, Lemma~2.1 and Appendix~II; its Section~5 supplies the
all-ideal decomposition whose fixed-twist proof is redone here.
No statement about all primitive characters of comparable conductor,
Artin holomorphy, Lindel\"of bounds or zero locations is assumed.

\begin{theorem}[The original-family numerator moment]
\label{mm:afm}
For every $\varepsilon>0$, uniformly for $H\ge1$ and $t\in\mathbb R$,
\begin{equation}\label{mm:central}
 \sum_{0<Nu\le H}|L(1/2+it,\psi_u)|^2
 \ll_{\nu,S,\varepsilon}H^{1+\varepsilon}(1+|t|)^{1+\varepsilon}.
\end{equation}
For every fixed $0<\delta<1/2$, the same estimate holds with
$1/2+it$ replaced by $1/2+\delta+it$, with the constant also depending
on $\delta$. If the central bound has constant $C_0$ at an exponent
$0<\varepsilon_0\le1/4$, an admissible off-central constant is
\begin{equation}\label{mm:poisson-constant}
 200 C_0\left(\frac{3/2+\delta}{1/2-\delta}\right)^2.
\end{equation}
\end{theorem}
\begin{proof}
For the canonical BGL character $\vartheta_{\mathfrak a}$,
Appendix~II gives a radicand $\alpha_{\mathfrak a}$ with
$(\alpha_{\mathfrak a})_S=\mathfrak a\mathfrak g^{-6}$.
Since $\mathcal O_S$ is principal, comparison with any literal $u$
whose outside-$S$ ideal is $\mathfrak a$ gives
$\kappa_u=\eta\vartheta_{\mathfrak a}$ for one of finitely many fixed
$S$-unit Kummer twists $\eta$. This is an equality of primitive
characters, established outside a finite set and extended by uniqueness.
For each outside-$S$ ideal $\mathfrak a$ and $S$-supported integral
ideal $\mathfrak b$, at most six elements generate
$\mathfrak a\mathfrak b$. The resulting $S$-supported multiplicity sum
converges at exponent $1+\varepsilon$.

For each fixed $\xi=\nu\eta$, split the canonical parameter as
$\mathfrak a=\mathfrak a_0\mathfrak a_1\mathfrak a_2^6$, where
$\mathfrak a_0$ is squarefree, $\mathfrak a_1$ is squarefull and
sixth-power-free, and $(\mathfrak a_0,\mathfrak a_1)=1$.
The primitive conductor is at most a fixed multiple of
$N(\mathfrak a_0\mathfrak a_1)$.
The common-gamma approximate functional equation, with
$X=H/(N\mathfrak a_2)^6$ and $T=1+|t|$, has common maximal length
$B=X^{1/2}T(HT)^\eta$ for any sufficiently small fixed $\eta>0$.
Mellin inversion removes its row-dependent conductor weight.
Split a column as $\mathfrak b_S\mathfrak b_1\mathfrak b_2$, with
$\mathfrak b_S$ supported in $S$, $\mathfrak b_1$ squarefree and
$\mathfrak b_2$ squarefull, the last two coprime and prime to $S$.
The original primitive values at $\mathfrak b_S$ remain bounded row
multipliers. The $S$-supported and squarefull column weights are summable
at $1/2+\eta$. Multiplying squarefree column coefficients by the fixed
$\xi$ preserves their squared norm, whatever its finite order.
BGL's large sieve therefore bounds the energy for fixed
$\mathfrak a_1,\mathfrak a_2$ by
\[
 (HT)^{O(\eta)}\{M+B+(MB)^{2/3}\},\qquad
 M=X/N\mathfrak a_1.
\]
Squarefull counting and summation give
$X+B\sqrt X+(XB)^{2/3}\ll (HT)^{O(\eta)}XT$.
The $\mathfrak a_2$ sum converges at exponent $6$.
Principal rows are separated before the pole-free functional equation
and bounded by scalar zeta convexity. This proves~\eqref{mm:central}.

For the off-central statement divide each $L(s,\psi)$ by
$(s+1)^{5/4}$ and, on principal rows, also multiply by $(s-1)/(s+1)$.
These are bounded analytic functions on $\operatorname{Re}s>1/2$,
continuous on the boundary. Poisson representation and Jensen's
inequality transfer the central squared moment. The convolution of
$P_\delta(y)=\delta/[\pi(\delta^2+y^2)]$ with
$(1+|t|)^{-3/2+\varepsilon_0}$ is at most eleven times that weight
when $0<\varepsilon_0\le1/4$. Restoring the denominator and the
principal pole-removal factor gives~\eqref{mm:poisson-constant}.
The extended edition gives the detailed derivation with all primitive local values.
\end{proof}

\paragraph{Separate fixed-unitary extension.}
The moment theorem also holds for any one fixed unitary Hecke character
$\nu$. For precision choose
$\nu(\mathfrak a)=\nu_0(\mathfrak a)(N\mathfrak a)^{-i\tau}$, with
angular parameter $m\in\mathbb Z$ and completed gamma factor
\[
 \gamma_\nu(s)=2(2\pi)^{-(s+|m|/2+i\tau)}
                    \Gamma(s+|m|/2+i\tau).
\]
All finite Kummer twists have this same infinity type. Use
$T_\nu=1+|t+\tau|+|m|$, comparable to $1+|t|$ with
$\nu$-dependent constants. Pure norm rows have
$L(s,\psi)=\zeta_K(s+i\tau)$ and are separated before the pole-free
functional equation; their pole at $1-i\tau$ is removed by
$(s-1+i\tau)/(s+1+i\tau)$. No uniformity in $m$, $\tau$ or the
conductor of $\nu$ is asserted. This extension of the moment is separate
from the following finite-order type-I residue formula.

\subsection{The exact masked remainder and the remaining loss}
Retain the original zero-extended symbol $\chi_d(u)=(u/d)_6$ on good
squarefree ideals. Let $\mu$ be the ideal M\"obius function and set
\[
 \beta_U(d)=\sum_{ab=d\atop Na,Nb\le U}\mu(a)\mu(b),\quad
 T(u,d)=S\cup\operatorname{supp}(u)\cup\operatorname{supp}(d).
\]
Let $\psi_{u,2}=\operatorname{prim}(\psi_u^2)$, formed independently
of squaring any zero sentinel. With $\mathfrak f_u$ and
$\mathfrak f_{u,2}$ the primitive conductors, the full-mask series is
\begin{align}
 F_{u,d}(s)
 &=\frac{L(s,\psi_u)}{L(2s,\psi_{u,2})}C_{T(u,d)}(s),\label{mm:masked}\\
 C_T(s)
 &=\frac{\prod_{p\in T,\ p\nmid\mathfrak f_u}
              (1-\psi_u(p)(Np)^{-s})}
         {\prod_{p\in T,\ p\nmid\mathfrak f_{u,2}}
              (1-\psi_{u,2}(p)(Np)^{-2s})}.\notag
\end{align}
Initially it equals
$\sum_{c\ \mathrm{squarefree},\ (c,T)=1}\psi_u(c)(Nc)^{-s}$
for $\operatorname{Re}s>1$. Every prime of $T$ remains deleted,
including those forgotten by either primitive conductor.

For finite-order $\nu$, $D\ge1$, $U\ge1$, $U^2\le D$, and
$W\in C_c^\infty((0,\infty))$ supported in $[1,2]$, put
\[
 \widehat W(s)=\int_0^\infty W(x)x^{s-1}\,dx,\qquad
 \mathcal M_4(W)=\sum_{j=0}^4\sup_x|(x\partial_x)^jW(x)|.
\]
On $\sigma=1/2+\delta$, the exact residue-subtracted type-I remainder is
\begin{equation}\label{mm:remainder}
 E_u=\frac{D^\delta}{2\pi}\int_{\mathbb R}
 \widehat W(\sigma+it)D^{it}
 \sum_{d\ \mathrm{squarefree}\atop(d,S)=1}
 \frac{\beta_U(d)\nu(d)\chi_d(u)}{(Nd)^{\sigma+it}}
 F_{u,d}(\sigma+it)\,dt.
\end{equation}
This definition retains the actual $\beta_U$, all phases, $d$-masks,
conductors and the independently primitive square.

\begin{theorem}[An improved masked residual estimate]
\label{mm:residual}
Under the preceding finite-order hypotheses, for every $\varepsilon>0$,
\begin{equation}\label{mm:energy}
 \sum_{0<Nu\le H}|E_u|^2
 \ll_{\nu,S,\delta,\varepsilon}
 \mathcal M_4(W)^2D^{2\delta}U^{2-4\delta}H(HD)^\varepsilon.
\end{equation}
\end{theorem}
\begin{proof}
The same deleted set in both factors of~\eqref{mm:masked} gives
$|F_{u,d}(\sigma+it)|\ll (Nu\,Nd)^\eta|L(\sigma+it,\psi_u)|$:
the reciprocal incomplete Euler product at $2\sigma>1$ is bounded
by $\zeta_K(2\sigma)$, and its remaining finite numerator is bounded
by a small power of its radical. Apply Minkowski to~\eqref{mm:remainder}
in the exact row space and use Theorem~\ref{mm:afm}. The actual coefficient
sum is at most
\[
 \sum_d\frac{|\beta_U(d)|}{(Nd)^{\sigma-\eta}}
 \le\left(\sum_{Na\le U}(Na)^{-\sigma+\eta}\right)^2
 \ll U^{1-2\delta+2\eta}.
\]
Four Mellin integrations by parts bound the remaining height integral
by a constant times $\mathcal M_4(W)$. Square the result and absorb
the small losses using $U^2\le D$.
\end{proof}

At $H=D^{3/4}$ and $U=D^a$ with fixed $0<a\le1/2$, the upper bound's excess
over $H$ is $2\delta+(2-4\delta)a$, up to small losses; as
$\delta\downarrow0$ it still leaves $2a$. This is a limitation of the
upper bound, not an energy lower bound. The extracted principal residue,
the original short term at smaller scales and the signed type-II minus
type-I covariance are not controlled by~\eqref{mm:energy}. The next
arithmetic task is to estimate the complete $d$-weighted moment before
its absolute coefficient sum, or establish sufficient signed covariance.
No positive arithmetic determinant, zero-location theorem or GRH proof
follows. The analytic argument uses the external classical inputs above;
finite diagnostics and selected existing Lean local identities do not
formalize the Hecke functions, approximate functional equation, Poisson
transfer or these moment bounds.

\section{Recombining the actual weighted cutoff}
\label{rc:core}
Keep the finite-order Eisenstein family and original masks of
Section~\ref{mm:core}. For $U\ge1$ and a squarefree ideal $m$, put
\[
 A_U(m)=\sum_{d\mid m}\beta_U(d)
 =\sum_{ab\mid m\atop(a,b)=1,\ Na,Nb\le U}\mu(a)\mu(b).
\]
Every summation variable here is an integral ideal. The last formula
uses squarefreeness of $m$; an unconstrained square of a M\"obius
polynomial would also include forbidden overlapping supports.
\begin{theorem}[Exact recombination and an absolute-half-plane rate]
\label{rc:recombination}
Write $T_0=S\cup\operatorname{supp}(u)$ and
$J_{u,U}(s)=\sum_{d\ \mathrm{squarefree},\ (d,S)=1}
\beta_U(d)\nu(d)\chi_d(u)(Nd)^{-s}F_{u,d}(s)$.
For $\operatorname{Re}s>1$,
\begin{align*}
 J_{u,U}(s)&=\sum_{m\ \mathrm{squarefree},\ (m,T_0)=1}
                   A_U(m)\psi_u(m)(Nm)^{-s},\\
 |A_U(m)|&\le3^{\omega(m)},\qquad
 A_U(m)=\mu(m)\quad(Nm\le U).
\end{align*}
For every $\kappa>0$ and $0<\epsilon<\kappa$, uniformly for
$\operatorname{Re}s\ge1+\kappa$, all original rows and all $U\ge1$,
\[
 \left|J_{u,U}(s)-\frac1{L^{T_0}(s,\psi_u)}\right|
 \le U^{-\kappa+\epsilon}
 \{\zeta_K(1+\epsilon)^3+\zeta_K(1+\epsilon)\}.
\]
Here $L^{T_0}$ removes every prime of $T_0$, including conductor-drop
primes; the bound is in the absolute-convergence half-plane only.
\end{theorem}
\begin{proof}
On a surviving term $m=dc$, the two squarefree ideals $c,d$ are
coprime, and both avoid $T_0$. The original product
$\nu(d)\chi_d(u)\psi_u(c)$ is then $\psi_u(m)$. Grouping by $m$
gives $A_U$. Each prime of $m$ has three possible positions in
$(a,b,m/(ab))$, so the absolute sum is at most $3^{\omega(m)}$.
When $Nm\le U$ all divisors fit both cutoffs; the primewise sum is
$1-1-1=-1$, giving $\mu(m)$.

The same bound justifies absolute convergence and regrouping for
$\operatorname{Re}s>1$, because on squarefree ideals
$3^{\omega(m)}=\tau_3(m)$ and
$\sum_m\tau_3(m)(Nm)^{-v}=\zeta_K(v)^3$ for $v>1$,
where this latter sum runs over all nonzero integral ideals.
Subtract the absolutely convergent M\"obius series for $1/L^{T_0}$.
Only $Nm>U$ remain. Rankin's inequality with exponent
$\eta=\kappa-\epsilon$ bounds their absolute sum by
\[
 U^{-\eta}\sum_m\{\tau_3(m)+1\}
                  (Nm)^{-1-\epsilon},
\]
which proves the stated constant independently of deleted primes,
row multiplicities or character phases.
\end{proof}

For $W$ supported in $[1,2]$, the same exact regrouping gives the
original normalized type-I polynomial
\[
 I_u=D^{-1/2}\sum_{m\ \mathrm{squarefree},\ (m,S)=1}
 A_U(m)\nu(m)\chi_m(u)W(Nm/D).
\]
This formula retains all original zero sentinels. It permits a direct
family large sieve before any absolute $d$ sum; its quantitative moment
is developed separately below. The displayed absolute-half-plane rate
alone is not that critical-strip moment.

% Original masked cutoff family; no new Lean or GRH theorem.
\section{A cutoff-uniform envelope for the original masked remainder}
\label{wc:core}
Fix $K=\mathbb Q(\omega)$, $\mathcal O=\mathbb Z[\omega]$, a finite-order
Hecke character $\nu$, and a finite set $S$ containing the primes above $6$
and the conductor of $\nu$. All unspecified ideal sums below are over
nonzero integral ideals. Let $\chi_m(u)=(u/m)_6$ be the original symbol,
zero when $(m,u)>1$, and put
\[
 \beta_U(d)=\sum_{ab=d\atop Na,Nb\le U}\mu(a)\mu(b),\qquad
 A_U(m)=\sum_{d\mid m}\beta_U(d).
\]
The columns $d,c,m$ are squarefree and prime to $S$ where indicated.
Retain the primitive characters $\psi_u=\operatorname{prim}(\nu\kappa_u)$,
$\psi_{u,2}=\operatorname{prim}(\psi_u^2)$, and the full deleted set
$T(u,d)=S\cup\operatorname{supp}(u)\cup\operatorname{supp}(d)$.
For $H,D,U\ge1$, $U^2\le D$, and $W\in C_c^\infty((0,\infty))$
supported in $[1,2]$, the original normalized type-I piece is
\begin{equation}\label{wc:polynomial}
 I_u=D^{-1/2}\sum_{m\ \mathrm{squarefree}\atop(m,S)=1}
     A_U(m)\nu(m)\chi_m(u)W(Nm/D).
\end{equation}
Its exact principal residue is
\begin{align}\label{wc:principal}
 P_u&=\mathbf1_{\psi_u=1}\sqrt D\,\widehat W(1)
 \sum_{d\ \mathrm{squarefree}\atop(d,S)=1}
 \frac{\beta_U(d)\nu(d)\chi_d(u)}{Nd}\rho_{T(u,d)},\notag\\
 \rho_T&=\frac{\operatorname*{res}_{s=1}\zeta_K(s)}{\zeta_K(2)}
               \prod_{p\in T}(1+(Np)^{-1})^{-1}.
\end{align}
Define $E_u=I_u-P_u$; this is the same residue-subtracted remainder as the
fully masked Mellin contour, with no cutoff limit or profile supremum.

\begin{theorem}[Uniformity in the actual truncated cutoff]
\label{wc:uniform}
For every $\varepsilon>0$,
\begin{equation}\label{wc:bound}
 \sum_{0<Nu\le H}|E_u|^2
 \ll_{\nu,S,\varepsilon}\|W\|_\infty^2(HD)^\varepsilon
       \{H+DH^{1/3}+(HD)^{2/3}\}.
\end{equation}
The sum is over every actual nonzero Eisenstein element, including units,
$S$-supported powers and sixth-power rows. The constant is independent of
$U$ under $U^2\le D$. This is an application of the classical sextic
large sieve with exact original coefficients, not a zero-location theorem.
\end{theorem}
\begin{proof}
On an absolutely convergent line, expand the restored squarefree series
$F_{u,d}$. Its condition $(c,T(u,d))=1$ includes $(c,d)=1$, and the
original symbols multiply. Grouping $m=cd$ gives~\eqref{wc:polynomial}
by Mellin inversion. Each prime of $m$ may belong to the first truncated
M\"obius divisor, to the second, or to neither. Consequently
$|A_U(m)|\le3^{\omega(m)}\ll_\eta(Nm)^\eta$, uniformly in $U$.
The length-$D$ coefficient sequence in~\eqref{wc:polynomial} therefore
has squared norm $\ll_\eta\|W\|_\infty^2D^\eta$.

The all-ideal row norm in Section~3 of
\href{https://arxiv.org/pdf/1112.1650v1}{Blomer--Goldmakher--Louvel,
arXiv:1112.1650v1} gives
\[
 B_2^1(X,D)\ll_\eta(XD)^\eta
                 \{X+DX^{1/3}+(XD)^{2/3}\}.
\]
For completeness, split $a=a_1a_2^2\cdots a_5^5a_6^6$ dyadically;
$a_1,\ldots,a_5$ are squarefree and pairwise coprime. If their dyadic
sizes are $A_j$, put $R=\prod_{j=1}^6A_j$ and $M=\max(A_1,A_2)$.
Pivot on the first or second variable, using the squarefree powered
sieve proved simultaneously in that paper's Section~3.2. The other
variables are bounded column multipliers. The block energy is at most
\[
 (XD)^\eta\{R+DR/M+D^{2/3}R/M^{1/3}\}\sum_m|\lambda_m|^2.
\]
The inequalities $R\le X$, $R^3\le XM^3$ and $R^3\le X^2M$
follow from $\prod_jA_j^j\le X$ and the pivot choice. Summing the
logarithmically many blocks proves the displayed bound.

Finite $S$-unit and ray-class partitions convert the original literal
symbol into the canonical column value $\vartheta_m(a)$ with fixed
modulus-one coefficient multipliers. On noncoprime pairs both values
are zero: squarefree $m$ has outside-$S$ conductor exactly $m$.
This orientation preserves sixth-power zeros. At most six elements
generate each outside-$S$ ideal times an $S$-supported ideal; the
$S$-supported sums converge at powers $1$, $1/3$ and $2/3$.
Thus~\eqref{wc:bound} holds first with $I_u$ in place of $E_u$.

A principal $\psi_u$ requires every outside-$S$ valuation of $u$ to
be divisible by $6$. There are $O_{\nu,S}(H^{1/6})$ such literal
rows. Equation~\eqref{wc:principal}, $|\rho_T|\ll_K1$, and the actual
coefficient sum
\[
 \sum_d\frac{|\beta_U(d)|}{Nd}
 \le\left(\sum_{Na\le U}\frac1{Na}\right)^2\ll\log^2(2U)
\]
give $\sum|P_u|^2\ll DH^{1/6}\|W\|_\infty^2\log^4(2U)$.
This is absorbed by the right side of~\eqref{wc:bound}, using
$U^2\le D$ and a smaller internal exponent. Finally
$|E_u|^2\le2|I_u|^2+2|P_u|^2$. The extended proof specifies the
canonical-family transfer and the powered-sieve input precisely.
\end{proof}

Combining this with the previously proved fixed-twist AFM transfer yields
\begin{equation}\label{wc:minimum}
 \sum_{0<Nu\le H}|E_u|^2\ll_{\nu,S,\delta,\varepsilon}
 \mathcal M_4(W)^2(HD)^\varepsilon
 \min\{D^{2\delta}U^{2-4\delta}H,\ H+DH^{1/3}+(HD)^{2/3}\}.
\end{equation}
At $H=D^{3/4}$ and $U=D^a$, $0\le a\le1/2$, the excess over $H$
is at most $\min\{2\delta+(2-4\delta)a,1/2\}$, up to arbitrarily
small losses. The new envelope improves the old one for $a>1/4$ as
$\delta\downarrow0$, but still leaves an exponent loss of $1/2$.
It does not settle the original signed short-row bound, type-II covariance,
positive arithmetic determinant comparison or GRH. The principal-residue
bound used here is sufficient for this envelope, not the stronger target
needed to reassemble the full signed argument.

\section{High cutoffs, actual bilinear factors and a signed covariance test}\label{sb:core}
Retain the original Eisenstein family, fixed finite-order $\nu$, fixed $S$,
common $W\in C_c^\infty((0,\infty))$ supported in $[1,2]$, and every raw
shared-prime zero. Put
\[
 C_U(m)=\sum_{\substack{abc=m\\Na>U,\,Nb>U}}\mu(a)\mu(b),\qquad
 r_U(m)=\#\{(a,b):ab=m,\ Na,Nb>U\}.
\]
All factors are ideals, squarefree and pairwise coprime; on the active
annulus $A_U=C_U-\mu$. A coefficient identity is different from a row-energy bound.

\begin{proposition}[The high-cutoff window]\label{sb:window}
If $D>1$ and $2D/7\le U^2\le D$, then for every surviving squarefree
column $D\le Nm\le2D$,
\[
 C_U(m)=\mu(m)r_U(m),\qquad
 A_U(m)=\mu(m)(r_U(m)-1),\qquad |A_U(m)|\ge1.       
\]
For each fixed nonzero $W$, the normalized squared coefficient norm satisfies
\begin{equation}
 D^{-1}\sum_{\substack{m\ {\rm sf}\\(m,S)=1}}
 |A_U(m)W(Nm/D)|^2\gg_{W,S}1                     \label{sb:coefficient-lower}
\end{equation}
for large $D$, uniformly in this window. This is not a lower bound for row energy.
\end{proposition}
\begin{proof}
The least norm of a nonunit ideal outside $S$ is at least $7$. A long-long
term with $c\ne1$ would have $Nm>7U^2\ge2D$. Thus $ab=m$, including both
endpoints. The fixed-point-free swap $(a,b)\mapsto(b,a)$ makes $r_U$ even,
so $A_U$ is an odd nonzero integer. The last claim follows by counting
squarefree ideals on a fixed interval where $|W|$ is bounded below; the
constant and error are proved in the extended proof.
\end{proof}

The exact remaining polynomial is genuinely bilinear, with $(a,b)=1$,
$U<Na,Nb<2D/U$, $D\le Na\,Nb\le2D$, and
$x_{ab}(u)=x_a(u)x_b(u)$, including zeros. Both factors see the same row
$u=\varepsilon\alpha\beta^2$. Cauchy followed by the accepted powered BGL
bound gives only
\begin{equation}
 Q_{B_X}(C_U,C_U)\ll_{\nu,S,\epsilon}
 (DX)^\epsilon\|W\|_\infty^2
 \{\sqrt D\,X^2+DX+D^{5/6}X^{5/3}\}.              \label{sb:factor-bound}
\end{equation}
At $X\asymp D^{1/4}$ the exponents are $1,5/4,5/4$; the earlier regrouped
bound has $1/2,5/4,13/12$ and is stronger. No $BT_\sigma$ saving is obtained.
Multiplying two marginal operator norms does not supply a coupled fourth moment.

\paragraph{A finite test in the original family.}
There is no automatic nonnegative signed covariance, even when
$C_U(m)\mu(m)\ge0$. Take $\nu=1$, $S$ the primes above $2,3$, $D=400$,
$U=17$, and all $72$ literal rows of $B_7$. Choose the common nonnegative
smooth profile constructed in the extended proof, with values $10$ at $439/400$, $1$ at
$475/400$, and zero at every other integer norm divided by $400$.
There are exactly four active squarefree ideal columns, of norms
$439,439,475,475$. In that order,
\[
 A=(1,1,1,1),\quad C=(0,0,2,2),\quad
 \mu=(-1,-1,1,1),\qquad
 \Re G_{B_7}=
 \begin{pmatrix}72&0&0&18\\0&72&18&0\\0&18&72&0\\18&0&0&72\end{pmatrix}.
\]
Exact raw-symbol enumeration, independently reproducible from the ideal
and residue-field definitions, gives
\begin{equation}
 Q(C,\mu)=-\frac{27}{25},\qquad
 Q(C,C)=\frac{36}{25},\quad
 Q(\mu,\mu)=\frac{864}{25},\quad
 Q(A,A)=\frac{954}{25}.                          \label{sb:covariance-example}
\end{equation}
The real Gram matrix is positive definite, with eigenvalues $54,90$;
its positivity does not force a cross covariance to be positive.
This fixed finite example refutes that shortcut, not the asymptotic
$BT_\sigma$ target or the original signed mean-square estimate.

\paragraph{The two-cutoff contour remains coupled.}
For $T_0=S\cup\operatorname{supp}(u)$, $\psi=\psi_u$ and independently
primitive $\psi_2=\operatorname{prim}(\psi^2)$, put $q_p=\psi(p)(Np)^{-s}$,
$\alpha_p=(Np)^{-z}$, $\beta_p=(Np)^{-w}$ outside $T_0$. The actual
three-placement Euler factor is $1+q_p-q_p\alpha_p-q_p\beta_p$. Its exact
factorization is
\begin{equation}
 \mathcal F(s,z,w)=
 \frac{L^{T_0}(s,\psi)}{L^{T_0}(s+z,\psi)L^{T_0}(s+w,\psi)}
 \prod_{p\notin T_0}\left(1-
 \frac{q_p^2(1-\alpha_p)(1-\beta_p)}
 {(1-q_p\alpha_p)(1-q_p\beta_p)}\right).           \label{sb:perron-factor}
\end{equation}
Exact double Perron inversion is valid on positive auxiliary lines
$c,d>1-\Re s$ for $1/2<\Re s<1$, retaining the original principal pole
exclusion. These safe lines reproduce the previous $U^{2-4\delta}$
energy loss at $\Re s=1/2+\delta$, with the same common-profile derivative
and fixed-margin dependencies. The correction product alone does not
justify moving those lines through reciprocal-$L$ poles. In fact
$\mathcal F(s,z,0)=1/L^{T_0}(s+z,\psi)$ as a meromorphic identity.
No zero-free region, reciprocal moment, analytic saving, new Lean
formalization or GRH proof is supplied by this construction.

\section{A quantitative covariance budget without a new exponent}\label{cbq:core}
Retain the original Eisenstein sextic symbols, fixed finite-order $\nu$,
fixed $S$, one possibly complex common $W\in C_c^\infty$ supported in
$[1,2]$, all squarefree columns prime to $S$, and every raw shared-prime
zero. In the accepted window $2D/7\le U^2\le D$, put $r_m=r_U(m)$,
$\lambda_m=\mu(m)\nu(m)W(Nm/D)$ and $v_m=\overline{\lambda_m}$.
Thus $C_U=\mu r$ and $A_U=\mu(r-1)$ on the active annulus.
For a finite literal row block $B$, let
\[
 G_{mn}=\sum_{u\in B}x_m(u)\overline{x_n(u)},\qquad
 R=\operatorname{diag}(r_m),\qquad
 M_j=\sum_mr_m^j|W(Nm/D)|^2\quad(0\le j\le2).
\]
The primitive $\psi_u,\psi_{u,2}$ and full deleted sets retain their
previous meanings; no conductor-drop value replaces a raw zero.

\begin{proposition}[An actual-vector lower budget]\label{cbq:budget}
If $0\le G\le\Lambda_B I$, then
\begin{equation}\label{cbq:vector}
 \Re Q_B(C_U,\mu)\ge
 -\frac{\Lambda_B}{2D}\{\sqrt{M_0M_2}-M_1\}.
\end{equation}
In this high window, uniformly in $U$,
\begin{equation}\label{cbq:moments}
 \sum_{D\le Nm\le2D}r_m\ll_K D,\qquad
 \sum_{D\le Nm\le2D}r_m^2\ll_K D\log(2D).
\end{equation}
Consequently
\[
 \Re Q_B(C_U,\mu)\ge
 -C_K\|W\|_\infty^2\Lambda_B\sqrt{\log(2D)}.
\]
\end{proposition}
The vector inequality is sharp among positive contractions with the
specified $v,Rv$; no such sharpness is asserted for the arithmetic Gram.
The moment bound counts actual balanced ideal factorizations, rather than
assuming independence of the two row factors.

On $B_X$, inserting the previously accepted powered-BGL norm gives only
\[
 \Re Q_{B_X}(C_U,\mu)\ge
 -C_{\nu,S,\epsilon}\|W\|_\infty^2(DX)^\epsilon
 \{X^2+DX+D^{2/3}X^{5/3}\}.
\]
At $X\asymp D^{1/4}$ the exponents remain $1/2,5/4,13/12$.
At $H=D^{3/4}$ this budget still has the old $D^{1/2}$ excess over $H$;
it supplies no $BT_\sigma$ saving. A negative-part budget is weaker than
$2\Re Q(C_U,\mu)\ge Q(C_U,C_U)+Q(\mu,\mu)-R_\sigma$.

For $S_R=\sqrt R$, the exact identity
\[
 \tfrac12(RG+GR)=S_RGS_R+\tfrac12[S_R,[S_R,G]]
\]
localizes the commutator to unequal factor-count levels, with entry
$(\sqrt{r_m}-\sqrt{r_n})^2G_{mn}$. If $B$ retains all six literal units,
the actual unit sum further gives $G_{mn}=0$ unless $Nm\equiv Nn\pmod{36}$.
These are exact deletions, not a power saving. Even zero different-level
pairings leave the nonnegative within-level energies
$\sum_k(k-1)^2Q_B(\mu\mathbf1_{r=k},\mu\mathbf1_{r=k})$.
The unequal-product coupled fourth moment and its mixed M\"obius term
remain unestimated. No new analytic exponent, energy lower bound,
Lean formalization or GRH conclusion follows.

\section{The fourth-moment conductor and its full raw zero set}
\label{upq:core}
Retain the original Eisenstein sextic symbols $x_m(u)=(u/m)_6$,
fixed finite-order $\nu$, fixed $S$, and one possibly complex common
smooth $W$ supported in $[1,2]$. Put
$\lambda_m=\mu(m)\nu(m)W(Nm/D)$ and $r_m=r_U(m)$.
Columns remain squarefree and prime to $S$. In $2D/7\le U^2\le D$, the actual polynomial is
\[
 \mathcal C_U(u)=
 \sum_{\substack{(a,b)=1,\ Na,Nb>U\\D\le Na\,Nb\le2D}}
 \mu(a)\mu(b)\nu(ab)W(Nab/D)x_a(u)x_b(u).
\]
Its squared row energy has four factor indices $a,b,c,d$, with
$(a,b)=(c,d)=1$. These internal masks are essential. At each prime,
\[
 e_p=\mathbf1_{p\mid a}+\mathbf1_{p\mid b}
       -\mathbf1_{p\mid c}-\mathbf1_{p\mid d}
 \quad\hbox{belongs to }\{0,\pm1\}.
\]
Thus, for the two squarefree products $m=ab,n=cd$, the primitive
oscillatory conductor outside fixed $S$ is $q=mn/(m,n)^2$.
The full raw zero-preserving period is $L=mn/(m,n)$:
\[
 x_m(u)\overline{x_n(u)}
 =\mathbf1_{(u,(m,n))=1}\theta_{m/(m,n),\,n/(m,n)}(u).
\]
No common-prime zero is restored by taking the primitive character.

A generic fourth moment $|\sum_m\lambda_mx_m(u)|^4$ has four individually
squarefree columns $m_1,\ldots,m_4$ with no such internal coprimality.
Its signed local exponent
$e_p=\sum_{i=1}^2\mathbf1_{p\mid m_i}-\sum_{i=3}^4\mathbf1_{p\mid m_i}$
can also be $\pm2$, giving a nontrivial cubic component. For $L=\operatorname{rad}(m_1m_2m_3m_4)$, the exact row-power
interface is
\[
 Q_j=\prod_{\substack{p\mid L\\je_p\not\equiv0\ (6)}}p,\qquad
 Z_j=L/Q_j,\qquad
 \prod_{i=1}^2x_{m_i}(v^j)
       \overline{\prod_{i=3}^4x_{m_i}(v^j)}
 =\theta_{j,e}(v)\mathbf1_{(v,Z_j)=1}.
\]
Here $\theta_{j,e}$ is primitive on $Q_j$, with its own zero extension.
At $j=3$ an exponent $\pm2$ becomes primitive-principal but its raw
zero remains. For the original two-squarefree-product expansion,
$Q_j=q$ for all $1\le j\le5$, including the quadratic $j=3$ direction.
For the accepted row blocks $Nr_j\asymp A_j$, put
$R_0=\prod_jA_j$, $P=\max(A_1,\ldots,A_5)$ and retain
$\prod_jA_j^j\le H$. The full $A_3$ saturation face and all row-factor
masks remain.

The small-conductor pair count recovers an already controlled high-gcd
sector; it is not a new analytic advance. On the remaining original
sector $Nq>P^2$, favorable smooth completion
costs $\min(A_j,\sqrt{Nq})=A_j$ for $1\le j\le5$.
The literal sharp cutoff is not replaced by this smooth trial.
The signed large-conductor fourth-moment and mixed M\"obius sums remain
open; no improved full exponent, $BT_\sigma$, new Lean theorem,
zero-location result or GRH conclusion follows.

% Standalone bibliography fragment for the dessin-resolution/spectral manuscripts.
% Exact reviewed arXiv versions are intentional; do not silently replace with latest.
% Unversioned primary texts are identified by publication record or retrieval pin.
\begin{thebibliography}{99}
\bibitem{SurfaceOrders} A. Schreiber, \emph{Surface Algebras I: Dessins D'enfants, Surface Algebras, and Dessin Orders}, \href{https://arxiv.org/abs/1810.06750v5}{arXiv:1810.06750v5} (2023).
\bibitem{SurfaceDuality} A. Schreiber, \emph{Surface Algebras and Surface Orders II: Affine Bundles on Curves}, \href{https://arxiv.org/abs/1812.00621v1}{arXiv:1812.00621v1}; record 2018, downloaded PDF internally dated November 27, 2024.
\bibitem{Persistence} A. Schreiber, \emph{Multiparameter Persistent Homology---Generic Structures and Quantum Computing}, \href{https://arxiv.org/abs/2210.11433v1}{arXiv:2210.11433v1} (2022).
\bibitem{BessenrodtHolm} C. Bessenrodt and T. Holm, \emph{Weighted locally gentle quivers and Cartan matrices}, \href{https://arxiv.org/abs/math/0511610v1}{arXiv:math/0511610v1} (2005).
\bibitem{LocallyGentleHomology} S. Ford, A. Oswald and J.J. Zhang, \emph{Homological conditions on locally gentle algebras}, \href{https://arxiv.org/abs/2409.08333v3}{arXiv:2409.08333v3} (5 September 2026).
\bibitem{GeometricGentle} S. Opper, P.-G. Plamondon and S. Schroll, \emph{A geometric model for the derived category of gentle algebras}, \href{https://arxiv.org/abs/1801.09659v7}{arXiv:1801.09659v7} (10 July 2025).
\bibitem{Chang} W. Chang, \emph{Geometric models for the algebraic hearts in the derived category of a gentle algebra}, \href{https://arxiv.org/abs/2305.03998v2}{arXiv:2305.03998v2} (4 September 2024).
\bibitem{GeometricResolutions} Y.-Z. Liu, C. Zhang and H. Zhang, \emph{Constructing projective resolution and taking cohomology for gentle algebras in the geometric model}, \href{https://arxiv.org/abs/2308.07220v1}{arXiv:2308.07220v1} (2023).
\bibitem{GeometricKoszul} Z. Li, Y. Qiu and Y. Zhou, \emph{A geometric realization of Koszul duality for graded gentle algebras}, \href{https://arxiv.org/abs/2403.15190v1}{arXiv:2403.15190v1} (2024).
\bibitem{SurfaceCuts} W. Chang, H. Jin and S. Schroll, \emph{Recollements of partially wrapped Fukaya categories and surface cuts}, \href{https://arxiv.org/abs/2206.11196v2}{arXiv:2206.11196v2} (2023).
\bibitem{ResolvingTreesI} B. Dequ\^ene and M. Schoonheere, \emph{Resolving subcategories for gentle algebras I: Monogeneous resolving subcategories for gentle trees}, \href{https://arxiv.org/abs/2502.20994v3}{arXiv:2502.20994v3} (2026).
\bibitem{BuchsbaumEisenbud} D.A. Buchsbaum and D. Eisenbud, \emph{What makes a complex exact?}, J. Algebra 25 (1973), 259--268, \href{https://eisenbud.github.io/papers/pdfs/1973-002.pdf}{primary paper}.
\bibitem{Bruns} W. Bruns, \emph{The existence of generic free resolutions and related objects}, Math. Scand. 55 (1984), 33--46.
\bibitem{Eisenbud} D. Eisenbud, \emph{Homological algebra on a complete intersection, with an application to group representations}, Trans. AMS 260 (1980), 35--64.
\bibitem{BeauvilleLaszlo} A. Beauville and Y. Laszlo, \emph{Un lemme de descente}, C. R. Acad. Sci. Paris 320 (1995), 335--340, \href{https://math.univ-cotedazur.fr/~beauvill/pubs/descente.pdf}{author text}.
\bibitem{Magyar} P. Magyar, \emph{Affine Schubert varieties and circular complexes}, \href{https://arxiv.org/abs/math/0210151v1}{arXiv:math/0210151v1} (2002).
\bibitem{KnudsenMumford} F.F. Knudsen and D. Mumford, \emph{The projectivity of the moduli space of stable curves I: Preliminaries on `det' and `Div'}, Math. Scand. 39 (1976), 19--55, \href{https://doi.org/10.7146/math.scand.a-11642}{doi:10.7146/math.scand.a-11642}.
\bibitem{Breuning} M. Breuning, \emph{Determinant functors on triangulated categories}, \href{https://arxiv.org/abs/math/0610435v1}{arXiv:math/0610435v1} (2006).
\bibitem{DeligneRiemannRoch} D. Eriksson and G. Freixas i Montplet, \emph{Deligne--Riemann--Roch and intersection bundles}, J. \'Ec. polytech. Math. 11 (2024), 247--361, \href{https://doi.org/10.5802/jep.254}{doi:10.5802/jep.254}.
\bibitem{DeligneWeilI} P. Deligne, \emph{La conjecture de Weil I}, Publ. Math. IH\'ES 43 (1974), 273--307, \href{https://www.numdam.org/item/PMIHES_1974__43__273_0.pdf}{primary text}.
\bibitem{DeligneWeilII} P. Deligne, \emph{La conjecture de Weil II}, Publ. Math. IH\'ES 52 (1980), 137--252, \href{https://numdam.org/articles/10.1007/BF02684780/}{primary text}.
\bibitem{SGAFourHalf} P. Deligne,
\emph{Cohomologie \'etale (SGA $4\frac12$)}, Lecture Notes in Mathematics 569,
Springer-Verlag, 1977,
\href{https://doi.org/10.1007/BFb0091516}{doi:10.1007/BFb0091516}.

\bibitem{Tabuada} G. Tabuada, \emph{Noncommutative Weil conjecture}, \href{https://arxiv.org/abs/1808.00950v7}{arXiv:1808.00950v7} (6 December 2019).
\bibitem{DeningerFoliated} C. Deninger, \emph{Arithmetic geometry and analysis on foliated spaces}, \href{https://arxiv.org/abs/math/0505354v1}{arXiv:math/0505354v1} (2005).
\bibitem{RaySinger} D.B. Ray and I.M. Singer, \emph{R-Torsion and the Laplacian on Riemannian manifolds}, Advances in Math. 7 (1971), 145--210, \href{https://doi.org/10.1016/0001-8708(71)90045-4}{doi:10.1016/0001-8708(71)90045-4}.
\bibitem{Bornemann} F. Bornemann, \emph{On the numerical evaluation of Fredholm determinants}, \href{https://arxiv.org/abs/0804.2543v1}{arXiv:0804.2543v1}; Math. Comp. 79 (2010), 871--915.



\bibitem{ResolvingTreesII} B. Dequ\^ene and M. Schoonheere,
\emph{Resolving subcategories for gentle algebras II: Resolving subcategories for gentle trees},
\href{https://arxiv.org/abs/2509.10120v2}{arXiv:2509.10120v2} (3 October 2025).

\bibitem{Procesi} C. Procesi, \emph{The invariant theory of $n\times n$ matrices},
Advances in Math. 19 (1976), 306--381,
\href{https://doi.org/10.1016/0001-8708(76)90027-X}{doi:10.1016/0001-8708(76)90027-X}.

\bibitem{LeBruynProcesi} L. Le Bruyn and C. Procesi,
\emph{Semisimple representations of quivers}, Trans. Amer. Math. Soc. 317 (1990), 585--598,
\href{https://doi.org/10.1090/S0002-9947-1990-0958897-0}{doi:10.1090/S0002-9947-1990-0958897-0}.

\bibitem{CrawYamagishi} A. Craw and R. Yamagishi,
\emph{The Le Bruyn--Procesi theorem following Lusztig},
\href{https://arxiv.org/abs/2312.08527v4}{arXiv:2312.08527v4} (14 November 2025).

\bibitem{Etingof} P. Etingof, O. Golberg, S. Hensel, T. Liu, A. Schwendner, D. Vaintrob and E. Yudovina,
\emph{Introduction to Representation Theory}, with historical interludes by S. Gerovitch,
Student Mathematical Library 59, American Mathematical Society, 2011,
\href{https://doi.org/10.1090/stml/059}{doi:10.1090/stml/059};
\href{https://math.mit.edu/~etingof/representationtheorybook.pdf}{author-hosted text}.

\bibitem{SchurComplexes} K. Akin, D.A. Buchsbaum and J. Weyman,
\emph{Schur functors and Schur complexes}, Advances in Math. 44 (1982), 207--278,
\href{https://doi.org/10.1016/0001-8708(82)90039-1}{doi:10.1016/0001-8708(82)90039-1}.

\bibitem{DerivedSchur} Q. Jiang, \emph{Derived Grassmannians and derived Schur functors},
\href{https://arxiv.org/abs/2212.10488v3}{arXiv:2212.10488v3} (6 July 2023).

\bibitem{SchurRSK} A. Price, A. Stelzer and A. Yong,
\emph{Representations from matrix varieties, and filtered RSK},
\href{https://arxiv.org/abs/2403.09938v2}{arXiv:2403.09938v2} (6 October 2025).

\bibitem{DerksenWeyman} H. Derksen and J. Weyman,
\emph{Semi-invariants of quivers and saturation for Littlewood--Richardson coefficients},
J. Amer. Math. Soc. 13 (2000), 467--479,
\href{https://sites.lsa.umich.edu/hderksen/wp-content/uploads/sites/614/2018/05/A.I.a.6.pdf}{primary author text}.

\bibitem{SchofieldVdB} A. Schofield and M. Van den Bergh,
\emph{Semi-invariants of quivers for arbitrary dimension vectors},
\href{https://arxiv.org/abs/math/9907174v1}{arXiv:math/9907174v1} (27 July 1999);
Indag. Math. 12 (2001), 125--138.

\bibitem{MilneReductive} J.S. Milne, \emph{Reductive Groups},
\href{https://www.jmilne.org/math/CourseNotes/RG.pdf}{author's course notes}, version 2.00 (10 March 2018), \S5.18.

\bibitem{CrawleyBoeveyStrings} W. Crawley-Boevey,
\emph{Classification of modules for infinite-dimensional string algebras},
Trans. Amer. Math. Soc. 370 (2018), 3289--3313,
\href{https://doi.org/10.1090/tran/7032}{doi:10.1090/tran/7032};
\href{https://arxiv.org/abs/1308.6410v2}{arXiv:1308.6410v2} (4 March 2016).

\bibitem{AuslanderReitenCM} M. Auslander and I. Reiten,
\emph{Almost split sequences for Cohen--Macaulay-modules}, Math. Ann. 277 (1987), 345--349,
\href{https://doi.org/10.1007/BF01457367}{doi:10.1007/BF01457367}.

\bibitem{BautistaLiuPaquette} R. Bautista, S. Liu and C. Paquette,
\emph{Representation Theory of an Infinite Quiver},
\href{https://arxiv.org/abs/1109.3176v1}{arXiv:1109.3176v1} (14 September 2011);
published as \emph{Representation theory of strongly locally finite quivers},
Proc. London Math. Soc. 106 (2013), 97--162.

\bibitem{Paquette} C. Paquette,
\emph{Irreducible morphisms and locally finite dimensional representations},
\href{https://arxiv.org/abs/1508.04353v3}{arXiv:1508.04353v3} (4 January 2016).

\bibitem{LiuNgPaquette} S. Liu, P. Ng and C. Paquette,
\emph{Almost split sequences and approximations},
\href{https://arxiv.org/abs/1203.5009v1}{arXiv:1203.5009v1} (19 March 2012).

% The Ringel URL is unversioned. Retrieved 2026-11-03, 57 pages.
% SHA256 604cea239cfeacd2e2a04fa47ebcceb43423d7120c14492b9e5759d8b4bf15c2
\bibitem{Ringel} C.M. Ringel,
\emph{The Minimal Representation-Infinite Algebras which are Special Biserial},
\href{https://www.math.uni-bielefeld.de/~ringel/opus/tokyo.pdf}{author-hosted ICRA Tokyo 2010 proceedings text},
57 pages, retrieved 3 November 2026.

\bibitem{PuninskiPrest} G. Puninski and M. Prest,
\emph{Ringel's conjecture for domestic string algebras},
\href{https://arxiv.org/abs/1407.7470v2}{arXiv:1407.7470v2} (13 September 2015).

\bibitem{LakingPrestPuninski} R. Laking, M. Prest and G. Puninski,
\emph{Krull--Gabriel dimension of domestic string algebras},
\href{https://arxiv.org/abs/1506.07703v2}{arXiv:1506.07703v2} (29 June 2015).

\bibitem{HuangLinSu} H.-L. Huang, Z. Lin and X. Su,
\emph{The Auslander--Reiten quivers of string algebras of affine type $\widetilde C$
and a conjecture by Geiss--Leclerc--Schr\"oer},
\href{https://arxiv.org/abs/2110.02777v2}{arXiv:2110.02777v2} (16 November 2021).

\bibitem{AbramsArandaPino} G. Abrams and G. Aranda Pino,
\emph{The Leavitt path algebra of a graph},
\href{https://arxiv.org/abs/math/0509494v1}{arXiv:math/0509494v1} (21 September 2005).

\bibitem{BPRS} T. Bates, D. Pask, I. Raeburn and W. Szyma\'nski,
\emph{The $C^*$-algebras of row-finite graphs}, New York J. Math. 6 (2000), 307--324,
\href{https://nyjm.albany.edu/j/2000/6-14.pdf}{primary journal text}.

\bibitem{FowlerRaeburn} N.J. Fowler and I. Raeburn,
\emph{The Toeplitz algebra of a Hilbert bimodule},
\href{https://arxiv.org/abs/math/9806093v3}{arXiv:math/9806093v3} (23 June 1998).

\bibitem{BrownloweLacaRobertsonSims} N. Brownlowe, M. Laca, D. Robertson and A. Sims,
\emph{Reconstructing directed graphs from generalised gauge actions on their Toeplitz algebras},
\href{https://arxiv.org/abs/1812.08903v1}{arXiv:1812.08903v1} (21 December 2018).

\bibitem{KakariadisShalit} E.T.A. Kakariadis and O.M. Shalit,
\emph{On operator algebras associated with monomial ideals in noncommuting variables},
\href{https://arxiv.org/abs/1501.06495v3}{arXiv:1501.06495v3} (21 November 2018).

\bibitem{ShalitSolel} O.M. Shalit and B. Solel, \emph{Subproduct systems},
\href{https://arxiv.org/abs/0901.1422v3}{arXiv:0901.1422v3} (26 May 2009).

\bibitem{ConnesConsaniKnots} A. Connes and C. Consani,
\emph{Knots, primes and class field theory},
\href{https://arxiv.org/abs/2501.06560v1}{arXiv:2501.06560v1} (11 January 2025).

\bibitem{Prolate} A. Connes, C. Consani and H. Moscovici,
\emph{Zeta zeros and prolate wave operators},
\href{https://arxiv.org/abs/2310.18423v2}{arXiv:2310.18423v2} (4 May 2024).

\bibitem{LocalMoments} A. Connes, C. Consani and H. Moscovici,
\emph{On $q$-series and the moment problem associated to local factors},
\href{https://arxiv.org/abs/2403.01247v1}{arXiv:2403.01247v1} (2 March 2024).

\bibitem{ZetaTriples} A. Connes, C. Consani and H. Moscovici,
\emph{Zeta Spectral Triples},
\href{https://arxiv.org/abs/2511.22755v1}{arXiv:2511.22755v1} (27 November 2025).

\bibitem{QuadraticForms} A. Connes and W.D. van Suijlekom,
\emph{Quadratic Forms, Real Zeros and Echoes of the Spectral Action},
\href{https://arxiv.org/abs/2511.23257v1}{arXiv:2511.23257v1} (28 November 2025).

\bibitem{ConnesSurvey} A. Connes,
\emph{The Riemann Hypothesis: Past, Present and a Letter Through Time},
\href{https://arxiv.org/abs/2602.04022v1}{arXiv:2602.04022v1} (3 February 2026).

\bibitem{ConnesConsaniJacobian} A. Connes and C. Consani,
\emph{On the Jacobian of $\overline{\operatorname{Spec}\mathbb Z}$},
\href{https://arxiv.org/abs/2602.15941v1}{arXiv:2602.15941v1} (17 February 2026).

\bibitem{DeningerPackets} C. Deninger, \emph{Dynamical systems for arithmetic schemes},
\href{https://arxiv.org/abs/1807.06400v4}{arXiv:1807.06400v4} (7 February 2024);
Indag. Math. 37 (2026), 25--136,
\href{https://doi.org/10.1016/j.indag.2024.05.007}{doi:10.1016/j.indag.2024.05.007}.

\bibitem{DeningerKnots} C. Deninger, \emph{Primes, knots and periodic orbits},
\href{https://arxiv.org/abs/2301.11643v1}{arXiv:2301.11643v1} (27 January 2023).

\bibitem{DeningerReal} C. Deninger,
\emph{There is no ``Weil-''cohomology theory with real coefficients for arithmetic curves},
\href{https://arxiv.org/abs/2204.02714v2}{arXiv:2204.02714v2} (27 January 2023).

\bibitem{PongeTian} R. Ponge and Y. Tian,
\emph{Nonclassical Weyl laws and Connes' Integration for weak Lorentz ideals, I},
\href{https://arxiv.org/abs/2605.25591v2}{arXiv:2605.25591v2} (28 September 2026).

\bibitem{DLMF} F.W.J. Olver et al. (eds.),
\emph{NIST Digital Library of Mathematical Functions},
\href{https://dlmf.nist.gov/2.10.E9}{finite summation by parts, Eq.~2.10.9} and
\href{https://dlmf.nist.gov/25.15}{Dirichlet $L$-functions, \S25.15}, consulted 3 November 2026.

\bibitem{GraphPullbacks} A. Chirvasitu, P.M. Hajac and M. Tobolski,
\emph{Non-surjective pullbacks of graph C*-algebras from non-injective pushouts of graphs},
\href{https://arxiv.org/abs/1907.10260v3}{arXiv:1907.10260v3} (31 December 2019), Theorems3.4 and4.3;
\href{https://doi.org/10.1112/blms.12389}{doi:10.1112/blms.12389}.
\bibitem{LekiliPolishchuk} Y. Lekili and A. Polishchuk,
\emph{Derived equivalences of gentle algebras via Fukaya categories},
\href{https://arxiv.org/abs/1801.06370v5}{arXiv:1801.06370v5} (27 August 2019),
Definition3.1.1, Lemma3.1.3 and Theorem3.2.2.
\bibitem{TomfordeEmbedding} M. Tomforde,
\emph{Uniqueness theorems and ideal structure for Leavitt path algebras},
\href{https://arxiv.org/abs/math/0612628v6}{arXiv:math/0612628v6} (27 February 2007), Theorem7.3, pp.29--30.
\bibitem{WeilJacobi} A. Weil, \emph{Numbers of solutions of equations in finite fields},
Bull. Amer. Math. Soc.55 (1949),497--508, especially pp.499--502 and Eq.(4),
\href{https://doi.org/10.1090/S0002-9904-1949-09219-4}{doi:10.1090/S0002-9904-1949-09219-4};
\href{https://www.sas.rochester.edu/mth/sites/doug-ravenel/otherpapers/weil.pdf}{primary-paper scan}.
\bibitem{StacksDeterminants} The Stacks Project,
\href{https://stacks.math.columbia.edu/tag/0FJ9}{Tag0FJ9} and
\href{https://stacks.math.columbia.edu/tag/0GSX}{Tag0GSX} (retrieved 3 November 2026).
\bibitem{AlperDescent} J. Alper, \emph{Good moduli spaces for Artin stacks},
Ann. Inst. Fourier63 (2013),2349--2402, Theorem10.3, pp.2386--2387,
\href{https://doi.org/10.5802/aif.2833}{doi:10.5802/aif.2833}.
\bibitem{RootStackGuide} The Stacks Project,
\href{https://stacks.math.columbia.edu/tag/04V8}{Tag04V8: Taking roots of line bundles} (retrieved 3 November 2026).

\bibitem{ZwaraCorrected} G. Zwara, \emph{Smooth morphisms of module schemes},
Proc. London Math. Soc. (3) 84 (2002), 539--558,
\href{https://doi.org/10.1112/S0024611502013382}{primary publication}.
\bibitem{RiedtmannZwaraCorrected} C. Riedtmann and G. Zwara,
\emph{Orbit closures and rank schemes}, Comment. Math. Helv. 88 (2013), 55--84,
\href{https://doi.org/10.4171/CMH/278}{primary publication}, Section6.
\bibitem{CoveringLecture2026} Y. Liu, N. Li, B. Xing and P. Chen,
\emph{A lecture note on covering theory in representation theory of algebras},
\href{https://arxiv.org/abs/2508.01563v5}{arXiv:2508.01563v5}, 28 May 2026.
\bibitem{TameGLS2026} Q. Dong, Z. Lin, M. Lu and S. Ruan,
\emph{Representation-tame Gei\ss--Leclerc--Schr\"oer algebras and a revised
GLS conjecture on root systems},
\href{https://arxiv.org/abs/2609.26503v1}{arXiv:2609.26503v1}, 22 September 2026.
\bibitem{SchrollGeometric2026} S. Schroll,
\emph{On geometric models in representation theory}, Proceedings of ICM 2026,
vol.3, pp.250--271,
\href{https://doi.org/10.1137/25M1810647}{doi:10.1137/25M1810647};
author preprint \href{https://arxiv.org/abs/2601.14396v1}{arXiv:2601.14396v1}.

\bibitem{KellerAinf} B. Keller, \emph{Introduction to $A_\infty$ algebras and modules},
\href{https://arxiv.org/abs/math/9910179v2}{arXiv:math/9910179v2} (2001), especially \S\S3.1,3.3,6.2.
\bibitem{StacksKoszulDG} The Stacks Project Authors, \emph{The Koszul complex},
\href{https://stacks.math.columbia.edu/tag/0621}{Tag 0621}, especially Definitions15.29.1--2 and Lemma15.29.5; snapshot November3,2026, hash recorded in the repository.
\bibitem{StacksKoszulRegular} The Stacks Project Authors, \emph{Regular sequences are Koszul regular},
\href{https://stacks.math.columbia.edu/tag/062F}{Tag 062F}, Lemma15.31.2; snapshot November3,2026, hash recorded in the repository.

\bibitem{WelchGeometry} S. Datta, S. Howard and D. Cochran, \emph{Geometry of the Welch Bounds}, \href{https://arxiv.org/abs/0909.0206v2}{arXiv:0909.0206v2} (2012), Section2.1, Theorem2.2.
\bibitem{WilsonKogutRG} K.G. Wilson and J. Kogut, \emph{The renormalization group and the epsilon expansion}, Physics Reports12(1974),75--199, \href{https://doi.org/10.1016/0370-1573(74)90023-4}{DOI10.1016/0370-1573(74)90023-4}; pinned institutional scan and review scope in \path{planning/ACTUAL_RG_SOURCES.json}.
\bibitem{LevinNaveRG} M. Levin and C.P. Nave, \emph{Tensor renormalization group approach to2Dclassical lattice models}, \href{https://arxiv.org/abs/cond-mat/0611687v2}{arXiv:cond-mat/0611687v2} (2008); Phys. Rev. Lett.99(2007),120601.
\bibitem{QTokenGT} J. Kim et al., \emph{Pure Transformers are Powerful Graph Learners}, \href{https://arxiv.org/abs/2207.02505v2}{arXiv:2207.02505v2} (2022). Our finite measured-feature prototype does not reproduce its full benchmarks.
\bibitem{QGraphCG} S. Liu et al., \emph{Unsupervised Discovery of Steerable Factors When Graph Deep Generative Models Are Entangled}, \href{https://arxiv.org/abs/2401.17123v1}{arXiv:2401.17123v1} (2024); scoped primary reading in the computational source manifest.
\bibitem{QContinuousGFN} S. Lahlou et al., \emph{A theory of continuous generative flow networks}, \href{https://arxiv.org/abs/2301.12594v2}{arXiv:2301.12594v2} (2023); state measures and trajectory balance.
\bibitem{QGrossSiebert} M. Gross and B. Siebert, \emph{Mirror Symmetry via Logarithmic Degeneration Data II}, \href{https://arxiv.org/abs/0709.2290v2}{arXiv:0709.2290v2}; used as a geometric research context, not an implemented mirror comparison.
\bibitem{GreenTaoVaughan} B. Green and T. Tao,
\emph{Quadratic uniformity of the M\"obius function},
\href{https://arxiv.org/abs/math/0606087v2}{arXiv:math/0606087v2},
Section4, Lemma4.1, printed pages12--13. Focused identity reading only;
no nilsequence estimate is asserted for the moving sextic family here.
\bibitem{KudlaTateTypeI} S.S. Kudla, \emph{Tate's thesis}, expanded Jerusalem lecture notes (2001), Theorem4.3 and Corollary4.4, \href{https://u.cs.biu.ac.il/~reznikov/courses/kudla-1.pdf}{pinned primary notes}.
\bibitem{MilneCFTTypeI} J.S. Milne, \emph{Class Field Theory}, version4.03 (6August2020), Theorem0.8 and ChapterI,\S1, \href{https://www.jmilne.org/math/CourseNotes/CFT.pdf}{author notes}.

\bibitem{LFNTK} A. Jacot, F. Gabriel and C. Hongler,
\emph{Neural Tangent Kernel: Convergence and Generalization in Neural Networks},
\href{https://arxiv.org/abs/1806.07572v4}{arXiv:1806.07572v4}; specified classical scaling hypotheses.
\bibitem{LFQNTK} J. Liu et al.,
\emph{Representation Learning via Quantum Neural Tangent Kernels},
\href{https://arxiv.org/abs/2111.04225v3}{arXiv:2111.04225v3}; focused version/abstract review.
\bibitem{LFAffine} M. Kontsevich and Y. Soibelman,
\emph{Affine structures and non-archimedean analytic spaces},
\href{https://arxiv.org/abs/math/0406564v1}{arXiv:math/0406564v1}; valued and integral-affine data.
\bibitem{LFFukaya} Y. Lekili and A. Polishchuk,
\emph{Derived equivalences of gentle algebras via Fukaya categories},
\href{https://arxiv.org/abs/1801.06370v5}{arXiv:1801.06370v5}; category and smoothness hypotheses retained.
\bibitem{LFQFM} Z. Cui, P. Zhang and Y. Tang,
\emph{Quantum Flow Matching},
\href{https://arxiv.org/abs/2508.12413v1}{arXiv:2508.12413v1}; abstract-level comparison, not an upstream reproduction.
\bibitem{LFAbrams} L. Abrams,
\emph{Two-dimensional topological quantum field theories and Frobenius algebras},
J. Knot Theory Ramifications 5 (1996), 569--587,
\href{https://www.math.umd.edu/~jmr/AbramsTQFT.pdf}{author-hosted paper}.
\bibitem{LFSegal} G. Segal,
\emph{The definition of conformal field theory},
in \emph{Topology, Geometry and Quantum Field Theory} (2004);
primary publisher metadata and sewing-scope comparison.
\bibitem{LFRyu} S. Ryu and T. Takayanagi,
\emph{Holographic Derivation of Entanglement Entropy from AdS/CFT},
\href{https://arxiv.org/abs/hep-th/0603001v2}{arXiv:hep-th/0603001v2}; external setting only.
\bibitem{LFAnabelian} S. Mochizuki,
\emph{The Absolute Anabelian Geometry of Hyperbolic Curves},
\href{https://www.kurims.kyoto-u.ac.jp/~motizuki/Absolute%20Anabelian%20Geometry.pdf}{author-hosted paper},
introduction and Theorem 1.3.4; arithmetic extension hypotheses retained.
\bibitem{LFMIP} Z. Ji, A. Natarajan, T. Vidick, J. Wright and H. Yuen,
\emph{MIP*=RE},
\href{https://arxiv.org/abs/2001.04383v4}{arXiv:2001.04383v4} (16 October 2026);
external theorem and version record, not rerun locally.


\bibitem{DSFHS} T. Fukui, Y. Hatsugai and H. Suzuki,
\emph{Chern Numbers in Discretized Brillouin Zone},
arXiv:cond-mat/0503172v2, 2005.
\bibitem{DSRDKit} RDKit contributors,
\emph{MMFF94 force fields and ETKDG conformer APIs},
\url{https://www.rdkit.org/docs/}; executed version2026.03.1.
\bibitem{DSPDB} Worldwide Protein Data Bank,
\emph{Usage Policies}, CC0-1.0 coordinate archive,
\url{https://www.wwpdb.org/about/usage-policies}.
\bibitem{DSCISS} F. Evers et al.,
\emph{Theory of Chirality Induced Spin Selectivity: Progress and Challenges},
arXiv:2108.09998v1, 2021.
\bibitem{DSBridge} A. Privitera et al.,
\emph{Direct detection of spin polarization in photoinduced charge transfer
through a chiral bridge}, arXiv:2205.05353v1, 2022.
\bibitem{VisPDB} wwPDB, Atomic Coordinate Entry Format3.3, secondary structure and coordinate sections, accessed2026-11-04. \url{https://www.wwpdb.org/documentation/file-format-content/format33/sect5.html}.
\bibitem{VisScore} Y.~Song et al., Score-Based Generative Modeling through Stochastic Differential Equations, arXiv:2011.13456v1 (2020).
\bibitem{VisFlow} Y.~Lipman et al., Flow Matching for Generative Modeling, arXiv:2210.02747v1 (2022).
\bibitem{VisRFM} R.~T.~Q.~Chen and Y.~Lipman, Flow Matching on General Geometries, arXiv:2302.03660v1 (2023).
\bibitem{HubsSpokes} The Univalent Foundations Program, \emph{Homotopy Type Theory: Univalent Foundations of Mathematics}, first-edition-257-g5561b73 (2013), Sections~6.6--6.8, especially Remarks~6.7.1--6.7.2, \href{https://homotopytypetheory.org/wp-content/uploads/2013/03/hott-online-257-g5561b73.pdf}{pinned official edition}.
\bibitem{HLNTK} A. Jacot, F. Gabriel and C. Hongler, \emph{Neural Tangent Kernel: Convergence and Generalization in Neural Networks}, \href{https://arxiv.org/abs/1806.07572v4}{arXiv:1806.07572v4} (2020). This is background; our diagonal holomorphic readout limit is derived separately.
\bibitem{HLDirichlet} NIST Digital Library of Mathematical Functions, \emph{Dirichlet L-functions}, Section~25.15, equations~25.15.1--25.15.5, \href{https://dlmf.nist.gov/25.15}{primary reference}, accessed 4 November 2026.
\bibitem{HLFredholm} F. Bornemann, \emph{On the Numerical Evaluation of Fredholm Determinants}, \href{https://arxiv.org/abs/0804.2543v2}{arXiv:0804.2543v2} (2008), Section~4, equation~(4.1); Math. Comp. 79 (2010), 871--915.
\bibitem{GameTilted} C. Bamps and S. Pironio, \emph{Sum-of-squares decompositions for a family of CHSH-like inequalities and their application to self-testing}, \href{https://arxiv.org/abs/1504.06960v1}{arXiv:1504.06960v1} (2015).
\bibitem{GameNPA} M. Navascu\'es, S. Pironio and A. Ac\'in, \emph{A convergent hierarchy of semidefinite programs characterizing the set of quantum correlations}, \href{https://arxiv.org/abs/0803.4290v1}{arXiv:0803.4290v1} (2008).
\bibitem{GameMIP} Z. Ji, A. Natarajan, T. Vidick, J. Wright and H. Yuen, \emph{MIP*=RE}, \href{https://arxiv.org/abs/2001.04383v4}{arXiv:2001.04383v4} (16 October 2026). External theorem, not re-proved or newly formalized here.
\end{thebibliography}

\end{document}
